% Generated by report/build_transcripts.py. Do not edit by hand; rerun the script instead.

\paragraph{Provenance of the replies.} The prompts and replies below are verbatim from chat sessions with Claude (Anthropic). The chats had no file uploads, code execution or web access, so Claude saw only what a prompt contained (and, for a follow-up, the earlier turns) and could not compute anything on our data; every number it gave was checked with code. The exchange 02 follow-up was sent in a fresh chat, without the first reply. Two frozen tests took their design from these replies, the 1994--2009 rule of exchange 01 and the 1931--1969 run of exchange 03; both were fixed before they ran and cannot be re-drawn, because their windows have been spent.

\paragraph{Layout.} Each exchange opens with a four-line box and my evaluation, followed by the prompts and replies and the fact-check. Non-ASCII characters in the transcripts are mapped to ASCII so that no glyph is lost, and line breaks inside a listing are marked with an arrow. Word counts are \texttt{wc}-style counts of the source file; where Markdown symbols (table bars, bullets) inflate that count, the count without them follows, and that is the count compared with a word cap.

\subsection{Exchange 01: idea generation (a design review of the halfway strategy)}\label{app:ex01}

Files: \texttt{exchange/\allowbreak{}01\_\allowbreak{}idea\_\allowbreak{}generation/\allowbreak{}}. Answered by: Claude (see the provenance note). Date: 2026-09-26. Purpose: idea generation.

\begin{center}\fbox{\begin{minipage}{0.95\linewidth}\small
\begin{tabularx}{\linewidth}{@{}lX@{}}
\textbf{Asked.} & A skeptical design review of the halfway strategy: at most six problems ranked by impact on the verdict, a ranking of the seven teammate proposals (A to G), exactly five falsifiable new ideas, a plan, and the three assumptions it was least sure of. \\
\textbf{Came back.} & A 1,609-word review (cap 1,800); after I sent five corrections and a data fact, a 675-word follow-up (cap 600) with one frozen rule for 1993--2009, a four-part pass bar and an attribution specification. \\
\textbf{Verified.} & 52 checkable claims scored against the data (26 correct, 15 partly correct, 6 wrong, 5 unverifiable); its proposed tests were run on seen data (checks c01 to c11), except idea 1 (kept for the frozen test), idea 4 (left to M4), idea 5 and proposal F. \\
\textbf{Kept.} & The always-short benchmark for every timing claim, the pre-registered test on 1994-03 to 2009-12 (M8), and Aero and Ships as suspect Brown members (M4). \\
\end{tabularx}\end{minipage}}\end{center}

\paragraph{My evaluation}

\begin{quote}\small
The first reply was a competent design review built on a variable it never examined. The fact-check scored 52 checkable claims: 26 correct, 15 partly correct, 6 wrong and 5 unverifiable. Its arithmetic, citations and labeling of guesses were sound.

\textbf{Three contributions changed the plan.} First, critique 1: timed returns had been tested against zero, when a timing claim's null is the always-short position with the same hedge. On seen data it settles the question (timed minus always-short FF5+UMD alpha peaks at t = 1.09). Second, Idea 1, a backward test on the one window with no computed strategy returns (pre-2010), was not in my plan. Third, it flagged Aero and Ships as odd Brown members and guessed that Aero, Util and Ships drove the COVID gains. Both held: EPA factors rank Aero 31st and Ships 20th of 49, and those three industries deliver 86\% of the Original 3-month rule's COVID gross P\&L. Its redirect of test C to the Brown leg was also right (HML +0.44 for Brown vs. +0.14 for Green since 2010).

\textbf{Six claims were wrong} (five corrections in my follow-up). The holdout was not lost to duration: the hedged Brown leg has no holdout rate exposure ($\Delta$10y t = 0.05), and holdout alphas stay at $-$2.3\% to $-$4.8\% annually after FF5+UMD, rate and commodity controls. Its inference fix fails: more Newey--West lags raise the 24-month COVID t from 3.21 to 5.32, and the binding problem is small-sample p-values (with t(n$-$4), 0 of 32 cells survive Holm). The IC flip is 1.3 to 1.6 standard errors for the raw and purified signals and 2.8 to 2.9 Newey--West standard errors for the continuous ones, not ``about one''. Idea 1's dates were wrong (two claims; first crossings 1993-01 and 1999-01), and the team's macro-state notebook had already used 192 pre-2010 months. The sixth, on coverage, is below. Two weaker points were not scored wrong: it asked where Coal ranks (9th of 41, in the file) but missed that EPA v1.3 lacks electricity codes, and ideas 4 and 5 were generic climate-premium portfolios with no stated alpha source.

\textbf{It never asked what ``attention'' measures.} My prompt passed on the team's label without a source, and neither reply asked who builds the index. The series is the EMV tracker for Energy and Environmental Regulation (since 2010, 88.5\% of crossings fall in months when overall EMV is above its past-only median, against 56.1\% of other months; Fisher p = 0.0012), and 31 of 48 holdout months are exactly zero. It put the coverage risk in the wrong period (``early index coverage may be thin''; the pre-2010 zero share is 0.8\% to 4.2\%). This is a failure of my prompt as much as of the replies.

\textbf{The follow-up turned the corrections into specifications.} It accepted every correction, kept ``Do not implement'', and gave a frozen 1993--2009 rule (share of overall EMV, zeros as missing, the 6-month rule, a four-part pass bar) plus an attribution against an exposure-matched always-short benchmark, with a calendar-shuffle null. It rightly declined credit for the holdout: its zero rule silences the signal in all 48 holdout months but was written after seeing them. Three flaws remain. The two zero thresholds are redundant (at most 6 zeros in 60 makes the 48-nonzero minimum inert). The share signal keeps most of the original timing (15 of its 25 post-2010 crossings sit within a month of the team's, against about 10 by chance). Its k = 14 undercounts the model's 15 coefficients.

\textbf{Where each piece went.}

\begin{itemize}
\item M1b: MCCC and CPU stay the primary alternatives on 2010--2026, with VIX and overall EMV as COVID placebos. The share series was already a robustness measure; the follow-up's point that MCCC starts in 2003 and CPU cannot signal before 1995 is why it carries the pre-2010 test.
\item M2 and M4: test C covers both legs, costs are reported as break-evens, and the EPA rebuild patches Util from USEEIO.
\item M3: its Ferson--Schadt terms match my planned conditional-beta step, and a GS10 bond return with duration and convexity supersedes its $-$8$\cdot$$\Delta$y10 guess.
\item M8: the exposure-matched benchmark, calendar-shuffle null, merged event counts and t(n$-$k) p-values apply to every timing claim. The 1993--2009 rule is frozen with corrected dates and a one-month publication lag, and is scored once.
\item M5 and M6 owe nothing here: neither reply proposed industry momentum or shrinkage timing.
\end{itemize}
\end{quote}

{\raggedright\paragraph{Prompt (verbatim)} \texttt{prompt.md}, 993 words (980 excluding Markdown symbols).\par}

\begin{lstlisting}[style=transcript]
You are a skeptical buy-side quant refereeing a student research project in climate-driven factor timing. Treat this as a design review: I want the weak points found before I spend time on them.

SETTING. This is the MFE 230GA (Active Asset Management) final project at Berkeley. The report needs a one-paragraph Executive Summary that says "Do not implement" if there is no alpha; "What did you try" (idea, data, risk model, turnover and costs, at least three substantial Claude interactions); "What did you learn" (factor exposures, performance over the full sample, post-2010 and the last 12 to 18 months, robustness and limits, a critical evaluation of Claude). Grading: thesis 40%, execution 40%, originality 20% (novel data use, prompt design). A well-documented "Do not implement" is an acceptable result. Our team of five wrote the first half; I am taking the second half forward.

HALFWAY DESIGN. Thesis: Green minus Brown has no permanent premium, but its factor-neutral part pays after spikes in climate-transition attention. Green is the five lowest emissions-intensity Fama-French 49 industries (Fun, RlEst, Drugs, Telcm, Fin), Brown the five highest (Util, Ships, Aero, Steel, BldMt), each leg equal-weighted across value-weighted industries. Emissions intensity is one static snapshot covering 41 of 49 industries (Oil, Chips, FabPr, Gold, Hshld, LabEq, Other, Toys missing). The traded position is not the spread: it shorts the Brown leg alone, with an FF3 overlay from rolling 60-month betas lagged one month. Signal: climate-transition attention (a monthly index starting in 1985) as a 60-month rolling z-score of log(1 + attention). A month is extreme when the signal exceeds its expanding, past-only 80th percentile, and each crossing opens a Short-Brown position held 3 or 6 months, sized to 5% residual volatility (cap 1x). Costs per unit of turnover are 10 bp on the leg, 5 bp on the market overlay and 25 bp on SMB/HML, all assumed. Variants: "purified" attention (the residual from a 120-month walk-forward ridge on rate, oil, inflation and activity shocks) and a continuous weight (floor 0.5, 3-month half-life). Validation ran 2010 to Jul 2022; Aug 2022 to Jul 2026 was a frozen holdout.

RESULTS.
- Validation: the 6-month rules earn about 2.5% a year net (FF3 alpha Newey-West t = 2.23 and 2.46); the 3-month rules have t of about 1.5.
- COVID carries it: the 3-month rule earns 5.94% net in 2020 to 2021 (alpha t = 3.21).
- Holdout: every version loses 2.1% to 3.8% a year (t = -1.7 to -2.2). The IC flips from about -0.08 (negative is the expected sign for a short-Brown rule) to +0.12 to +0.15.
- Over 2010 to 2026 the best timed version earns 1.3% a year net against 1.0% for always-short Brown.
- The raw spread loads on HML (-0.24 full sample, -0.35 since 2010). After hedging, residual HML is about -0.05 with t of about -2 in every version.
- The continuous variant cuts volatility 55 to 65% and drawdown 60 to 75%, but a calendar-shuffle placebo gives p of about 0.12.
- 2 of 32 strategy-period cells survive Holm, both in COVID.
- A follow-up notebook finds the attention slope turns positive in high-rate months (difference t = 1.83, Holm p of about 0.6); 47 of those 74 months are in the holdout.

TEAMMATE PROPOSALS. From Cristhian: (A) rebuild with 8 Green and 8 Brown industries; (B) add momentum and an explicit commodity exposure to the controls; (C) regress Fin, Telcm, Drugs, Fun and RlEst individually to find where the HML exposure comes from; (D) uniform costs at 5, 10 and 25 bp. From Charishma: (E) a Part B on whether this is alpha or beta at the portfolio level. From the group: (F) macro controls, since exposures look regime dependent (COVID, then inflation and rates); (G) a momentum control, in case the signal is only persistence in past returns.

DATA. In hand: the four team files (FF49 monthly returns, Jul 1926 to Jul 2026; FF3 and rf; the emissions snapshot; a macro file with attention, the 10-year yield, WTI, CPI, CFNAI, the NBER recession flag and GSCPI). Downloadable: Ken French FF5, UMD, short- and long-term reversal, FF49 value- and equal-weighted returns, industry BE/ME, firm counts and average size; FRED Treasury yields (3-month, 2-year, 10-year), BAA and AAA, VIX, commodity indices, activity series, breakevens and TIPS yields; the Gavriilidis Climate Policy Uncertainty index, the Ardia et al. Media Climate Change Concerns index, and EPA supply-chain GHG factors by NAICS. No firm-level data. Every result so far starts in 2010, and the holdout has now been seen, so it is no longer a clean test.

WHAT I NEED, in this order and under these headings.
1. Design critique. At most six problems, ranked by how much each could change the verdict: the problem, why it matters here, the fix.
2. Proposal ranking. One table for A to G: rank, proposal, the result that would change our conclusion, what it cannot tell us, overlap with other proposals, effort in hours. One line of reasoning per row.
3. New ideas. Exactly five, at least two extending the current strategy and at least two alternatives on the same data, each aimed at a small defensible alpha or a more rigorous "Do not implement". For each: (a) a falsifiable hypothesis and the result that rejects it; (b) the exact test: sample, return series, specification, benchmark, pass bar; (c) the data, and whether it is listed above; (d) main pitfalls (look-ahead, multiple testing, small N, overlapping holds); (e) effort in hours; (f) what it adds to the report if it fails.
4. Plan. The three items from sections 2 and 3 you would run first, in order, and why.
5. The three assumptions in your answer you are least confident about.

RULES. Work from our numbers, not generic ESG results. Cite a paper only with authors, year and the specific claim. Do not invent backtest results; label any number you have not computed as a guess. Stay under 1,800 words.
\end{lstlisting}

{\raggedright\paragraph{Reply (verbatim; Claude)} \texttt{claude\_\allowbreak{}response.md}, 1,772 words (1,609 excluding Markdown symbols).\par}

\begin{lstlisting}[style=transcriptsmall]
# Design Review: Climate-Attention Timing of Short-Brown

**Bottom line:** The holdout failed, validation rests on one episode, and timing beats always-on Short-Brown by roughly 0.3% a year (1.3% vs 1.0%, arithmetic on your figures). Make "Do not implement" airtight, and test pre-specified ideas on the one period nobody has looked at, **1990-2009**.

---

## 1. Design critique

**1. The benchmark is wrong for a timing claim.**
- *Problem:* Timed returns are tested against zero. For a timing thesis the null is "always short Brown, same hedge, same vol target."
- *Why it matters:* The signal adds about 0.3%/yr over 2010-2026. Also, the traded position has no Green leg, so "Green minus Brown" isn't what you're testing.
- *Fix:* Make the primary statistic the alpha of (timed - always-on), or regress hedged Brown returns on an in-position dummy. Restate the thesis as "Brown underperforms after attention spikes."

**2. Validation is one episode.**
- *Problem:* Both Holm survivors are COVID; the 3-month rule's 2020-21 alpha (t = 3.21) carries the result.
- *Why it matters:* Aero (Boeing-heavy: 737 MAX grounding plus the travel collapse), Ships and Util took COVID-specific hits unrelated to transition risk. With five equal-weighted industries, each is 20% of the leg.
- *Fix:* Leave-one-industry-out, and rerun excluding Mar 2020-Dec 2021. If the 6-month t-stats drop below ~1.5, the Executive Summary should say so.

**3. The holdout is spent; everything after it is post-hoc.**
- *Problem:* The high-rate interaction was found with 47 of its 74 months in the holdout.
- *Why it matters:* It can't be confirmed on the data that suggested it. The holdout also doesn't prove a *reversal*: with ~48 months the IC's standard error is roughly 1/sqrt48 ~ 0.14 (a rough approximation that ignores overlap), so -0.08 -> +0.13 is about one standard error of the difference. The honest reading is "failed to confirm," not "flipped."
- *Fix:* Declare 1990-2009 the new frozen holdout and write down one rule before running it.

**4. The Brown leg may not measure what the thesis claims.**
- *Problem:* One static snapshot, eight industries missing, including Oil, the most obvious transition-risk industry. Aircraft and Shipbuilding in the top five by emissions intensity is surprising; check the mapping and where Coal ranks.
- *Why it matters:* If the ranking is off, a "climate" result is an "industrials and utilities" result.
- *Fix:* Rebuild intensity from EPA supply-chain GHG factors (NAICS -> SIC -> FF49) covering all 49, and rerun with Oil and Coal in.

**5. Hedge residuals line up with the regime story.**
- *Problem:* Residual HML ~ -0.05 (t ~ -2) in *every* version, and no control for momentum, profitability/investment, rates or commodities.
- *Why it matters:* Util, BldMt and Steel are rate- and commodity-sensitive. A 2022-2026 holdout dominated by rates can lose money through unhedged duration, not climate.
- *Fix:* Attribute ex post against FF5 + UMD + a rate factor (e.g., Delta10-year yield) + a commodity return. Keep the ex-ante hedge unchanged so you aren't re-optimizing.

**6. Inference is weaker than the t-stats suggest.**
- *Problem:* Overlapping 3-6-month holds, few independent crossings, 4+ variants, and a purified signal built from CPI, CFNAI and GSCPI, which are published with a lag and revised.
- *Why it matters:* Too few Newey-West lags overstate t; final-vintage macro data is look-ahead. The continuous variant's lower vol and drawdown is partly mechanical (lower average exposure), which is consistent with placebo p ~ 0.12.
- *Fix:* NW lags >= hold length; report the count of independent events; block-bootstrap by event; lag macro inputs 1-2 months; compare the continuous variant with always-on scaled to the same average exposure.

---

## 2. Proposal ranking

| Rank | Proposal | Result that would change our conclusion | What it cannot tell us | Overlap | Hours | Reasoning |
|---|---|---|---|---|---|---|
| 1 | **E** Alpha vs beta | Timed - always-on alpha significant after FF5+UMD | Whether the mechanism is climate | B, F, G | 4-6 | The question graders ask, and the one your numbers are closest to answering |
| 2 | **A** 8+8 industries (add leave-one-out) | Timing alpha survives broader legs and dropping Aero | Whether emissions matter vs any 5 industries | C | 3-4 | Cheap; directly tests the COVID concentration |
| 3 | **G** Momentum control | Signal loses a horse race to Brown's own trailing 3-12m return | Whether attention causes the persistence | B | 2-3 | Cheap; kills the simplest alternative story |
| 4 | **B** Momentum + commodity controls | Alpha t falls below ~1.5 with a commodity factor | Climate vs commodity channel (Oil is missing) | G, F | 3-5 | Useful but mostly subsumed by E |
| 5 | **F** Macro controls | Alpha vanishes when conditioning on rate/inflation state | Anything confirmatory; regimes were picked after seeing the holdout | B, purified variant | 6-10 | Most expensive, most overfit-prone |
| 6 | **C** Green industries one by one | Little: Green isn't in the traded position | Source of the *traded* HML residual (that sits in Brown) | A | 2 | Redirect it to the five Brown industries |
| 7 | **D** Uniform 5/10/25 bp | Only a break-even cost below ~10 bp, unlikely with 3-6-month holds (guess) | Anything about the signal | none | 1 | Costs can't rescue a -2% to -4% holdout; report break-even instead |

---

## 3. New ideas

### Idea 1 (extension): Backward out-of-sample test, 1990-2009
- **(a) Hypothesis:** The frozen 6-month rule earns positive alpha over always-on before 2010. *Rejected* if timed - always-on alpha <= 0 or t < 1.
- **(b) Test:** Signal from Jan 1990 (60-month z-score uses 1985-89); purified variant from ~1995 (120-month ridge). Hedged Short-Brown, 5% vol target, validation parameters unchanged. FF3 and FF5+UMD alpha, NW 6 lags. Benchmark: always-short Brown. Pass: t > 2. Pre-register the high-rate interaction as a second hypothesis, Holm across both.
- **(c) Data:** In hand.
- **(d) Pitfalls:** Static snapshot applied backward (classification look-ahead); transition risk was arguably unpriced before ~2005, so failure doesn't fully kill the mechanism; early index coverage may be thin.
- **(e) Effort:** 3-5 h.
- **(f) If it fails:** Twenty more years behind "Do not implement," from the cleanest test left.

### Idea 2 (extension): Timing increment and event-time profile
- **(a) Hypothesis:** Hedged Brown residuals are more negative in t+1...t+6 after crossings than unconditionally. *Rejected* if the post-event cumulative residual is within one SE of zero, or the move sits in spike month t itself.
- **(b) Test:** 2010-2026 only until Idea 1 is run (don't contaminate the new holdout). Event-time averages t-3...t+6; regression of residual on a post-event dummy, NW 6 lags. Benchmark: always-on. Pass: dummy t > 2 excluding 2020-21.
- **(c) Data:** In hand.
- **(d) Pitfalls:** Overlapping events (merge crossings within 6 months), small N (report event count), COVID dominance.
- **(e) Effort:** 3-4 h.
- **(f) If it fails:** Evidence attention is priced on impact. Ardia, Bluteau, Boudt & Inghelbrecht (2023) find green stocks rise and brown stocks fall on days of unexpected increases in climate-change concerns; Pastor, Stambaugh & Taylor (2022) attribute green outperformance over 2012-2020 to unanticipated strengthening of climate concerns. If the repricing is contemporaneous, a post-spike rule has nothing left to harvest.

### Idea 3 (extension): Replicate the signal across attention indices
- **(a) Hypothesis:** If the effect is about climate attention, the frozen rule works with the Climate Policy Uncertainty (CPU) and Media Climate Change Concerns (MCCC) indices. *Rejected* if neither gives timed - always-on t > 1 with the same sign.
- **(b) Test:** Identical rule and parameters, swap the signal; samples set by index start dates (CPU late 1980s, MCCC ~2003; verify). Pass: same sign in all three, pooled t > 2.
- **(c) Data:** Downloadable, listed.
- **(d) Pitfalls:** Different samples per index; three tests (Holm); the 60-month z-score burns five years of a short index.
- **(e) Effort:** 4-6 h.
- **(f) If it fails:** Shows the result is index-specific noise; earns originality credit either way.

### Idea 4 (alternative): Rebuilt carbon-intensity spread across all 49 industries
- **(a) Hypothesis:** No unconditional premium: an intensity-sorted industry spread has FF5+UMD alpha of zero. Bolton & Kacperczyk (2021) find a positive carbon premium at the firm level; this tests the industry analogue. *Rejected* if alpha t > 2 in either direction.
- **(b) Test:** 1990-2026, monthly; long top-third/short bottom-third intensity, VW industries. FF5+UMD alpha, NW 12 lags; report pre/post-2010 and the last 18 months.
- **(c) Data:** EPA factors (listed); a NAICS->SIC crosswalk (not listed; public concordances exist).
- **(d) Pitfalls:** Single-vintage intensity applied backward; crosswalk error; supply-chain factors include upstream emissions, a different concept from your snapshot.
- **(e) Effort:** 6-10 h.
- **(f) If it fails:** A rigorous version of the "no permanent premium" half of your thesis, with Oil and Coal included.

### Idea 5 (alternative): Climate-news hedge portfolio on FF49
- **(a) Hypothesis:** Following Engle, Giglio, Kelly, Lee & Stroebel (2020), who build mimicking portfolios that hedge innovations in a climate news index, an FF49 portfolio mimicking attention innovations hedges out of sample. *Rejected* if its out-of-sample correlation with next-month innovations is <= 0.
- **(b) Test:** 1995-2026. Monthly rolling 60-month regressions of each industry's excess return on attention innovations (past-only AR(1) residual) plus FF3; long top-10/short bottom-10 by beta, one-month hold. Pass: OOS correlation t > 2; report alpha as the price of the hedge.
- **(c) Data:** In hand.
- **(d) Pitfalls:** Noisy betas, only 49 assets, innovation definition choice, full-sample AR fits (look-ahead).
- **(e) Effort:** 6-8 h.
- **(f) If it fails:** Reframes attention as a risk to hedge, not a signal to trade.

---

## 4. Plan
1. **Idea 2 + leave-one-out (Critique #2).** Cheapest, and decides whether timing adds anything over always-on, on data already seen.
2. **Idea 1.** The only clean out-of-sample test left. Freeze the rule in writing first, and don't touch pre-2010 attention data before then.
3. **E with B and G folded in.** FF5 + UMD + rates + commodities attribution for full sample, post-2010 and last 18 months; the report needs it whatever the outcome.

---

## 5. Assumptions I'm least confident about
1. That the attention index has usable coverage in 1985-2009 and the emissions ranking is stable enough to apply backward; if not, Idea 1 is weak.
2. That the COVID gains come from Aero, Util and Ships shocks; this is a hypothesis, not something I computed.
3. The index start dates (CPU, MCCC), the NAICS->FF49 mapping effort, and every hour estimate.
\end{lstlisting}

{\raggedright\paragraph{Follow-up prompt (verbatim)} \texttt{followup\_\allowbreak{}prompt.md}, 396 words (392 excluding Markdown symbols).\par}

\begin{lstlisting}[style=transcript]
I ran most of your proposed tests on the seen data. Five of your claims are wrong, and you missed a data fact that changes the review.

The attention series is not what we called it. I traced the team's column to its source: it is FRED EMVENRGYENVREG, the Baker, Bloom, Davis and Kost Equity Market Volatility tracker for Energy and Environmental Regulation. It counts newspaper articles on stock-market volatility that mention the category, scaled to the VIX. Since 2010 its z-score correlates +0.39 with overall EMV, +0.11 with CPU and -0.06 with MCCC, and 38% of crossing months are high-EMV months, against 12% of other months (odds ratio 4.5, p = 0.002). Your coverage worry also points the wrong way. The series is exactly zero in 12 of 441 months from 1985-01 to 2021-09, but in 39 of 59 months from 2021-10 (31 of 48 holdout months). The holdout tested a near-binary signal, not the one we validated.

Corrections:
1. Duration. The hedged Brown leg has no holdout rate exposure (Delta10y t = 0.05), and the strategies gain when yields rise. After FF5+UMD+Delta10y+commodity controls, holdout alphas are still -2.3% to -4.8% (t about -2).
2. Newey-West lags. More lags do not shrink the validation t (2.23 at 6 lags, 2.30 at 12); they push the 24-month COVID t to 5.3. The binding problem is small-sample p-values: with t(n-4), 0 of 32 cells survive Holm.
3. IC flip. It is 1.3 to 1.6 SE for raw and purified, but 2.8 to 2.9 NW SE for the continuous signals.
4. Idea 1 dates. The frozen rule first crosses in 1993-01 (raw) and 1999-01 (purified).

Your tests rescue nothing: timed minus always-on FF5+UMD alpha peaks at t = 1.09, and the Idea 2 post-event CAR is -1.29% (t = -1.08).

Two questions.
(a) Do I salvage this series (as a share of total EMV, or zero-robust) or replace it with CPU or MCCC, given both end in 2025 and burn about eight years under our rule? Give the one rule you would freeze for the unseen 1993 to 2009 test, its pass bar, and how it treats zeros.
(b) Give an exact attribution spec that separates timing alpha from duration, value and volatility exposure: dependent variable, regressors, an exposure-matched benchmark, lags, and what counts as alpha.

Under 600 words. Label anything you have not computed as a guess.
\end{lstlisting}

{\raggedright\paragraph{Follow-up reply (verbatim; Claude)} \texttt{claude\_\allowbreak{}followup\_\allowbreak{}response.md}, 733 words (675 excluding Markdown symbols).\par}

\begin{lstlisting}[style=transcriptsmall]
# Follow-up: Signal Choice and Attribution

Your corrections are right, including the ones on duration and the IC. The EMV finding matters more than any of them. The signal we validated partly tracks news about market volatility, and the holdout ran on a different, near-binary signal. That explains the holdout but does not rescue the strategy. On the data you have already seen, timed minus always-on peaks at t = 1.09, so **"Do not implement" stands**.

---

## (a) Salvage the series; don't replace it

- **MCCC can't test 1993-2009** because it starts in 2003.
- **CPU would only cover about 1995-2009** once the burn-in is used up (a guess). Its correlation with your series is +0.11, so switching to it tests a new hypothesis rather than a cleaner version of yours.
- Run CPU and MCCC only after the primary test, and label them exploratory.
- **Use the series as a share of total EMV.** As I understand the construction, the category tracker equals overall EMV times the fraction of EMV articles that mention the category. Dividing by overall EMV removes the part scaled to the VIX and leaves the attention part. This goes straight at your odds ratio of 4.5.

**The rule to freeze.** Write it down before you load any pre-2010 data.
1. **Signal:** s_t = EMVENRGYENVREG_t / EMVOVERALLEMV_t.
2. **Zeros:** treat a zero month as missing, not as low attention.
   - A zero month cannot trigger a position.
   - Drop zero months from both the z-score window and the percentile history.
   - The z-score needs at least 48 nonzero months out of the trailing 60.
   - If more than 10% of the trailing 60 months are zero, the signal is off.
3. **Extreme month:** z_t is the 60-month rolling z-score of log(s_t). A month is extreme when z_t exceeds its expanding, past-only 80th percentile.
4. **Trade:** the validated 6-month Short-Brown rule, unchanged: lagged 60-month FF3 hedge, 5% vol target, cap 1x, same costs. A new crossing during a hold extends the hold; positions never stack.
5. **Sample:** Jan 1993 to Dec 2009, run once.
6. **Pass bar:** all four of the following. Anything less confirms "Do not implement."
   - The timing alpha from (b) is above zero with t >= 2, using t(n-k) p-values.
   - The calendar-shuffle p is 0.05 or lower.
   - There are at least 8 independent episodes (a guessed threshold).
   - The alpha keeps its sign when each Brown industry is dropped in turn.

**Caveat:** the zero rule (rule 2) would probably have kept you flat for most of the holdout (a guess). It was written after seeing the holdout, though, so take no credit for that period.

---

## (b) Attribution spec

**Dependent variable.** D_t = R^T_t - pi*R^AO_t, net of costs.
- R^T is the timed strategy.
- R^AO is always-on Short-Brown with the same hedge and vol target.
- pi = mean timed weight / mean always-on weight. The benchmark then has the same average exposure as the timed strategy but no timing.

**Regressors.** These are for ex-post attribution only; keep the ex-ante hedge unchanged.
- FF5 + UMD.
- **Duration:** TERM_t ~ -8*Deltay10_t, an approximate 10-year Treasury return. The duration of 8 is a guess.
- **Commodity:** the WTI log return.
- **Volatility:** DeltaVIX_t (VIX data start in 1990) and Deltalog EMV_t.
- **Conditional terms:** I_{t-1} x {Mkt, HML, TERM, DeltaVIX}, where I_{t-1} = 1 if the strategy is in position. This follows Ferson and Schadt (1996), who let betas vary with lagged public information so that changing exposure is not counted as skill. Here it means any return from loading on value, duration or volatility only while in position counts as beta.

**Model.** D_t = alpha + beta'F_t + gamma'(I_{t-1}*F_t) + eps_t.

**Inference.**
- Newey-West standard errors with 6 lags, and 12 as a check.
- p-values from t(n-k), with k = 14 parameters.
- The null with matched exposure is a calendar shuffle: move the same hold windows to random dates 5,000 times and recompute alpha each time.

**What counts as alpha.** Only alpha.
- Report the decomposition mean(D) = alpha + Sigmabeta*mean(F) + Sigmagamma*mean(I*F). That gives duration, value and volatility a number each.
- Timing alpha exists only if alpha passes the bar in (a) and survives the volatility block.
- If alpha drops when DeltaVIX and DeltaEMV enter, the signal was a volatility bet, not a climate one.
\end{lstlisting}

\paragraph{Fact-check (verbatim excerpt)} Sections ``Summary'', ``Proposed tests that I ran'', ``What Claude missed'' of \texttt{factcheck.md} (4,299 words in full).

\begin{table}[H]\centering\small
\caption{Exchange 01 fact-check, verdict counts by section of the reply (C correct, P partly correct, W wrong, U unverifiable). Counted from the fact-check's claims table.}\label{tab:a1counts}
\begin{tabular}{@{}lrrrrr@{}}\toprule
Section of the reply & C & P & W & U & Total \\\midrule
Bottom line and design critique & 14 & 11 & 3 & 1 & 29 \\
Proposal ranking (section 2) & 2 & 0 & 0 & 0 & 2 \\
New ideas (section 3) & 7 & 3 & 3 & 2 & 15 \\
Plan and stated assumptions (sections 4, 5) and compliance & 3 & 1 & 0 & 2 & 6 \\
\midrule
All claims & 26 & 15 & 6 & 5 & 52 \\
\bottomrule\end{tabular}\end{table}

\begin{lstlisting}[style=transcriptsmall]
## Summary

52 checkable claims: 26 correct, 15 partly correct, 6 wrong, 5 unverifiable. The review is well organised, the arithmetic on our figures is right, and its literature citations are accurate. Its main weaknesses are these:

1. **It never asked what "attention" is.** The series is the EMV tracker for Energy and Environmental Regulation, a count of news articles about stock-market volatility. From 2010 to 2026 its z-score correlates +0.39 with overall EMV, +0.11 with CPU and -0.06 with MCCC. Crossings are 4.5 times as likely (odds ratio) in months of high overall EMV (p = 0.002).
2. **It got the coverage problem backwards.** It worried that early coverage "may be thin". Before 2010 the zero share is 0.8% to 4.2%. In the holdout, 31 of 48 months are exactly zero, so the holdout tests a different, near-binary signal.
3. **Idea 1 has the wrong start dates.** Under the frozen rule the first raw crossing is Jan 1993 and the first purified crossing Jan 1999, not 1990 and about 1995. Also, the team's macro-state notebook already estimated attention-to-residual slopes over 1994 to 2026.
4. **The duration story for the holdout is wrong.** The hedged Brown leg has no rate exposure in the holdout (t = 0.05). The strategies gain when yields rise. With FF5 + UMD + Delta10y + commodity controls, holdout alphas are still -2.3% to -4.8% (t about -2).
5. **Two statistical diagnoses are off.**
   - The IC flip is 1.3 to 1.6 SE for raw and purified attention, not "about one". For the continuous signals it is 2.8 to 2.9 NW SE (p < 0.01).
   - More Newey-West lags do not shrink the validation t (2.23 becomes 2.30 at 12 lags). They inflate the 24-month COVID t to 5.3. The real problem there is small-sample p-values.

Every test it proposed that I could run on seen data fails to rescue the strategy. The timed-minus-always-on FF5+UMD alpha peaks at t = 1.09. The Idea 2 post-event CAR has t = -1.08 (-0.63 ex-COVID). Momentum and commodity controls do not absorb the signal.

## Proposed tests that I ran on seen data (2010-2026), and what they show

| Test (Claude's spec) | Result | Script |
|---|---|---|
| E / critique 1: alpha of (timed minus always-on) after FF5+UMD | Never significant. 2010-2026 t -0.93 to 0.78; Val at most 1.09; HO -1.11 to 1.10. The in-position dummy (unit position) gives Val t of 2.68 to 2.75 for 6m holds, but at most 1.14 for 2010-2026 excluding Mar 2020 to Dec 2021, and -1.5 to -2.1 in the HO. | c01 |
| Critique 2: drop Mar 2020 to Dec 2021; leave one industry out | 6m Val t 1.42 and 1.65 (Claude's bar is 1.5); 3m 0.41 and 0.50. Leave-one-out 6m Val t 1.62 to 2.77; COVID 6m t collapses without Aero (0.03, 0.07). HO stays negative in every case (-1.17 to -2.56). | c02 |
| A: 8+8 legs (Hardw and MedEq tie for the 8th Green slot) | 6m Val t 2.16 and 2.31; 3m 1.02 and 1.08. Without Aero: 6m 1.61 and 1.85, 3m 0.67 and 0.69. HO -0.72 to -2.13. The 6m validation edge survives broader legs; nothing else changes. | c11 |
| B / critique 5: FF5 + UMD + Delta10y + commodity attribution | Val 6m alpha 2.2% to 2.5% (t 2.46 and 2.72, slightly higher than FF3); HO -2.3% to -4.8% (t -2.0 to -2.3). Commodity beta t -0.4 to -0.7. It does not "fall below ~1.5". | c04 |
| G: horse race with Brown's trailing 3m and 12m returns | Attention is not absorbed. Val z t is -1.26 alone and -1.84 with the controls; trailing 12m Brown return t is +2.0 (continuation). 2010-2026 z t -0.39 alone, -0.90 with controls. The p80 state correlates -0.25 with trailing 12m Brown return. | c08 |
| Idea 2: event-time profile, merged events, post-event dummy, pass if t > 2 ex 2020-21 | Rejected by its own rule. CAR over t+1 to t+6: -1.29% (t -1.08, raw) and -1.01% (t -0.94, pure); ex-2020-21 -0.82% (t -0.63) and -0.67% (t -0.57), within one SE of zero. Post-event dummy t -0.40 and +0.54; ex-COVID -0.17 and +0.29. | c07 |
| Critique 6: continuous rule vs always-on at the same exposure | Sharpe over 2010-2026: 0.24 and 0.30 vs 0.20. HO: -0.60 and -0.47 vs -0.21. | c01 |
| Idea 3 premise: do the attention indices agree? | Team z vs CPU z +0.11 and vs MCCC z -0.06 (2010-2026). CPU and MCCC levels correlate 0.70 with each other, but only 0.15 and -0.04 with the team series. Running the frozen rule on CPU and MCCC belongs to module M1b and was not run here. | c06, c10 |

Not run: idea 1 (kept clean on purpose), idea 4 (covered by M4), idea 5, proposal F.

## What Claude missed

1. **What the attention variable measures.** The prompt gave the variable as "climate-transition attention" with no source, and Claude accepted that label. The series is FRED EMVENRGYENVREG, a count of newspaper articles about stock-market volatility that mention energy or environmental regulation, scaled to the VIX.
   - Its z-score correlates +0.41 with overall EMV and +0.29 with VIX (1993-2026), but +0.13 with CPU and 0.00 with MCCC.
   - Since 2010, crossings fall in high overall-EMV months 38% of the time, against 12% otherwise (odds ratio 4.5, p = 0.002).
   - The rule partly times general market-volatility news. A careful referee would have ranked idea 3 (swap in CPU or MCCC) first. (c10; M1_signal_audit_decomposition)
2. **The holdout signal is degenerate.** There are 39 zeros in 59 months from 2021-10, against 12 in 441 before. In the holdout, zero months map to z between -1.39 and -0.59, and entries require any nonzero print. The IC "flip" and the holdout losses test a near-binary signal, not the validated one. Claude placed the coverage risk in the wrong period. (c06)
3. **Same-month look-ahead in the signal itself.** Month-t EMV is used to trade at the close of month t. Claude flagged vintage issues in the macro controls but not in the traded series, which moves the headline numbers (row 27).
4. **Small-sample inference in COVID.** Both Holm survivors are NW(6) t-stats on 24 months. With t(n-4) p-values none survive. Asking for more lags makes that t larger, not smaller.
5. **Unequal exposure in the benchmark comparison.** The "0.3% a year" compares a rule with mean exposure 0.43 to one with 0.66, and picks the best of six. The comparison should be exposure- or volatility-matched and should correct for selection.
6. **Leg weighting.** Equal weights give Ships (9 to 10 firms, about $30bn) the same 20% as Util (about $1.1tn). Ships alone accounts for 1.60 pp of the Original 3m rule's 3.36% annual holdout gross loss. (c02)
7. **Power of a 48-month holdout.** Eskildsen, Ibert, Jensen and Pedersen (2026) give the relevant arithmetic: expected t = 2 x Sharpe over 4 years, so even a true Sharpe of 0.5 gives an expected t of 1.0. This would have sharpened critique 3 more than the IC standard error does. (lit_digest section 7)
8. **Emissions-measure gaps.** EPA v1.3 has no electricity codes, so Util needs a patch. No intensity-per-output measure captures the downstream emissions that make Oil and Coal transition-exposed. Under EPA, Oil still does not make a five-industry Brown leg. (c05)
9. **Alternative indices end early.** CPU ends Sep 2025 and MCCC Jun 2025, leaving 10 and 13 holdout months uncovered. With the team's rule each also burns about 8 years at the start. (c06)
10. Not inferable from the prompt, listed for completeness: the missing Oct 2025 CPI print explains the whole purified-vs-original holdout gap (replication.md item 2).
\end{lstlisting}

\subsection{Exchange 02: coding support (a second implementation of the frozen test)}\label{app:ex02}

Files: \texttt{exchange/\allowbreak{}02\_\allowbreak{}coding\_\allowbreak{}support/\allowbreak{}}. Answered by: Claude (see the provenance note). Date: 2026-09-26. Purpose: coding support.

\begin{center}\fbox{\begin{minipage}{0.95\linewidth}\small
\begin{tabularx}{\linewidth}{@{}lX@{}}
\textbf{Asked.} & A pandas and statsmodels implementation of the frozen 1993--2009 test, one function per step with docstrings, a fixed seed, every silent reading as a parameter, and pre-run checks with stop rules. \\
\textbf{Came back.} & A 985-line script with 25 numbered assumptions and 13 pre-run checks with stop rules; after a follow-up naming six problems, six replaced functions, each with a unit test. \\
\textbf{Verified.} & Both versions run on the real data against my own implementation; with the two readings my prompt left open set to mine, all 15 coefficients and all 5,000 shuffle draws agree to 3.3e-13; mutation tests of its look-ahead audit. \\
\textbf{Kept.} & A second implementation that confirms the M8 verdict, and its null-seed harness, which showed that the pre-registered NW(6) t over-rejects (13 of 100 seeds) while the shuffle is correctly sized. \\
\end{tabularx}\end{minipage}}\end{center}

\paragraph{My evaluation}

\begin{quote}\small
The first reply wrote a correct backtest inside a flawed test harness. I never used it as the run of record: I ran it on the real data as a second implementation and compared every component with my own (M8), which reproduces the team pipeline bitwise.

\textbf{The core mathematics matched mine.} The 985-line script ran the first time, with only file paths and \texttt{fac} supplied, and reached the same verdict (FAIL, 1 of 4; alpha $-$0.19\%/yr, NW(6) t $-$0.25, vs. $-$0.18\%/yr, t $-$0.27). Two readings explain every numeric gap. With both switched to mine, it reproduces all 15 coefficients, the decomposition, the drop-one alphas and all 5,000 shuffle draws, to 3.3e-13 at worst. The timing conventions (lagged hedge, publication lag, extend-not-stack holds, costs paid in t+1) needed no correction.

\textbf{Three of the six problems came from my prompt.} It described only the FF5 \texttt{fac}, so the code hedged with FF5 SMB (up to 3.5 pp a month off the team FF3 SMB); the hedge betas moved by up to 0.09, and the reply's own check that reproduces the team trade stopped at a gap of 2.3e-3 a month. ``Keeps its sign'' let (iv) pass a negative alpha, and ``the seen 2010-2022 window'' let seen mode run to 2022-12, where five holdout months flip the seen alpha from $-$0.16\%/yr (t $-$0.24) to +0.26\%/yr (t 0.37). The other three are the reply's own: a broken instruction, a bug and a blind spot. It hard-coded a 60-calendar-month burn-in and disclosed it (assumption 11), although the prompt asked for every silent reading to be a parameter; the window start moved from 1994-03 to 1995-02, $\pi$ from 0.729 to 0.809 and the shuffle p from 0.478 to 0.536. Its synthetic GS10, a random walk clipped at 0.5, sat on the floor (76\% of 1995--2009 months for seed 0), so BOND went constant and its interaction with the in-position indicator became a plain in-position dummy that absorbed the planted effect: the positive control failed its own stop rule (4\% of the planted effect in alpha, t 0.98), and seed 16 crashed the null-seed check. Its look-ahead test compared end-of-month states only and missed the two return-side bugs among the six I injected.

\textbf{The follow-up fixed all six and broke one interface.} All six unit tests passed the first time without edits. The defaults reproduce my implementation to 3.3e-13 and the seen dry run ($-$0.12\%/yr, t $-$0.19), the positive control puts 97--104\% of the planted shift in alpha (t 3.5 to 6.1), and the new audit flags 7 of 8 injected timing bugs, including both earlier misses. But the new \texttt{attention\_signal} changed its output layout while the reply said everything else was unchanged; I sent the follow-up in a fresh chat without the first reply, so Claude wrote its six replacement functions without seeing the code they replace. A literal drop-in crashes, and the obvious one-line patch runs without error on a window 14 months too long (from 1993-01, n = 204, alpha $-$0.13\%/yr). It also told me to report a stale seen-window alpha ($-$0.16\%/yr, from its first version). Claude had no code execution in these chats (see the provenance note), so the numbers it quotes about its own code are not computed.

\textbf{Trust its formulas and audit its interfaces.} Where the spec was exact, the code was exact. The errors sat outside the formulas: in readings the spec left open, in test scaffolding, and at the seam between old and new functions. None changed the verdict or the frozen rule. The independent reference made the exchange useful: each discrepancy became a switch tied to specific lines, and mutation tests measured what its checks could see. Its own acceptance test for the fix (``should again match your implementation to 1e-13'') presumes that reference exists. The repaired harness also audited our own inference: on 100 null seeds the calendar shuffle is correctly sized (6 rejections at 5\%), but the pre-registered NW(6) t over-rejects (13), a bias that cannot rescue a negative alpha. Run against a reference and never trusted blind, Claude is a cheap second implementation for backtest code, not a verifier of it.

\end{quote}

{\raggedright\paragraph{Prompt (verbatim)} \texttt{prompt.md}, 997 words (957 excluding Markdown symbols).\par}

\begin{lstlisting}[style=transcript]
Write Python code for a pre-registered backtest I will run exactly once. You do not have the data, so report no results.

CONTEXT. This is my MFE 230GA final project at Berkeley. Our team's strategy shorts a factor-hedged "Brown" industry leg after attention spikes. It failed its 2022-2026 holdout, and we recommend "Do not implement". I froze one rule for 1993-2009, the one tradable window nobody has backtested. I will debug your code on synthetic data and the seen 2010-2022 window first, so the test window is a parameter and nothing is tuned.

INPUTS. Returns are monthly decimals.
1. `ff49_industry_monthly.csv`: a `date` column (month-end, e.g. 1926-07-31) and 49 Fama-French industry columns of value-weighted returns, NaN before an industry exists. Brown = Util, Ships, Aero, Steel, BldMt. The Brown leg R^B_t is their equal-weighted mean minus RF.
2. `fac`: a DataFrame parsed from Ken French's FF5 and momentum files, month-end DatetimeIndex from 1963-07, columns Mkt-RF, SMB, HML, RMW, CMA, RF, UMD.
3. FRED CSVs with columns `observation_date` (first of month; map to month-end) and the series id: EMVENRGYENVREG and EMVOVERALLEMV (monthly from 1985-01; the first has exact zeros, the second none), GS10 (monthly average 10-year yield, percent), MCOILWTICO (monthly average WTI price, from 1986-01), VIXCLS (daily closes from 1990-01-02, blank on holidays).

THE TEAM TRADE (validated, unchanged).
- Hedge: OLS of R^B on a constant and f = (Mkt-RF, SMB, HML) over months t-59 to t gives (a_t, b_t); residual e_t = R^B_t - a_{t-1} - b_{t-1}'f_t.
- Size: m_t = min(1, 0.05 / (sqrt(12) * sd_t)), sd_t = 36-month rolling std (ddof=1) of e through t.
- Position w_t = -m_t in a hold month, else 0, set at the end of month t; it earns w_t*R^B_{t+1} - w_t*b_t'f_{t+1}.
- Costs per unit of turnover: 10 bp on |Deltaw|, 5 bp on the Mkt-RF overlay change, 25 bp on SMB and HML overlay changes; a trade at the end of t pays in t+1.
- A crossing is an extreme month whose previous month was not extreme. The hold covers the crossing month and the next five. With the publication lag, the hold state at the end of t uses the extreme flag of t-1.
- Always-on benchmark R^AO: every month is a hold month.

THE FROZEN RULE.
1. Signal: s_t = EMVENRGYENVREG_t / EMVOVERALLEMV_t, observed one month late (the value for month t is usable at the end of month t+1).
2. Zeros: a zero month is missing, not low attention. It cannot trigger a position, and it is dropped from both the z-score window and the percentile history. The z-score needs at least 48 nonzero months in the trailing 60; if more than 10% of the trailing 60 months are zero, the signal is off.
3. Extreme month: z_t = 60-month rolling z-score of log(s_t) over nonzero months; extreme when z_t exceeds its expanding, past-only 80th percentile (team convention: at least 60 prior z observations).
4. Trade: the 6-month rule above. A new crossing during a hold extends the hold; positions never stack.
5. Sample: 1993-01 to 2009-12 (or from the first month the rule can fire, if later), run once.
6. Pass bar, all four required: (i) the timing alpha below is positive with t >= 2 and p from t(n-k); (ii) calendar-shuffle p <= 0.05 (move the same hold windows to random dates, 5,000 times, recompute alpha); (iii) at least 8 independent episodes (merge holds that touch or overlap); (iv) alpha keeps its sign when each Brown industry is dropped in turn (rebuild leg, hedge, sizing, R^AO and pi each time).

ATTRIBUTION. D_t = R^T_t - pi*R^AO_t, net of costs, with R^T the timed strategy and pi = mean |timed position| / mean |always-on position| over the test window. Regress D_t on FF5 + UMD, BOND, the WTI log return, dVIX_t (change in the monthly mean of daily VIX), dlog EMVOVERALLEMV_t, and I_{t-1} x {Mkt-RF, HML, BOND, dVIX}, where I_{t-1} = 1 if the strategy held a position entering month t. Alpha is the intercept. Newey-West with 6 lags (12 as a check); p from t(n-k), k = number of estimated parameters. Report mean(D) = alpha + sum beta*mean(F) + sum gamma*mean(I*F).

BOND is the monthly excess return of a constant-maturity 10-year Treasury from GS10 (percent to decimal): r_t = y_{t-1}/12 - D_{t-1}*(y_t - y_{t-1}) + 0.5*C_{t-1}*(y_t - y_{t-1})^2, with D and C the modified duration and convexity of a 10-year semiannual par bond at yield y_{t-1}, minus RF.

CONSTRAINTS.
- No look-ahead: every input used at the end of month t is known then (the signal a month later). Attribution regressors are contemporaneous by design; do not lag them.
- pandas, numpy and statsmodels only. `sm.OLS(...).fit(cov_type="HAC", cov_kwds={"maxlags": 6}, use_t=True)` matches our team's Newey-West t and gives p from t(n-k). Count k from the design matrix.
- One function per step, each with a docstring giving inputs, outputs and timing convention. Fixed seed, no parameter search, no plots.

Where the spec is silent, pick one reading, make it a parameter and say why. Three I see: whether a zero month between two extreme months makes the next one a new crossing; how the shuffle places windows and whether pi and I are recomputed per draw; what happens to a hold still open in 2009-12.

OUTPUTS, as DataFrames, also printed.
1. Pass/fail table: one row per component (i) to (iv), with statistic, bar, value, pass/fail; then the overall verdict.
2. Attribution: coefficient, NW(6) t, NW(12) t and p per regressor; n, k, R-squared, pi, window dates.
3. Decomposition: mean(D), alpha, every beta*mean(F) and gamma*mean(I*F) term, and the identity gap (zero to machine precision).
4. Diagnostics: first month the rule can fire, zero-month counts, crossing dates, episode list, leave-one-out alphas and t.
5. Shuffle: actual alpha, null quantiles, p.

AFTER THE CODE.
(a) Number every assumption you made beyond this spec.
(b) The checks you would run before my single real run, each with the result that would make you stop, including a look-ahead test.
\end{lstlisting}

{\raggedright\paragraph{Reply (verbatim; Claude)} \texttt{claude\_\allowbreak{}response.md}, 7,110 words (6,149 excluding Markdown symbols).\par}

\begin{lstlisting}[style=transcriptsmall]
Below is one self-contained script (pandas, numpy and statsmodels only). Each step is a separate function, and each docstring gives the step's inputs, outputs and timing. All choices live in a frozen `Config`. The test window is a parameter. Nothing is estimated from the test window except the attribution regression itself. The script also contains a synthetic-data generator, a look-ahead truncation test and a planted-effect positive control, so you can debug end to end before the real run.

I have not executed this code. Treat the checks in (b) as mandatory before the single real run.

```python
"""
brown_attention_prereg.py

Pre-registered, run-once backtest of the frozen Brown-attention rule (MFE 230GA).
pandas / numpy / statsmodels only. No plots, no parameter search, fixed seed.

Order of use
  1. MODE = "synthetic": debug end to end on random data (+ look-ahead test, planted effect).
  2. MODE = "seen":      2010-01..2022-12, compare with the team's existing numbers.
  3. MODE = "final":     1993-01..2009-12, exactly once, no code change after step 2.
"""
from __future__ import annotations

from dataclasses import asdict, dataclass, replace

import numpy as np
import pandas as pd
import statsmodels.api as sm

pd.set_option("display.width", 220)
pd.set_option("display.max_columns", 40)
pd.set_option("display.max_rows", 500)
pd.set_option("display.float_format", "{:.6g}".format)

BROWN = ("Util", "Ships", "Aero", "Steel", "BldMt")


# =============================================================================
# 0. Frozen configuration
# =============================================================================
@dataclass(frozen=True)
class Config:
    """Every number the test uses. Nothing here is estimated or tuned."""
    # test window in RETURN months (inclusive); start moves later if the rule cannot fire yet
    test_start: str = "1993-01"
    test_end: str = "2009-12"
    brown: tuple = BROWN
    hedge_factors: tuple = ("Mkt-RF", "SMB", "HML")
    # team trade (validated, unchanged)
    hedge_window: int = 60
    vol_window: int = 36
    vol_target: float = 0.05
    cost_w: float = 0.0010
    cost_mkt: float = 0.0005
    cost_smb: float = 0.0025
    cost_hml: float = 0.0025
    hold_months: int = 6
    pub_lag: int = 1
    # frozen signal
    z_window: int = 60
    z_min_nonzero: int = 48
    max_zero_share: float = 0.10
    pct: float = 80.0
    min_pct_history: int = 60
    # attribution
    attr_factors: tuple = ("Mkt-RF", "SMB", "HML", "RMW", "CMA", "UMD",
                           "BOND", "WTI", "dVIX", "dlogEMV")
    interact: tuple = ("Mkt-RF", "HML", "BOND", "dVIX")
    nw_lags: int = 6
    nw_lags_check: int = 12
    # pass bars
    t_bar: float = 2.0
    p_bar: float = 0.05
    min_episodes: int = 8
    # readings of the points the spec leaves open -- see notes (a)1-(a)3
    bridge_missing: bool = True       # extreme, zero, extreme = one run (no new crossing)
    bridge_limit: int = 6             # ... unless more than 6 consecutive missing months
    end_policy: str = "truncate"      # "truncate" or "liquidate" a hold still open at test_end
    shuffle_draws: int = 5000
    shuffle_min_gap: int = 1          # shuffled episodes never touch, so K is preserved
    shuffle_recompute_pi_I: bool = True
    seed: int = 230


# =============================================================================
# Calendar / IO helpers
# =============================================================================
def _me(x) -> pd.Timestamp:
    """Month-end timestamp (midnight) of any date-like."""
    return (pd.Timestamp(x) + pd.offsets.MonthEnd(0)).normalize()


def _to_month_end(values) -> pd.DatetimeIndex:
    """Vectorised month-end mapping; FRED first-of-month dates map to the same month's end."""
    return (pd.DatetimeIndex(pd.to_datetime(values)) + pd.offsets.MonthEnd(0)).normalize()


def _month_range(a, b) -> pd.DatetimeIndex:
    return pd.date_range(_me(a), _me(b), freq=pd.offsets.MonthEnd())


def _shift_months(ts, k: int) -> pd.Timestamp:
    return (_me(ts) + pd.offsets.MonthEnd(k)).normalize()


def _position_months(window: pd.DatetimeIndex) -> pd.DatetimeIndex:
    """Position months whose positions earn inside the window (return month t <- position t-1)."""
    return _month_range(_shift_months(window[0], -1), _shift_months(window[-1], -1))


def _assert_monthly(idx: pd.DatetimeIndex, name: str) -> None:
    if idx.has_duplicates:
        raise ValueError(f"{name}: duplicate months")
    if len(idx) != len(_month_range(idx.min(), idx.max())):
        raise ValueError(f"{name}: gaps in the monthly index")


def _assert_decimal(frame: pd.DataFrame, name: str) -> None:
    q = np.nanpercentile(np.abs(frame.to_numpy(float)), 99)
    if q > 1.0:
        raise ValueError(f"{name}: 99th pct |return| = {q:.2f}; looks like percent, not decimals")


def _read(src) -> pd.DataFrame:
    return src.copy() if isinstance(src, pd.DataFrame) else pd.read_csv(src)


def _show(title: str, obj) -> None:
    print("\n" + "=" * 100 + f"\n{title}\n" + "=" * 100)
    print(obj.to_string() if hasattr(obj, "to_string") else obj)


# =============================================================================
# 1. Loaders
# =============================================================================
def load_ff49(src) -> pd.DataFrame:
    """
    Step 1a. Fama-French 49 industries.
    Input : path or DataFrame with `date` (month-end) + 49 industry columns, monthly decimals.
    Output: DataFrame on a sorted, gap-free month-end index; NaN where an industry does not
            exist; values <= -0.99 (French's -99.99 code) set to NaN.
    Timing: row t = value-weighted return over calendar month t.
    """
    df = _read(src)
    df.columns = [str(c).strip() for c in df.columns]
    idx = _to_month_end(df.pop("date"))
    df = df.apply(pd.to_numeric, errors="coerce")
    df.index = idx
    df = df.sort_index().mask(lambda x: x <= -0.99)
    _assert_monthly(df.index, "ff49")
    _assert_decimal(df, "ff49")
    return df


def load_factors(fac: pd.DataFrame) -> pd.DataFrame:
    """
    Step 1b. French FF5 + momentum.
    Input : DataFrame with Mkt-RF, SMB, HML, RMW, CMA, RF, UMD (decimals), month-end index.
    Output: those columns on a sorted, gap-free month-end index.
    Timing: row t = return over month t.
    """
    cols = ["Mkt-RF", "SMB", "HML", "RMW", "CMA", "RF", "UMD"]
    out = fac[cols].astype(float).copy()
    out.index = _to_month_end(out.index)
    out = out.sort_index()
    _assert_monthly(out.index, "fac")
    _assert_decimal(out, "fac")
    return out


def load_fred_monthly(src, series_id: str) -> pd.Series:
    """
    Step 1c. Monthly FRED series (EMV trackers, GS10, MCOILWTICO).
    Input : path or DataFrame with `observation_date` (first of month) and `series_id`.
    Output: Series on a complete month-end index; '.'/blank -> NaN; exact zeros kept.
    Timing: the value dated month t describes month t (publication lag handled downstream).
    """
    df = _read(src)
    s = pd.Series(pd.to_numeric(df[series_id], errors="coerce").to_numpy(float),
                  index=_to_month_end(df["observation_date"]), name=series_id).sort_index()
    if s.index.has_duplicates:
        raise ValueError(f"{series_id}: more than one observation per month")
    return s.reindex(_month_range(s.index.min(), s.index.max()))


def load_vix_monthly(src, series_id: str = "VIXCLS") -> pd.Series:
    """
    Step 1d. Monthly mean of daily VIX closes.
    Input : path or DataFrame with daily `observation_date` and VIXCLS (blank on holidays).
    Output: Series on a complete month-end index = mean of the non-missing closes in month t.
    Timing: known at the end of month t.
    """
    df = _read(src)
    d = pd.Series(pd.to_numeric(df[series_id], errors="coerce").to_numpy(float),
                  index=pd.DatetimeIndex(pd.to_datetime(df["observation_date"]))).dropna()
    m = d.groupby(_to_month_end(d.index)).mean()
    m.index = pd.DatetimeIndex(m.index)
    return m.reindex(_month_range(m.index.min(), m.index.max())).rename("VIX")


# =============================================================================
# 2-4. Brown leg, hedge, sizing
# =============================================================================
def brown_leg(ff49: pd.DataFrame, rf: pd.Series, industries) -> pd.Series:
    """
    Step 2. Brown leg.
    Inputs : industry returns, RF, list of Brown industries.
    Output : R^B_t = equal-weighted mean of the listed industries available in t, minus RF_t.
    Timing : month-t return.
    """
    missing = [c for c in industries if c not in ff49.columns]
    if missing:
        raise KeyError(f"industries not found in FF49 file: {missing}")
    leg = ff49[list(industries)].mean(axis=1, skipna=True)
    return (leg - rf.reindex(leg.index)).rename("RB")


def rolling_hedge(rb: pd.Series, f: pd.DataFrame, window: int = 60):
    """
    Step 3. Rolling factor hedge.
    Inputs : R^B_t and hedge factors f_t = (Mkt-RF, SMB, HML) on the same complete index.
    Output : a (Series) and b (DataFrame, columns = f.columns).
    Timing : (a_t, b_t) = OLS of R^B on [1, f] over months t-window+1..t, i.e. data through t,
             known at the end of t. NaN unless all `window` months are complete.
    """
    idx = rb.index
    X = np.column_stack([np.ones(len(idx)), f.reindex(idx).to_numpy(float)])
    y = rb.to_numpy(float)
    coef = np.full((len(idx), X.shape[1]), np.nan)
    for i in range(window - 1, len(idx)):
        Xi, yi = X[i - window + 1:i + 1], y[i - window + 1:i + 1]
        if np.isfinite(Xi).all() and np.isfinite(yi).all():
            coef[i] = np.linalg.lstsq(Xi, yi, rcond=None)[0]
    a = pd.Series(coef[:, 0], index=idx, name="a")
    b = pd.DataFrame(coef[:, 1:], index=idx, columns=list(f.columns))
    return a, b


def hedge_residual_and_size(rb: pd.Series, f: pd.DataFrame, a: pd.Series, b: pd.DataFrame,
                            cfg: Config):
    """
    Step 4. Out-of-sample residual and position size.
    Inputs : R^B, f, (a, b) from step 3.
    Output : e_t = R^B_t - a_{t-1} - b_{t-1}'f_t;
             sd_t = std(e_{t-35..t}, ddof=1), all 36 required;
             m_t = min(1, vol_target / (sqrt(12) sd_t)).
    Timing : e_t, sd_t, m_t use data through t, known at the end of t.
    """
    fitted = a.shift(1) + (b.shift(1) * f[b.columns]).sum(axis=1, min_count=b.shape[1])
    e = (rb - fitted).rename("e")
    sd = e.rolling(cfg.vol_window, min_periods=cfg.vol_window).std(ddof=1).rename("sd")
    m = np.minimum(1.0, cfg.vol_target / (np.sqrt(12.0) * sd)).rename("m")
    return e, sd, m


# =============================================================================
# 5-7. Signal, crossings, hold
# =============================================================================
def attention_signal(emv_cat: pd.Series, emv_all: pd.Series, cfg: Config) -> pd.DataFrame:
    """
    Step 5. Frozen signal, indexed by SIGNAL month tau.
    Inputs : EMVENRGYENVREG and EMVOVERALLEMV (monthly, month-end index).
    Output : DataFrame with
        s            EMVENRGYENVREG / EMVOVERALLEMV
        missing      s is zero (or NaN): missing, never "low attention"
        n_missing60  missing months in tau-59..tau (NaN until 60 calendar months exist)
        off          month is nonzero but > 10% of the trailing 60 months are missing
        z            (log s_tau - mean) / sd over the NONZERO months of tau-59..tau
                     (needs >= 48 nonzero and <= 6 missing; NaN otherwise)
        thr          80th percentile of all valid z strictly before tau (needs >= 60)
        extreme      1 if z > thr, 0 if not, NaN if z or thr undefined
    Timing : row tau uses data dated <= tau only, but tau's data are published at the end of
             tau+1, so row tau may drive a position no earlier than the end of tau+pub_lag.
             Missing months are dropped from the z window and the percentile history.
    """
    start = max(emv_cat.first_valid_index(), emv_all.first_valid_index())
    end = min(emv_cat.last_valid_index(), emv_all.last_valid_index())
    idx = _month_range(start, end)
    s = emv_cat.reindex(idx) / emv_all.reindex(idx)
    missing = ~(s > 0)                                   # exact zero or NaN
    W = cfg.z_window
    max_missing = int(np.floor(cfg.max_zero_share * W + 1e-9))   # = 6
    logs = np.log(s.where(~missing))
    n_miss = missing.astype(float).rolling(W, min_periods=W).sum()
    n_nz = W - n_miss
    mu = logs.rolling(W, min_periods=cfg.z_min_nonzero).mean()
    sd = logs.rolling(W, min_periods=cfg.z_min_nonzero).std(ddof=1)
    ok = (~missing) & (n_nz >= cfg.z_min_nonzero) & (n_miss <= max_missing) & (sd > 0)
    z = ((logs - mu) / sd).where(ok)

    zv = z.to_numpy(float)
    thr = np.full(len(zv), np.nan)
    hist: list = []
    for i, zi in enumerate(zv):
        if len(hist) >= cfg.min_pct_history:
            thr[i] = np.percentile(hist, cfg.pct)        # past-only: z_tau not yet appended
        if np.isfinite(zi):
            hist.append(zi)
    defined = np.isfinite(zv) & np.isfinite(thr)
    extreme = np.where(defined, (np.nan_to_num(zv) > np.nan_to_num(thr)).astype(float), np.nan)

    return pd.DataFrame({
        "s": s, "missing": missing, "n_missing60": n_miss,
        "off": (~missing) & (n_miss > max_missing),
        "z": z, "thr": thr, "extreme": extreme}, index=idx)


def crossings(extreme: pd.Series, cfg: Config) -> pd.Series:
    """
    Step 6. Crossings.
    Input  : extreme flag (1/0/NaN) by signal month.
    Output : bool Series, True in signal month tau if extreme_tau == 1 and the previous month
             was not extreme. With bridge_missing, a missing flag (zero or off month) carries
             the last defined flag forward for up to bridge_limit months, so
             extreme -> zero -> extreme is ONE run. Undefined previous state = not extreme.
    Timing : uses flags dated <= tau; usable at the end of tau+pub_lag.
    """
    prev = extreme.ffill(limit=cfg.bridge_limit) if cfg.bridge_missing else extreme
    prev = prev.shift(1)
    return ((extreme == 1) & (prev != 1)).rename("crossing")


def hold_positions(cross: pd.Series, index: pd.DatetimeIndex, cfg: Config) -> pd.Series:
    """
    Step 7. Hold state in position time.
    Inputs : crossings (signal time) and the master monthly index.
    Output : bool Series on `index`; True = a short position is SET at the end of month t.
    Timing : hold in signal time covers tau..tau+5 after each crossing (a new crossing extends,
             positions never stack); with the publication lag the hold at the end of t uses
             signal month t-1: hold_pos_t = hold_sig_{t-pub_lag}. So a crossing in tau gives
             positions set at the end of tau+1..tau+6, earning tau+2..tau+7.
    """
    c = cross.reindex(index, fill_value=False).astype(float)
    hold_sig = c.rolling(cfg.hold_months, min_periods=1).max() > 0
    return hold_sig.shift(cfg.pub_lag, fill_value=False).astype(bool).rename("hold")


# =============================================================================
# 8. Trade engine (timed and always-on)
# =============================================================================
def run_trade(rb: pd.Series, f: pd.DataFrame, b: pd.DataFrame, m: pd.Series,
              hold: pd.Series, cfg: Config) -> pd.DataFrame:
    """
    Step 8. Team trade.
    Inputs : R^B_t, hedge factors, b_t, m_t, hold (bool, set at end of t).
    Output : DataFrame on rb.index:
             hold, m, w (w_t = -m_t if hold else 0, set at end of t),
             h_<factor> (overlay = -w_t b_t, set at end of t), trade_cost (trade at end of t),
             w_prev = w_{t-1}, hedged_t = R^B_t - b_{t-1}'f_t, gross_t = w_{t-1} hedged_t,
             cost_t = trade_cost_{t-1} (paid in t), net_t = gross_t - cost_t.
    Timing : row t's net return is earned in month t by the position set at the end of t-1.
    """
    idx = rb.index
    hf = list(cfg.hedge_factors)
    f = f.reindex(idx)[hf]
    b = b.reindex(idx)[hf]
    hold = hold.reindex(idx, fill_value=False).astype(bool)
    mv = m.reindex(idx).to_numpy(float)
    w = np.where(hold.to_numpy(), -mv, 0.0)
    bv = b.to_numpy(float)
    h = np.where(w[:, None] == 0.0, 0.0, -w[:, None] * bv)
    dw = np.abs(np.diff(w, prepend=0.0))
    dh = np.abs(np.diff(h, axis=0, prepend=np.zeros((1, h.shape[1]))))
    unit = np.array([cfg.cost_mkt, cfg.cost_smb, cfg.cost_hml])
    trade_cost = cfg.cost_w * dw + dh @ unit

    out = pd.DataFrame({"hold": hold, "m": mv, "w": w}, index=idx)
    for j, c in enumerate(hf):
        out[f"h_{c}"] = h[:, j]
    out["trade_cost"] = trade_cost
    out["w_prev"] = out["w"].shift(1)
    out["hedged"] = rb - (b.shift(1) * f).sum(axis=1, min_count=len(hf))
    out["gross"] = out["w_prev"] * out["hedged"]
    out["cost"] = out["trade_cost"].shift(1)
    out["net"] = out["gross"] - out["cost"]
    return out


def leg_pipeline(ff49: pd.DataFrame, fac: pd.DataFrame, industries, hold: pd.Series,
                 index: pd.DatetimeIndex, cfg: Config) -> dict:
    """
    Steps 2-4 + 8 for one set of Brown industries (used for the main run and for each
    leave-one-out rebuild).
    Output : dict with rb, f, a, b, e, sd, m, tr_T (timed book), tr_AO (always-on book).
    Timing : as in steps 2-4 and 8; the hold series is taken as given.
    """
    f = fac.reindex(index)[list(cfg.hedge_factors)]
    rb = brown_leg(ff49.reindex(index), fac["RF"].reindex(index), industries)
    a, b = rolling_hedge(rb, f, cfg.hedge_window)
    e, sd, m = hedge_residual_and_size(rb, f, a, b, cfg)
    tr_T = run_trade(rb, f, b, m, hold, cfg)
    tr_AO = run_trade(rb, f, b, m, pd.Series(True, index=index), cfg)
    return dict(rb=rb, f=f, a=a, b=b, e=e, sd=sd, m=m, tr_T=tr_T, tr_AO=tr_AO)


# =============================================================================
# 9-10. Attribution regressors
# =============================================================================
def par_bond_duration_convexity(y, maturity: float = 10.0, freq: int = 2):
    """
    Step 9a. Modified duration (years) and convexity (years^2) of a par bond.
    Input : annual yield(s) y in decimals, semiannual compounding and coupons.
    Output: (D, C) arrays, from exact cash flows.
    """
    y = np.asarray(y, dtype=float).reshape(-1, 1)
    n = int(round(maturity * freq))
    i = np.arange(1, n + 1, dtype=float)[None, :]
    g = 1.0 + y / freq
    cf = np.repeat(y / freq, n, axis=1)
    cf[:, -1] += 1.0
    disc = g ** (-i)
    price = (cf * disc).sum(axis=1)
    dur = (cf * (i / freq) * disc).sum(axis=1) / g[:, 0] / price
    conv = (cf * (i / freq) * ((i + 1) / freq) * disc).sum(axis=1) / g[:, 0] ** 2 / price
    return dur, conv


def bond_excess_return(gs10: pd.Series, rf: pd.Series) -> pd.Series:
    """
    Step 9b. BOND_t = y_{t-1}/12 - D_{t-1} dy_t + 0.5 C_{t-1} dy_t^2 - RF_t,
    y = GS10/100, dy_t = y_t - y_{t-1}, D and C of a 10y semiannual par bond at y_{t-1}.
    Timing: month-t excess return (contemporaneous regressor).
    """
    y = (gs10 / 100.0).astype(float)
    y_prev = y.shift(1)
    dy = y - y_prev
    dur, conv = par_bond_duration_convexity(y_prev.to_numpy())
    r = y_prev / 12.0 - dur * dy + 0.5 * conv * dy ** 2
    return (r - rf.reindex(r.index)).rename("BOND")


def attribution_factors(fac: pd.DataFrame, bond: pd.Series, wti: pd.Series, vix_m: pd.Series,
                        emv_all: pd.Series, index: pd.DatetimeIndex) -> pd.DataFrame:
    """
    Step 10. Month-t regressors, contemporaneous by design (NOT lagged):
    FF5 + UMD, BOND_t, WTI_t = log(P_t / P_{t-1}), dVIX_t = VIXbar_t - VIXbar_{t-1},
    dlogEMV_t = log EMVOVERALLEMV_t - log EMVOVERALLEMV_{t-1}.
    All differences are taken on complete monthly calendars before reindexing.
    """
    F = fac.reindex(index)[["Mkt-RF", "SMB", "HML", "RMW", "CMA", "UMD"]].copy()
    F["BOND"] = bond.reindex(index)
    F["WTI"] = np.log(wti).diff().reindex(index)
    F["dVIX"] = vix_m.diff().reindex(index)
    F["dlogEMV"] = np.log(emv_all).diff().reindex(index)
    return F


# =============================================================================
# 11. Window, timing difference D
# =============================================================================
def test_window(sig: pd.DataFrame, index: pd.DatetimeIndex, cfg: Config):
    """
    Step 11a. Test window.
    Output : (first_signal, first_position, first_return, window).
             first_signal   = first signal month with a defined flag (valid z, >= 60 prior z)
             first_position = first month-end a position can be set = first_signal + pub_lag
             first_return   = first return month that position affects = first_position + 1
             window         = RETURN months [max(test_start, first_return), test_end]
    """
    elig = sig["extreme"].dropna()
    if elig.empty:
        raise ValueError("the signal never produces a defined flag")
    first_sig = elig.index[0]
    first_pos = _shift_months(first_sig, cfg.pub_lag)
    first_ret = _shift_months(first_pos, 1)
    start, end = max(_me(cfg.test_start), first_ret), _me(cfg.test_end)
    if end > index[-1] or start > end:
        raise ValueError(f"window {start:%Y-%m}..{end:%Y-%m} not covered by the data")
    return first_sig, first_pos, first_ret, _month_range(start, end)


def _liquidation_cost(tr: pd.DataFrame, t_end: pd.Timestamp, cfg: Config) -> float:
    """Cost of closing, at the end of t_end, the book set at the end of t_end - 1."""
    prev = _shift_months(t_end, -1)
    h = np.abs(tr.loc[prev, [f"h_{c}" for c in cfg.hedge_factors]].to_numpy(float))
    unit = np.array([cfg.cost_mkt, cfg.cost_smb, cfg.cost_hml])
    return float(cfg.cost_w * abs(tr.loc[prev, "w"]) + h @ unit)


def timing_difference(tr_T: pd.DataFrame, tr_AO: pd.DataFrame, window: pd.DatetimeIndex,
                      cfg: Config, pi: float | None = None):
    """
    Step 11b. D_t = R^T_t - pi R^AO_t over the window's return months, net of costs.
    pi = mean|w^T_{t-1}| / mean|w^AO_{t-1}| over the same months (positions that earn there),
    unless pi is passed in. end_policy="liquidate" charges both books' closing cost in the last
    window month; "truncate" does not (that trade would be paid after the window).
    Output : D (Series), pi (float), I_prev (0/1: position held entering month t).
    """
    T, A = tr_T.loc[window], tr_AO.loc[window]
    cols = ["w_prev", "gross", "cost", "net"]
    if not (np.isfinite(T[cols].to_numpy(float)).all() and np.isfinite(A[cols].to_numpy(float)).all()):
        raise ValueError("non-finite strategy values inside the test window")
    if pi is None:
        pi = T["w_prev"].abs().mean() / A["w_prev"].abs().mean()
    RT, RA = T["net"].copy(), A["net"].copy()
    if cfg.end_policy == "liquidate":
        RT.iloc[-1] -= _liquidation_cost(tr_T, window[-1], cfg)
        RA.iloc[-1] -= _liquidation_cost(tr_AO, window[-1], cfg)
    elif cfg.end_policy != "truncate":
        raise ValueError(f"unknown end_policy {cfg.end_policy}")
    D = (RT - pi * RA).rename("D")
    I_prev = tr_T["hold"].shift(1, fill_value=False).loc[window].astype(float).rename("I_prev")
    return D, float(pi), I_prev


# =============================================================================
# 12. Attribution regression and decomposition
# =============================================================================
def design_matrix(F: pd.DataFrame, I_prev: pd.Series, cfg: Config) -> pd.DataFrame:
    """
    Step 12a. X_t = [const, FF5+UMD, BOND, WTI, dVIX, dlogEMV, I_{t-1} x {Mkt-RF, HML, BOND, dVIX}].
    No I_{t-1} main effect (the spec lists only the interactions). k = X.shape[1] = 15.
    """
    F = F.loc[I_prev.index, list(cfg.attr_factors)]
    X = F.copy()
    for c in cfg.interact:
        X[f"I*{c}"] = I_prev.to_numpy() * F[c].to_numpy()
    return sm.add_constant(X, has_constant="add")


def attribution_regression(D: pd.Series, X: pd.DataFrame, cfg: Config):
    """
    Step 12b. OLS of D on X; Newey-West (Bartlett) HAC with nw_lags and nw_lags_check.
    use_t=True -> t statistics and p-values from t(n-k), k counted from X.
    Output : coefficient table, NW(6) fit, NW(12) fit. Raises on missing data or rank deficiency.
    """
    bad = X.isna().any(axis=1) | D.isna()
    if bad.any():
        raise ValueError(f"missing regressors/returns in {list(X.index[bad.to_numpy()].strftime('%Y-%m'))}")
    if np.linalg.matrix_rank(X.to_numpy(float)) < X.shape[1]:
        raise ValueError("design matrix is rank deficient (e.g. I_{t-1} never or always 1)")
    f6 = sm.OLS(D, X).fit(cov_type="HAC", cov_kwds={"maxlags": cfg.nw_lags}, use_t=True)
    f12 = sm.OLS(D, X).fit(cov_type="HAC", cov_kwds={"maxlags": cfg.nw_lags_check}, use_t=True)
    n, k = X.shape
    if int(round(f6.df_resid)) != n - k:
        raise RuntimeError("df_resid != n - k")
    tab = pd.DataFrame({
        "coef": f6.params,
        f"t NW({cfg.nw_lags})": f6.tvalues,
        f"p NW({cfg.nw_lags})": f6.pvalues,
        f"t NW({cfg.nw_lags_check})": f12.tvalues,
        f"p NW({cfg.nw_lags_check})": f12.pvalues,
    }).rename(index={"const": "alpha"})
    return tab, f6, f12


def decompose_mean(D: pd.Series, X: pd.DataFrame, fit) -> pd.DataFrame:
    """
    Step 12c. mean(D) = alpha + sum beta_j mean(F_j) + sum gamma_j mean(I F_j).
    Exact for OLS with a constant (residuals sum to zero); the gap is reported.
    """
    contrib = fit.params * X.mean()
    rows = {"mean(D)": D.mean(), "alpha": contrib["const"]}
    for c in X.columns.drop("const"):
        rows[f"gamma*mean({c})" if c.startswith("I*") else f"beta*mean({c})"] = contrib[c]
    rows["sum of terms"] = contrib.sum()
    rows["identity gap"] = D.mean() - contrib.sum()
    return pd.Series(rows, name="value").to_frame()


# =============================================================================
# 13. Episodes and crossings
# =============================================================================
def hold_episodes(hold: pd.Series, pos_months: pd.DatetimeIndex) -> pd.DataFrame:
    """
    Step 13a. Independent episodes.
    Inputs : hold (bool, position time) and the position months that earn inside the window.
    Output : one row per maximal run of consecutive hold months (holds that touch or overlap
             merge), with position/return dates and whether the run was already open before,
             or is still open after, the window.
    """
    cols = ["position_from", "position_to", "months", "return_from", "return_to",
            "open_before_window", "open_after_window"]
    h = hold.reindex(pos_months, fill_value=False).to_numpy(bool)
    before = bool(hold.get(_shift_months(pos_months[0], -1), False))
    after = bool(hold.get(_shift_months(pos_months[-1], 1), False))
    rows, i = [], 0
    while i < len(h):
        if not h[i]:
            i += 1
            continue
        j = i
        while j + 1 < len(h) and h[j + 1]:
            j += 1
        rows.append({"position_from": pos_months[i], "position_to": pos_months[j],
                     "months": j - i + 1,
                     "return_from": _shift_months(pos_months[i], 1),
                     "return_to": _shift_months(pos_months[j], 1),
                     "open_before_window": i == 0 and before,
                     "open_after_window": j == len(h) - 1 and after})
        i = j + 1
    return pd.DataFrame(rows, columns=cols)


def crossing_table(sig: pd.DataFrame, cross: pd.Series, pos_months: pd.DatetimeIndex,
                   cfg: Config) -> pd.DataFrame:
    """
    Step 13b. Crossings whose hold overlaps the window's position months.
    Signal month tau -> positions set at the end of tau+pub_lag .. tau+pub_lag+hold_months-1.
    """
    cols = ["signal_month", "z", "threshold", "position_from", "position_to"]
    rows = []
    for tau in cross.index[cross.to_numpy(bool)]:
        p0 = _shift_months(tau, cfg.pub_lag)
        p1 = _shift_months(p0, cfg.hold_months - 1)
        if p1 < pos_months[0] or p0 > pos_months[-1]:
            continue
        rows.append({"signal_month": tau, "z": sig.at[tau, "z"], "threshold": sig.at[tau, "thr"],
                     "position_from": p0, "position_to": p1})
    return pd.DataFrame(rows, columns=cols)


def diagnostics_summary(sig, hold, leg, ff49, window, first_sig, first_pos, first_ret,
                        eps, xtab, pi, cfg: Config) -> pd.DataFrame:
    """
    Step 13c. Key facts about the run: fire dates, zero counts, hold months, pi, settings.
    Signal months "feeding" the window = the latest flags available at each position month.
    """
    pos_months = _position_months(window)
    sig_months = _month_range(_shift_months(pos_months[0], -cfg.pub_lag),
                              _shift_months(pos_months[-1], -cfg.pub_lag))
    sw = sig.reindex(sig_months)
    trw = leg["tr_T"].loc[pos_months]
    hist = _month_range(_shift_months(window[0], -(cfg.hedge_window + cfg.vol_window + 1)), window[-1])
    n_brown = int(ff49.reindex(hist)[list(cfg.brown)].notna().sum(axis=1).min())
    info = {
        "first signal month with a defined flag": f"{first_sig:%Y-%m}",
        "first month a position can be set (end of)": f"{first_pos:%Y-%m}",
        "first return month the rule can affect": f"{first_ret:%Y-%m}",
        "test window (return months)": f"{window[0]:%Y-%m} to {window[-1]:%Y-%m}",
        "n return months": len(window),
        "EMVENRGYENVREG zero months, whole file": int((sig["s"] == 0).sum()),
        "signal months feeding the window": f"{sig_months[0]:%Y-%m} to {sig_months[-1]:%Y-%m}",
        "  zero months among them": int((sw["s"] == 0).sum()),
        "  signal-off months (>10% zeros in trailing 60)": int(sw["off"].eq(True).sum()),
        "  months with a defined flag": int(sw["extreme"].notna().sum()),
        "  extreme months": int((sw["extreme"] == 1).sum()),
        "crossings affecting the window": len(xtab),
        "hold months (position months in window)": int(trw["hold"].sum()),
        "independent episodes": len(eps),
        "pi": pi,
        "mean m_t in hold months": float(trw.loc[trw["hold"], "m"].mean()) if trw["hold"].any() else np.nan,
        "min # Brown industries available (hedge history + window)": n_brown,
    }
    info.update({f"cfg.{k}": v for k, v in asdict(cfg).items()})
    return pd.Series(info, name="value").to_frame()


# =============================================================================
# 14. Calendar shuffle
# =============================================================================
def place_blocks(lengths, n: int, min_gap: int, rng: np.random.Generator) -> np.ndarray:
    """
    Step 14a. Uniform random placement of blocks with the given lengths in n slots, in random
    order, with >= min_gap empty slots between consecutive blocks (stars and bars).
    Output: bool array of length n.
    """
    lengths = np.asarray(lengths, dtype=int)
    K = len(lengths)
    free = n - int(lengths.sum()) - min_gap * (K - 1)
    if free < 0:
        raise ValueError("episodes do not fit in the window with the required gaps")
    lens = lengths[rng.permutation(K)]
    slots = np.sort(rng.choice(free + K, size=K, replace=False))
    starts = (slots - np.arange(K)
              + np.concatenate(([0], np.cumsum(lens)[:-1]))
              + min_gap * np.arange(K))
    out = np.zeros(n, dtype=bool)
    for s, L in zip(starts, lens):
        out[s:s + L] = True
    return out


def calendar_shuffle(alpha_actual: float, leg: dict, hold: pd.Series, F: pd.DataFrame,
                     window: pd.DatetimeIndex, pi_actual: float, I_actual: pd.Series,
                     cfg: Config):
    """
    Step 14b. Calendar-shuffle null.
    Inputs : actual alpha, the Brown pipeline `leg`, actual hold (position time), regressors,
             window, actual pi and I (used only if shuffle_recompute_pi_I is False).
    Method : the K actual episodes (maximal hold runs among the position months that earn in
             the window) are placed at random (place_blocks) inside those position months;
             months before the window keep their actual state. Each draw rebuilds the timed
             book (w, overlay, costs); pi and I_{t-1} are recomputed if shuffle_recompute_pi_I;
             alpha is re-estimated by OLS (point estimate only).
    Output : (summary table, array of null alphas);
             p = (1 + #{alpha_null >= alpha_actual}) / (1 + draws), one-sided.
    """
    pos_months = _position_months(window)
    lengths = hold_episodes(hold, pos_months)["months"].to_numpy(int)
    if len(lengths) == 0:
        raise ValueError("no episodes to shuffle")
    loc = hold.index.get_indexer(pos_months)
    if (loc < 0).any():
        raise ValueError("position months outside the hold index")
    rng = np.random.default_rng(cfg.seed)
    base = hold.to_numpy(bool)
    Fw = F.loc[window]
    null = np.empty(cfg.shuffle_draws)
    for d in range(cfg.shuffle_draws):
        hv = base.copy()
        hv[loc] = place_blocks(lengths, len(pos_months), cfg.shuffle_min_gap, rng)
        tr = run_trade(leg["rb"], leg["f"], leg["b"], leg["m"], pd.Series(hv, index=hold.index), cfg)
        if cfg.shuffle_recompute_pi_I:
            D, _, I = timing_difference(tr, leg["tr_AO"], window, cfg)
        else:
            D, _, _ = timing_difference(tr, leg["tr_AO"], window, cfg, pi=pi_actual)
            I = I_actual
        X = design_matrix(Fw, I, cfg).to_numpy(float)
        null[d] = np.linalg.lstsq(X, D.to_numpy(float), rcond=None)[0][0]
    p = (1 + np.sum(null >= alpha_actual)) / (1 + cfg.shuffle_draws)
    q = [0.01, 0.05, 0.10, 0.50, 0.90, 0.95, 0.99]
    vals = ([alpha_actual, null.mean(), null.std(ddof=1)] + list(np.quantile(null, q))
            + [p, cfg.shuffle_draws, len(lengths), int(lengths.sum()), cfg.seed])
    names = (["actual alpha", "null mean", "null sd"] + [f"null q{int(round(x * 100)):02d}" for x in q]
             + ["p (one-sided)", "draws", "episodes K", "hold months", "seed"])
    return pd.DataFrame({"value": vals}, index=names), null


# =============================================================================
# 15. Leave one industry out
# =============================================================================
def leave_one_out(ff49: pd.DataFrame, fac: pd.DataFrame, hold: pd.Series, F: pd.DataFrame,
                  window: pd.DatetimeIndex, index: pd.DatetimeIndex, alpha_full: float,
                  cfg: Config) -> pd.DataFrame:
    """
    Step 15. Drop each Brown industry in turn; rebuild leg, hedge, sizing, timed and always-on
    books, pi and D; refit the attribution. Hold timing is unchanged (the signal does not use
    returns). Output: one row per dropped industry with alpha, NW(6) t, p, pi, same_sign.
    """
    rows = []
    for drop in cfg.brown:
        keep = tuple(c for c in cfg.brown if c != drop)
        leg = leg_pipeline(ff49, fac, keep, hold, index, cfg)
        D, pi, I = timing_difference(leg["tr_T"], leg["tr_AO"], window, cfg)
        _, f6, _ = attribution_regression(D, design_matrix(F, I, cfg), cfg)
        a = float(f6.params["const"])
        rows.append({"dropped": drop, "alpha": a, f"t NW({cfg.nw_lags})": float(f6.tvalues["const"]),
                     "p": float(f6.pvalues["const"]), "pi": pi,
                     "same_sign": bool(np.sign(a) == np.sign(alpha_full))})
    return pd.DataFrame(rows)


# =============================================================================
# 16. Pass / fail
# =============================================================================
def pass_fail_table(f6, shuffle_p: float, n_episodes: int, loo: pd.DataFrame, cfg: Config):
    """
    Step 16. All four components required.
    Output : (table with one row per component + OVERALL, verdict "PASS"/"FAIL").
    """
    a, t, p = float(f6.params["const"]), float(f6.tvalues["const"]), float(f6.pvalues["const"])
    p1 = p / 2 if t > 0 else 1 - p / 2
    same, nl = int(loo["same_sign"].sum()), len(loo)
    rows = [
        ["(i) timing alpha", f"alpha/month, NW({cfg.nw_lags}) t, p from t(n-k)",
         f"alpha > 0 and t >= {cfg.t_bar:g}",
         f"alpha={a:.6f}, t={t:.3f}, p2={p:.4f}, p1={p1:.4f}", bool(a > 0 and t >= cfg.t_bar)],
        ["(ii) calendar shuffle", f"one-sided p, {cfg.shuffle_draws} draws", f"p <= {cfg.p_bar:g}",
         f"{shuffle_p:.4f}", bool(shuffle_p <= cfg.p_bar)],
        ["(iii) independent episodes", "merged hold runs in window", f">= {cfg.min_episodes}",
         str(n_episodes), bool(n_episodes >= cfg.min_episodes)],
        ["(iv) leave one industry out", "LOO alphas with the full alpha's sign", f"{nl}/{nl}",
         f"{same}/{nl}", bool(same == nl)],
    ]
    n_pass = sum(r[4] for r in rows)
    tab = pd.DataFrame(rows, columns=["component", "statistic", "bar", "value", "pass"])
    tab["pass"] = np.where(tab["pass"].astype(bool), "PASS", "FAIL")
    verdict = "PASS" if n_pass == 4 else "FAIL"
    tab.loc[len(tab)] = ["OVERALL", "all four required", "4/4", f"{n_pass}/4", verdict]
    return tab, verdict


# =============================================================================
# 17. Run
# =============================================================================
def run_backtest(inputs: dict, cfg: Config = Config(), verbose: bool = True) -> dict:
    """
    Step 17. Runs steps 1-16 once; returns and prints the five outputs.
    inputs: ff49 (path/DataFrame), fac (DataFrame), emv_cat, emv_all, gs10, wti, vix
            (FRED CSV paths, or DataFrames in the FRED CSV layout).
    """
    ff49 = load_ff49(inputs["ff49"])
    fac = load_factors(inputs["fac"])
    emv_cat = load_fred_monthly(inputs["emv_cat"], "EMVENRGYENVREG")
    emv_all = load_fred_monthly(inputs["emv_all"], "EMVOVERALLEMV")
    gs10 = load_fred_monthly(inputs["gs10"], "GS10")
    wti = load_fred_monthly(inputs["wti"], "MCOILWTICO")
    vix = load_vix_monthly(inputs["vix"], "VIXCLS")
    index = _month_range(fac.index.min(), min(fac.index.max(), ff49.index.max()))

    # signal and hold: independent of the Brown leg
    sig = attention_signal(emv_cat, emv_all, cfg)
    cross = crossings(sig["extreme"], cfg)
    hold = hold_positions(cross, index, cfg)
    first_sig, first_pos, first_ret, window = test_window(sig, index, cfg)
    pos_months = _position_months(window)
    eps = hold_episodes(hold, pos_months)
    xtab = crossing_table(sig, cross, pos_months, cfg)

    # books, D, regressors
    leg = leg_pipeline(ff49, fac, cfg.brown, hold, index, cfg)
    D, pi, I_prev = timing_difference(leg["tr_T"], leg["tr_AO"], window, cfg)
    bond = bond_excess_return(gs10, fac["RF"])
    F = attribution_factors(fac, bond, wti, vix, emv_all, index).loc[window]
    summary = diagnostics_summary(sig, hold, leg, ff49, window, first_sig, first_pos, first_ret,
                                  eps, xtab, pi, cfg)

    if eps.empty:
        pf = pd.DataFrame([["(iii) independent episodes", "merged hold runs in window",
                            f">= {cfg.min_episodes}", "0", "FAIL"],
                           ["OVERALL", "all four required", "4/4", "-", "FAIL (no position in window)"]],
                          columns=["component", "statistic", "bar", "value", "pass"])
        if verbose:
            _show("1. PASS / FAIL", pf)
            _show("4a. DIAGNOSTICS", summary)
        return {"pass_fail": pf, "verdict": "FAIL", "diagnostics": summary,
                "crossings": xtab, "episodes": eps}

    X = design_matrix(F, I_prev, cfg)
    coef, f6, f12 = attribution_regression(D, X, cfg)
    alpha = float(f6.params["const"])
    alpha_ls = np.linalg.lstsq(X.to_numpy(float), D.to_numpy(float), rcond=None)[0][0]
    if abs(alpha_ls - alpha) > 1e-10:
        raise RuntimeError("shuffle alpha path disagrees with statsmodels")
    n, k = X.shape
    meta = pd.Series({"n": n, "k": k, "n - k": n - k, "R-squared": f6.rsquared,
                      "adj. R-squared": f6.rsquared_adj, "pi": pi,
                      "window start": f"{window[0]:%Y-%m}", "window end": f"{window[-1]:%Y-%m}",
                      "HAC": f"Bartlett, maxlags {cfg.nw_lags} ({cfg.nw_lags_check} check), no small-sample corr."},
                     name="value").to_frame()
    decomp = decompose_mean(D, X, f6)

    shuf, null = calendar_shuffle(alpha, leg, hold, F, window, pi, I_prev, cfg)
    loo = leave_one_out(ff49, fac, hold, F, window, index, alpha, cfg)
    pf, verdict = pass_fail_table(f6, float(shuf.at["p (one-sided)", "value"]), len(eps), loo, cfg)

    if verbose:
        _show("1. PASS / FAIL", pf)
        _show("2a. ATTRIBUTION (D_t on FF5+UMD, BOND, WTI, dVIX, dlogEMV, I_{t-1} x F)", coef)
        _show("2b. REGRESSION FACTS", meta)
        _show("3. DECOMPOSITION OF mean(D)", decomp)
        _show("4a. DIAGNOSTICS", summary)
        _show("4b. CROSSINGS AFFECTING THE WINDOW", xtab)
        _show("4c. EPISODES", eps)
        _show("4d. LEAVE ONE INDUSTRY OUT", loo)
        _show("5. CALENDAR SHUFFLE", shuf)

    return {"pass_fail": pf, "verdict": verdict, "attribution": coef, "attribution_meta": meta,
            "decomposition": decomp, "diagnostics": summary, "crossings": xtab, "episodes": eps,
            "leave_one_out": loo, "shuffle": shuf, "null_alphas": null,
            "D": D, "X": X, "hold": hold, "signal": sig, "fit_nw6": f6, "fit_nw12": f12}


# =============================================================================
# 18. Pre-run checks: synthetic data, look-ahead test, planted effect
# =============================================================================
def make_synthetic_inputs(seed: int = 0, end: str = "2025-12") -> dict:
    """
    Debug data only: random inputs in exactly the raw layouts run_backtest expects.
    No effect is planted (use plant_timing_effect for a positive control).
    """
    rng = np.random.default_rng(seed)
    idx = _month_range("1926-07", end)
    n = len(idx)
    mkt, smb, hml = rng.normal(0.006, 0.045, n), rng.normal(0.002, 0.03, n), rng.normal(0.003, 0.03, n)
    rmw, cma, umd = rng.normal(0.003, 0.02, n), rng.normal(0.003, 0.02, n), rng.normal(0.006, 0.04, n)
    rf = np.full(n, 0.003)
    names = list(BROWN) + [f"Ind{i:02d}" for i in range(1, 45)]
    beta_m, beta_h = rng.uniform(0.6, 1.4, 49), rng.normal(0.0, 0.5, 49)
    ind = rf[:, None] + mkt[:, None] * beta_m + hml[:, None] * beta_h + rng.normal(0.0, 0.04, (n, 49))
    ff49 = pd.DataFrame(ind, columns=names)
    ff49.loc[idx < _me("1970-01"), "Ind44"] = np.nan              # an industry that starts late
    ff49.insert(0, "date", idx.strftime("%Y-%m-%d"))
    fac = pd.DataFrame({"Mkt-RF": mkt, "SMB": smb, "HML": hml, "RMW": rmw, "CMA": cma,
                        "RF": rf, "UMD": umd}, index=idx).loc[idx >= _me("1963-07")]

    def fred(dates, values, sid):
        return pd.DataFrame({"observation_date": dates.to_period("M").to_timestamp().strftime("%Y-%m-%d"),
                             sid: values})

    e_idx = _month_range("1985-01", end)
    ne = len(e_idx)
    x = np.zeros(ne)
    for i in range(1, ne):
        x[i] = 0.85 * x[i - 1] + rng.normal(0.0, 0.35)
    emv_all = 20.0 * np.exp(rng.normal(0.0, 0.3, ne))
    emv_cat = emv_all * np.exp(-3.5 + x)
    emv_cat[rng.random(ne) < np.where(e_idx < _me("1996-01"), 0.05, 0.01)] = 0.0
    g_idx = _month_range("1953-04", end)
    gs10 = np.clip(5.0 + np.cumsum(rng.normal(0.0, 0.2, len(g_idx))), 0.5, 15.0)
    w_idx = _month_range("1986-01", end)
    wti = 20.0 * np.exp(np.cumsum(rng.normal(0.0, 0.08, len(w_idx))))
    days = pd.bdate_range("1990-01-02", _me(end))
    vix = 20.0 * np.exp(rng.normal(0.0, 0.25, len(days)))
    vix[rng.random(len(days)) < 0.03] = np.nan
    return {"ff49": ff49, "fac": fac,
            "emv_cat": fred(e_idx, emv_cat, "EMVENRGYENVREG"),
            "emv_all": fred(e_idx, emv_all, "EMVOVERALLEMV"),
            "gs10": fred(g_idx, gs10, "GS10"),
            "wti": fred(w_idx, wti, "MCOILWTICO"),
            "vix": pd.DataFrame({"observation_date": days.strftime("%Y-%m-%d"), "VIXCLS": vix})}


def check_no_lookahead(inputs: dict, cfg: Config = Config(), n_random: int = 20,
                       cut_from: str | None = None, cut_to: str | None = None,
                       seed: int = 7) -> pd.DataFrame:
    """
    Look-ahead (truncation) test. For each cut month T keep French data dated <= T and EMV
    data dated <= T - pub_lag (what is published by the end of T), rerun signal, hold, hedge,
    sizing and trade, and compare everything set at the end of T (hold, w, m, overlay, b) with
    the full-sample run. Cuts = every crossing month tau and tau+pub_lag (where a lag bug shows)
    plus random months, restricted to [cut_from, cut_to]. Any row with ok == False = look-ahead.
    """
    ff49 = load_ff49(inputs["ff49"])
    fac = load_factors(inputs["fac"])
    emv_cat = load_fred_monthly(inputs["emv_cat"], "EMVENRGYENVREG")
    emv_all = load_fred_monthly(inputs["emv_all"], "EMVOVERALLEMV")
    index = _month_range(fac.index.min(), min(fac.index.max(), ff49.index.max()))

    def state(ff, fc, ec, ea, idx):
        s = attention_signal(ec, ea, cfg)
        hd = hold_positions(crossings(s["extreme"], cfg), idx, cfg)
        lg = leg_pipeline(ff, fc, cfg.brown, hd, idx, cfg)
        cols = ["hold", "w", "m"] + [f"h_{c}" for c in cfg.hedge_factors]
        return lg["tr_T"][cols].astype(float).join(lg["b"].add_prefix("b_"))

    full = state(ff49, fac, emv_cat, emv_all, index)
    sig = attention_signal(emv_cat, emv_all, cfg)
    cr = crossings(sig["extreme"], cfg)
    taus = cr.index[cr.to_numpy(bool)]
    first = _shift_months(sig["extreme"].first_valid_index(), cfg.pub_lag)
    lo = max(first, _me(cut_from)) if cut_from else first
    hi = min(index[-1], _me(cut_to)) if cut_to else index[-1]
    pool = index[(index >= lo) & (index <= hi)]
    rng = np.random.default_rng(seed)
    cuts = set(taus) | {_shift_months(t, cfg.pub_lag) for t in taus}
    cuts |= set(pd.DatetimeIndex(rng.choice(pool.to_numpy(), size=min(n_random, len(pool)), replace=False)))
    cuts = sorted(t for t in cuts if lo <= t <= hi)

    rows = []
    for T in cuts:
        L = _shift_months(T, -cfg.pub_lag)
        part = state(ff49.loc[:T], fac.loc[:T], emv_cat.loc[:L], emv_all.loc[:L], index[index <= T])
        a, b = part.loc[T], full.loc[T]
        diff = (a - b).abs().max()
        rows.append({"cut": T, "max_abs_diff": float(diff) if pd.notna(diff) else 0.0,
                     "nan_mismatch": bool((a.isna() != b.isna()).any())})
    out = pd.DataFrame(rows, columns=["cut", "max_abs_diff", "nan_mismatch"])
    out["ok"] = (out["max_abs_diff"] <= 1e-12) & ~out["nan_mismatch"]
    return out


def plant_timing_effect(inputs: dict, cfg: Config = Config(), delta: float = 0.03) -> dict:
    """
    Positive control (synthetic only). Lowers each Brown industry's return by `delta` in every
    month t entered with a position (hold set at the end of t-1). The signal is untouched, so
    the hold series is unchanged. Expected: alpha > 0, roughly
    delta * (1 - pi) * mean_t(I_{t-1} m_{t-1}), with a small shuffle p.
    """
    emv_cat = load_fred_monthly(inputs["emv_cat"], "EMVENRGYENVREG")
    emv_all = load_fred_monthly(inputs["emv_all"], "EMVOVERALLEMV")
    raw = _read(inputs["ff49"])
    months = _to_month_end(raw["date"])
    sig = attention_signal(emv_cat, emv_all, cfg)
    hold = hold_positions(crossings(sig["extreme"], cfg), _month_range(months.min(), months.max()), cfg)
    on = hold.shift(1, fill_value=False).reindex(months, fill_value=False).to_numpy(bool)
    raw.loc[on, list(cfg.brown)] = raw.loc[on, list(cfg.brown)] - delta
    out = dict(inputs)
    out["ff49"] = raw
    return out


if __name__ == "__main__":
    MODE = "synthetic"      # "synthetic" -> "seen" -> "final" (final: once, no edits after "seen")

    if MODE == "synthetic":
        cfg = replace(Config(), shuffle_draws=500)
        syn = make_synthetic_inputs(seed=0)
        _show("LOOK-AHEAD TEST (synthetic)", check_no_lookahead(syn, cfg))
        out = run_backtest(syn, cfg)
        out_planted = run_backtest(plant_timing_effect(syn, cfg, delta=0.03), cfg)
    else:
        fac = pd.read_pickle("fac.pkl")     # your parsed French FF5 + UMD DataFrame (decimals)
        inputs = dict(ff49="ff49_industry_monthly.csv", fac=fac,
                      emv_cat="EMVENRGYENVREG.csv", emv_all="EMVOVERALLEMV.csv",
                      gs10="GS10.csv", wti="MCOILWTICO.csv", vix="VIXCLS.csv")
        if MODE == "seen":
            cfg = replace(Config(), test_start="2010-01", test_end="2022-12")
            _show("LOOK-AHEAD TEST (seen data only)", check_no_lookahead(inputs, cfg, cut_from="2010-01"))
        elif MODE == "final":
            cfg = Config()
        else:
            raise ValueError(MODE)
        out = run_backtest(inputs, cfg)
```

---

## (a) Assumptions beyond the spec

**The three silent points (each is a parameter)**

1. **Zero month between two extreme months** (`bridge_missing=True`, `bridge_limit=6`). Extreme -> zero -> extreme is one run, so the second extreme month is **not** a new crossing. The last defined flag is carried across up to 6 consecutive missing months (zero or signal-off). Beyond 6, the earlier state counts as "not extreme". Why: the spec says a zero is missing data, not low attention, so it should neither end an episode nor start one. A new crossing would also extend the hold, which would let missing data lengthen positions. The cap of 6 equals the hold length and the largest number of zeros a valid 60-month window can contain. Setting `False` gives the other reading.
2. **Shuffle** (`shuffle_min_gap=1`, `shuffle_recompute_pi_I=True`). The K actual episodes, at the lengths they have inside the window, are placed uniformly at random over all non-overlapping arrangements, in random order. They are placed within the position months that earn inside the window, with at least one flat month between blocks. Months before the window keep their actual state. Each draw rebuilds w, the overlay, costs, pi and I_{t-1}, then refits the full regression and records the alpha. Why: the null statistic has to be computed exactly like the actual one. Keeping pi and I fixed would pair a shuffled D with the actual hold dummies, which is a different statistic. The gap keeps K and the total number of hold months fixed, so draws differ only in timing. p = (1 + #{null >= actual}) / (1 + 5000), one-sided because the bar requires a positive alpha.
3. **Hold still open at 2009-12** (`end_policy="truncate"`). It stops contributing. The position set at the end of 2009-12 earns in 2010-01 and its trade cost is paid in 2010-01, so both fall outside the window. The episode still counts toward (iii). Why: this follows the "trade at end of t pays in t+1" convention and uses no 2010 data, which you have already seen. `"liquidate"` instead charges both books' closing cost in 2009-12.

**Other assumptions**

4. All inputs are decimals, and the loaders stop if the data look like percentages. FF49 values <= -0.99 (French's -99.99 missing code) are set to NaN.
5. The Brown leg averages the Brown industries available in each month. All five should exist in every month the test uses, and the diagnostics print the minimum count.
6. Every series is mapped to month-end on a complete monthly calendar. Missing months stay NaN and nothing is forward-filled.
7. The hedge needs all 60 months present, and sizing needs all 36 residuals present. Both are estimated on the full history from 1963-07, so the results do not depend on where the window starts.
8. The overlay is h_t = -w_t b_t per factor. Its turnover |h_t - h_{t-1}| includes beta and sizing drift during a hold. There is no drift of weights within a month, no financing or borrow fee, and no cost on the cash/RF leg.
9. The timed book starts flat, so its first entry pays full cost. The always-on book has been on since the 1970s, so it pays no entry cost at the start of the window.
10. The z-score uses the mean and the sample standard deviation (ddof=1) of log s over the nonzero months in calendar months t-59..t, including t.
11. The z window must be a full 60 calendar months of EMV history, so the first z is 1989-12. In a full window, "no more than 10% zeros" (at most 6) already implies "at least 48 nonzero", so the 48 condition never binds. Both conditions are coded.
12. A NaN EMV value (none are expected) is treated like a zero.
13. Signal-off months are treated like zero months: they get no flag, cannot trigger, and are not added to the percentile history. A hold that is already running continues through zero or off months, because the hold is a calendar count.
14. The percentile uses numpy's linear interpolation, and "exceeds" is strict (>). The history is every prior valid z from 1989-12 on, and it is not restarted at the window.
15. The first month with a defined flag can itself be a crossing, because the undefined month before it counts as "not extreme".
16. "First month the rule can fire" is the end of (first flagged signal month + pub_lag). The window starts at the first return month that position can affect, or at 1993-01 if that is later. On complete data this should give roughly 1995-02 rather than 1993-01.
17. pi uses the positions that earn inside the window (w_{t-1}), so it is an ex-post scale for the benchmark. That is fine for attribution, but pi is not a tradable quantity.
18. I_{t-1} is the hold state set at the end of t-1.
19. As the spec is written, the regression has **no I_{t-1} main effect**, so alpha is the average timing alpha across all months. If the team intended a main effect, alpha would mean something different.
20. BOND uses exact cash-flow duration and convexity. GS10 is a monthly average but is used as if it were the month-end yield, as in the spec's formula. RF is French's.
21. dVIX is the difference between monthly means of the available daily closes. WTI and EMV changes are month-over-month log differences.
22. HAC uses the Bartlett kernel with no small-sample correction, which is the statsmodels default and your team's convention. Pass (i) uses the NW(6) t, and NW(12) is reported only. Both the two-sided and one-sided t(n-k) p-values are printed.
23. For (iii), an episode is a maximal run of hold months among the position months that earn in the window. A run cut by the window edge counts once.
24. For (iv), the hold timing stays fixed because the signal does not use returns. "Keeps its sign" means the same sign as the full-sample alpha, for all five industries.
25. Current data vintages are used, and French's data are revised slightly each year. As I understand their construction, the EMV category trackers are the overall EMV times a share of articles. So any full-sample scaling cancels in s_t and in the log z-score, and it creates no look-ahead.

## (b) Checks before the single real run (and what makes you stop)

1. **Units and dates.** Print the head and tail of every parsed input. EMV 1985-01 must sit at 1985-01-31, and the first monthly VIX value must be 1990-01. **Stop** if any series is shifted by a month, if the loaders flag percent units, or if any Brown industry is missing in the window.
2. **Reproduce the team trade** on 2010-2022. Compare m_t, b_t, R^AO and your validated timed series against the team's code. **Stop** if any month differs by more than 1e-8.
3. **Look-ahead test.** Run `check_no_lookahead` on synthetic data and on the seen window (cuts >= 2010-01). It cuts at every crossing month, the month after it, and random months. **Stop** if any row has `ok == False`. Also run a lag probe: set `pub_lag=0`. **Stop** if the hold series does not move exactly one month earlier, whether because nothing changes or because more than a shift changes.
4. **Signal unit tests** on hand-built series:
   - extreme -> zero -> extreme gives one crossing with bridging and two without;
   - 6 zeros in the trailing 60 leaves the signal on, and 7 turns it off;
   - a zero month has NaN z and does not lengthen the percentile history.

   **Stop** on any mismatch.
5. **Hold tests.** One crossing at tau must give exactly six position months, tau+1..tau+6. Crossings at tau and tau+3 must give one 9-month run. w must stay in [-1, 0]. **Stop** otherwise.
6. **Cost test.** With all costs set to 0, net must equal gross exactly. For one hand-picked entry, the cost must equal 10bp*m + 5bp*|h_Mkt| + 25bp*(|h_SMB| + |h_HML|) and be paid the next month. **Stop** on a mismatch.
7. **BOND.** At y = 5%, the modified duration should be about 7.79, from the analytic value (1/y)(1 - 1.025^-20). The correlation of BOND with -Deltay should be above 0.99. **Stop** otherwise.
8. **Regression mechanics.** Check that k = 15, n - k = df_resid, the identity gap is below 1e-12, the design matrix has full rank, and the lstsq alpha equals the statsmodels alpha (the code raises if not). **Stop** on any failure.
9. **Shuffle validity** on synthetic data:
   - every draw has K episodes and the same total of hold months;
   - the same seed gives the same p on two runs;
   - across 20 synthetic seeds with no planted effect, p <= 0.05 occurs about 5% of the time.

   **Stop** if it happens in more than 4 of the 20.
10. **Positive control.** Run `plant_timing_effect(delta=0.03)`. Alpha should be positive with t well above 2, and the shuffle p should be small. With `delta=0`, you should get the original alpha back exactly. **Stop** if the planted effect is missed or has the wrong sign.
11. **Seen-window run** on 2010-2022, end to end. Numbers should be consistent with the team's existing results. Any change after this point must be a bug fix, and a bug fix means rerunning checks 1-10.
12. **Coverage without peeking.** Before the real run, check only that the 1993-2009 regressors and industries have no NaNs. Do not compute anything from 1993-2009 returns. **Stop** if data are missing, and fix the data, not the rule.
13. **Freeze.** Record the SHA-256 of the script and of every input file, the printed `Config`, the seed and the package versions, then commit. **Stop** if anything differs at run time. Then set `MODE = "final"` and run once.
\end{lstlisting}

{\raggedright\paragraph{Follow-up prompt (verbatim)} \texttt{followup\_\allowbreak{}prompt.md}, 444 words (437 excluding Markdown symbols).\par}

\begin{lstlisting}[style=transcript]
I ran your script on the real data and through your checks. It ran the first time, with only file paths and `fac` supplied, and reaches the same verdict as my own implementation (FAIL, 1 of 4). The numbers differ, and one of your stop rules fires.

1. `make_synthetic_inputs`. GS10 is a random walk clipped at 0.5 that sits on the floor (76% of 1995-2009 months for seed 0, every month for seed 16). BOND is then constant, I x BOND acts as an I main effect, and your positive control fails check (b)10: 4% of the planted 3% effect reaches alpha (t 0.98), while I x BOND has t -10.8. Seed 16 crashes the null-seed loop (rank-deficient design).

2. `attention_signal`. Requiring 60 full calendar months before the first z is a silent reading you hard-coded. I treat pre-1985 months as unavailable, not zero, so the 48-nonzero minimum is the burn-in. Your reading moves the first threshold from 1994-01 to 1994-12, adds a 2002-12 crossing and starts the window at 1995-02, not 1994-03: pi 0.809 vs. 0.729, shuffle p 0.536 vs. 0.478.

3. `leg_pipeline`. You hedge with the Mkt-RF, SMB and HML in `fac`, the FF5 file. The team hedge uses its FF3 file, which my prompt never said (my gap). SMB differs by up to 3.5 pp a month, b_t moves by up to 0.09, and your check (b)2 stops at 2.3e-3 a month (6e-17 with FF3).

4. `pass_fail_table`. (iv) passes a negative alpha whose five drop-one alphas are also negative.

With readings 2 and 3 switched to mine, your code matches my implementation to 1e-13, all 5,000 shuffle draws included.

Two smaller gaps. `check_no_lookahead` compares end-of-T states only, so it missed both return-side bugs I injected (month-t return hedged with b_t; a position earning its own month). Seen mode ends at 2022-12; its five holdout months flip the seen alpha from -0.16%/yr (t -0.24) to +0.26%/yr (t 0.37).

Please send corrected versions of these functions only, each with a unit test:
- `make_synthetic_inputs`: BOND varies in every window for seeds 0-20, and the positive control puts at least 80% of the planted shift in alpha with t > 2.
- `attention_signal`: a `z_burn_in` parameter with both readings; on a zero-free series from 1985-01, the first z falls in month 48 under mine and month 60 under yours.
- `leg_pipeline`: a separate `hedge_fac` input; changing its SMB moves b_t but not the attribution regressors.
- `pass_fail_table`: a negative alpha fails (iv).
- `check_no_lookahead`: a return-side comparison that flags both injected bugs.
- `__main__`: seen mode ends at 2022-07, with `cut_to` on its look-ahead call.

These fixes are for reproducibility; the frozen rule stays.
\end{lstlisting}

{\raggedright\paragraph{Follow-up reply (verbatim; Claude)} \texttt{claude\_\allowbreak{}followup\_\allowbreak{}response.md}, 4,745 words (4,066 excluding Markdown symbols).\par}

\begin{lstlisting}[style=transcriptsmall]
The replacements below cover the six functions you listed, and each one comes with a unit test. They share a few small helpers that handle the raw input files, and one adapter, `run_panel`, which is the only place the tests call the rest of the script. I can't run code, so none of this has been tested yet. Please run the tests before anything else, and if one fails, send me the assertion message.

**Changes needed in the driver**
1. `run_backtest` gets two new arguments, `hedge_fac` and `z_burn_in`. Pass `hedge_fac` to every `leg_pipeline` call: the main run and all five leave-one-out rebuilds. Pass `z_burn_in` to `attention_signal`.
2. The overlay return `-w_t b_t'f_{t+1}` and the overlay costs must use the leg's `f_*` and `b_*` columns (FF3), not `fac`. If the old overlay used `fac`, the FF5 SMB was getting in there as well as in the hedge. The attribution still uses `fac`.
3. `pass_fail_table` now takes plain numbers.
4. `run_panel` assumes this driver signature: `run_backtest(paths, fac, start, end, *, hedge_fac, z_burn_in, n_shuffle, run_loo, verbose)`. It also assumes the output keys listed in its docstring. If your names differ, change them in `run_panel` only.

```python
# =============================================================================
# Replacements. Everything not defined here is unchanged from the earlier script.
# =============================================================================
import os
import sys
import tempfile

import numpy as np
import pandas as pd

# Constants (already in the script; repeated so this block stands alone)
BROWN = ["Util", "Ships", "Aero", "Steel", "BldMt"]
HEDGE_COLS = ["Mkt-RF", "SMB", "HML"]
FRED_IDS = ["EMVENRGYENVREG", "EMVOVERALLEMV", "GS10", "MCOILWTICO", "VIXCLS"]
FF49_NAMES = [
    "Agric", "Food", "Soda", "Beer", "Smoke", "Toys", "Fun", "Books", "Hshld", "Clths",
    "Hlth", "MedEq", "Drugs", "Chems", "Rubbr", "Txtls", "BldMt", "Cnstr", "Steel", "FabPr",
    "Mach", "ElcEq", "Autos", "Aero", "Ships", "Guns", "Gold", "Mines", "Coal", "Oil",
    "Util", "Telcm", "PerSv", "BusSv", "Hardw", "Softw", "Chips", "LabEq", "Paper", "Boxes",
    "Trans", "Whlsl", "Rtail", "Meals", "Banks", "Insur", "RlEst", "Fin", "Other",
]
LATE_START = {"Soda": "1963-07-31", "Hlth": "1969-07-31", "Rubbr": "1963-07-31", "FabPr": "1963-07-31",
              "Guns": "1963-07-31", "Gold": "1963-07-31", "Softw": "1965-07-31"}
STATE_COLS = ["w", "w_ao", "m", "hold"] + [f"b_{k}" for k in HEDGE_COLS]
POS_OF = {"RT": "w", "RAO": "w_ao"}          # return column -> the position that earns it
TEST_WINDOW = ("1993-01-31", "2009-12-31")
SEEN_WINDOW = ("2010-01-31", "2022-07-31")   # the holdout starts 2022-08
MONTH_END = pd.offsets.MonthEnd()


# -----------------------------------------------------------------------------
# Raw-input helpers. Checks and tests work on `raw` (the files as read from disk),
# so synthetic, perturbed and truncated data all pass through your unchanged loader.
# -----------------------------------------------------------------------------
def read_raw(paths, fac, hedge_fac):
    """Inputs: dict of CSV paths (keys 'ff49' and the FRED ids), fac (FF5+UMD), hedge_fac (team FF3).
    Output: dict of raw frames exactly as on disk plus the two factor frames. Timing: none."""
    raw = {"ff49": pd.read_csv(paths["ff49"], parse_dates=["date"], float_precision="round_trip"),
           "fac": fac, "hedge_fac": hedge_fac}
    for sid in FRED_IDS:
        df = pd.read_csv(paths[sid], parse_dates=["observation_date"], float_precision="round_trip")
        df[sid] = pd.to_numeric(df[sid], errors="coerce")        # FRED blanks -> NaN
        raw[sid] = df
    return raw


def write_raw_csvs(raw, folder):
    """Inputs: raw dict, folder. Output: dict of paths in the same layout as the real files. Timing: none."""
    paths = {"ff49": os.path.join(folder, "ff49_industry_monthly.csv")}
    raw["ff49"].to_csv(paths["ff49"], index=False, date_format="%Y-%m-%d")
    for sid in FRED_IDS:
        paths[sid] = os.path.join(folder, f"{sid}.csv")
        raw[sid].to_csv(paths[sid], index=False, date_format="%Y-%m-%d")
    return paths


def truncate_raw(raw, cut_to):
    """Inputs: raw dict, month-end cut_to. Output: copy with every row dated after cut_to removed
    (FRED monthly rows are dated the 1st, so month cut_to itself is kept). Timing: nothing after cut_to survives."""
    cut = pd.Timestamp(cut_to)
    out = {}
    for k, v in raw.items():
        if not isinstance(v, pd.DataFrame):
            out[k] = v
        elif "date" in v.columns:
            out[k] = v[v["date"] <= cut].copy()
        elif "observation_date" in v.columns:
            out[k] = v[v["observation_date"] <= cut].copy()
        else:
            out[k] = v.loc[v.index <= cut].copy()
    return out


def _last_date(raw):
    dates = []
    for v in raw.values():
        if isinstance(v, pd.DataFrame) and len(v):
            col = "date" if "date" in v.columns else ("observation_date" if "observation_date" in v.columns else None)
            dates.append(v[col].max() if col else v.index.max())
    return max(dates)


def _bump(raw, key, T, cols, fn):
    """Copy of raw with fn applied to `cols` of frame `key` in the calendar month of T."""
    r = dict(raw)
    df = raw[key].copy()
    p = pd.Timestamp(T).to_period("M")
    if "date" in df.columns:
        rows = (df["date"].dt.to_period("M") == p).to_numpy()
    elif "observation_date" in df.columns:
        rows = (df["observation_date"].dt.to_period("M") == p).to_numpy()
    else:
        rows = df.index.to_period("M") == p
    df.loc[rows, cols] = fn(df.loc[rows, cols])
    r[key] = df
    return r


def _maxdiff(a, b, cols):
    """Max |a - b| over cols on a's rows; NaN vs NaN counts as equal, NaN vs number as inf."""
    av = a[cols].astype(float).to_numpy()
    bv = b.reindex(a.index)[cols].astype(float).to_numpy()
    if (np.isnan(av) ^ np.isnan(bv)).any():
        return np.inf
    d = np.abs(av - bv)[~(np.isnan(av) & np.isnan(bv))]
    return float(d.max()) if d.size else 0.0


def _ret_err(base, pert, T, expected):
    errs = []
    for col, e in expected.items():
        d = pert.at[T, col] - base.at[T, col] - e if T in pert.index else np.nan
        errs.append(abs(d) if np.isfinite(d) else np.inf)
    return max(errs)


def run_panel(raw, start, end, z_burn_in="nonzero", n_shuffle=0, loo=False, verbose=False):
    """
    ADAPTER to the unchanged driver: the only place the tests touch it. Edit names here if yours differ.
    Inputs : raw dict; window; burn-in reading; shuffle draws (0 = skip); leave-one-out on/off.
    Output : (panel, out).
             panel, indexed by month-end t: w, w_ao (positions set at end of t), RT, RAO (net returns earned
             in t), m, hold, b_Mkt-RF, b_SMB, b_HML (end-of-t states), I (1 if a timed position was held
             entering t).
             out: D (Series), X (design DataFrame incl. constant), alpha, alpha_t (NW(6)), plus the
             driver's printed tables.
    Timing : the window end is clipped to the last month in raw, so truncated inputs run cleanly.
    """
    end = min(pd.Timestamp(end), pd.Timestamp(raw["ff49"]["date"].max()))
    with tempfile.TemporaryDirectory() as d:
        paths = write_raw_csvs(raw, d)
        out = run_backtest(paths, raw["fac"], start, end, hedge_fac=raw["hedge_fac"],   # noqa: F821
                           z_burn_in=z_burn_in, n_shuffle=n_shuffle, run_loo=loo, verbose=verbose)
    return out["panel"], out


# =============================================================================
# 1. make_synthetic_inputs
# =============================================================================
def make_synthetic_inputs(seed=0, out_dir=None, plant=0.0, plant_months=None, end="2026-06-30"):
    """
    Synthetic stand-ins for every input, shaped exactly like the real files.

    Inputs
      seed         : fixed RNG seed (nothing is tuned on it).
      out_dir      : if given, also write the six CSVs there; their paths are returned under 'paths'.
      plant        : positive-control effect (decimal). Every Brown industry return is lowered by `plant`
                     in each month of `plant_months`, so the short Brown leg gains there.
      plant_months : month-ends t whose return is earned by a timed position (I_{t-1} = 1) in an unplanted
                     run with the same seed. The signal files do not depend on returns and the plant draws
                     no random numbers, so the hold schedule and all other data are identical with and without it.
    Outputs
      dict: 'ff49' (date column + 49 industries, NaN before an industry starts), 'fac' (Mkt-RF, SMB, HML,
      RMW, CMA, RF, UMD from 1963-07), 'hedge_fac' (FF3: Mkt-RF, SMB, HML, RF from 1926-07; its SMB differs
      from fac's, as in the real files), one frame per FRED id (observation_date + id).
    Timing
      Returns at month-ends. FRED monthly rows dated the 1st. GS10 from 1953-04, EMV from 1985-01,
      WTI from 1986-01, VIXCLS daily business days from 1990-01-02 with ~3% blanks.
    Fix
      GS10 used to be a random walk clipped at 0.5, so it could sit on the clip and make BOND constant.
      log(GS10) is now a stationary AR(1): always positive, never clipped, and it moves every month.
    """
    rng = np.random.default_rng(seed)
    idx = pd.date_range("1926-07-31", end, freq=MONTH_END)
    n = len(idx)

    def ar1(nobs, mu, phi, sig, x0):
        eps = rng.standard_normal(nobs)
        x = np.empty(nobs)
        x[0] = x0
        for i in range(1, nobs):
            x[i] = mu + phi * (x[i - 1] - mu) + sig * eps[i]
        return x

    y10 = np.exp(ar1(n, np.log(5.0), 0.985, 0.045, np.log(3.5)))            # percent, roughly 2-11
    rf = y10 * 0.6 * np.exp(ar1(n, 0.0, 0.9, 0.05, 0.0)) / 1200.0              # monthly decimal

    zf = rng.standard_normal((n, 7))
    mkt, smb3, hml = 0.006 + 0.045 * zf[:, 0], 0.002 + 0.030 * zf[:, 1], 0.003 + 0.030 * zf[:, 2]
    rmw, cma, umd = 0.003 + 0.020 * zf[:, 3], 0.003 + 0.020 * zf[:, 4], 0.006 + 0.040 * zf[:, 5]
    smb5 = smb3 + 0.004 * zf[:, 6]
    hedge_fac = pd.DataFrame({"Mkt-RF": mkt, "SMB": smb3, "HML": hml, "RF": rf}, index=idx)
    fac = pd.DataFrame({"Mkt-RF": mkt, "SMB": smb5, "HML": hml, "RMW": rmw, "CMA": cma,
                        "RF": rf, "UMD": umd}, index=idx).loc["1963-07-31":]

    brown_common = 0.015 * rng.standard_normal(n)
    rets = {}
    for name in FF49_NAMES:
        bm, bs, bh = rng.normal(1.0, 0.2), rng.normal(0.2, 0.3), rng.normal(0.2, 0.3)
        r = rf + bm * mkt + bs * smb3 + bh * hml + 0.05 * rng.standard_normal(n)
        if name in BROWN:
            r = r + brown_common
        if name in LATE_START:
            r = np.where(idx < pd.Timestamp(LATE_START[name]), np.nan, r)
        rets[name] = r
    if plant != 0.0:
        if plant_months is None:
            raise ValueError("plant != 0 needs plant_months")
        hit = idx.to_period("M").isin(pd.DatetimeIndex(plant_months).to_period("M"))
        for name in BROWN:
            rets[name] = np.where(hit, rets[name] - plant, rets[name])
    ff49 = pd.DataFrame(rets, index=idx)[FF49_NAMES]
    ff49.insert(0, "date", idx)
    ff49 = ff49.reset_index(drop=True)

    def fred(mask, values, sid):
        return pd.DataFrame({"observation_date": idx[mask].to_period("M").to_timestamp(), sid: values})

    g = idx >= pd.Timestamp("1953-04-30")
    e = idx >= pd.Timestamp("1985-01-31")
    o = idx >= pd.Timestamp("1986-01-31")
    ne = int(e.sum())
    overall = np.exp(ar1(ne, np.log(25.0), 0.8, 0.15, np.log(25.0)))            # never zero
    env = overall * np.exp(ar1(ne, np.log(0.02), 0.6, 0.35, np.log(0.02)))
    p_zero = np.where(idx[e] < pd.Timestamp("1995-01-01"), 0.05, 0.01)          # zeros cluster early
    env = np.where(rng.random(ne) < p_zero, 0.0, env)
    wti = np.exp(ar1(int(o.sum()), np.log(40.0), 0.98, 0.09, np.log(22.0)))
    days = pd.bdate_range("1990-01-02", end)
    vix = np.round(np.exp(ar1(len(days), np.log(19.0), 0.98, 0.06, np.log(19.0))), 2)
    vix[rng.random(len(days)) < 0.03] = np.nan

    raw = {"ff49": ff49, "fac": fac, "hedge_fac": hedge_fac,
           "GS10": fred(g, y10[g], "GS10"),
           "EMVOVERALLEMV": fred(e, overall, "EMVOVERALLEMV"),
           "EMVENRGYENVREG": fred(e, env, "EMVENRGYENVREG"),
           "MCOILWTICO": fred(o, wti, "MCOILWTICO"),
           "VIXCLS": pd.DataFrame({"observation_date": days, "VIXCLS": vix})}
    if out_dir is not None:
        os.makedirs(out_dir, exist_ok=True)
        raw["paths"] = write_raw_csvs(raw, out_dir)
    return raw


def positive_control(seed=0, plant=0.03, window=TEST_WINDOW):
    """
    Check (b)10. Inputs: seed, plant (Brown leg lower by `plant` in every held month), window.
    Output: dict with the planted shift = mean(D_planted) - mean(D_unplanted), the share of it in alpha,
    and the planted run's alpha and NW(6) t. Timing: the plant sits in month t when I_{t-1} = 1.
    """
    raw0 = make_synthetic_inputs(seed)
    p0, o0 = run_panel(raw0, *window)
    w0 = p0.loc[window[0]:window[1]]
    held = w0.index[(w0["I"] == 1).to_numpy()]
    raw1 = make_synthetic_inputs(seed, plant=plant, plant_months=held)
    p1, o1 = run_panel(raw1, *window)
    shift = o1["D"].mean() - o0["D"].mean()
    return dict(seed=seed, same_I=bool(p1["I"].equals(p0["I"])), n_held=len(held), shift=shift,
                share=(o1["alpha"] - o0["alpha"]) / shift, alpha=o1["alpha"], t=o1["alpha_t"])


def test_make_synthetic_inputs(seeds=range(21), pc_seeds=(0, 1, 2)):
    for seed in seeds:
        raw = make_synthetic_inputs(seed)
        g = raw["GS10"]["GS10"]
        assert g.notna().all() and g.min() > 0.5, (seed, g.min())               # no floor
        assert (g.diff().iloc[1:] != 0).all(), seed                              # moves every month
        assert g.diff().rolling(24).std().dropna().min() > 0.02, seed            # and in every window
        for start, end in (TEST_WINDOW, SEEN_WINDOW):
            panel, out = run_panel(raw, start, end)
            X = out["X"]
            I = panel["I"].reindex(X.index).astype(float)
            assert np.linalg.matrix_rank(X.to_numpy()) == X.shape[1], (seed, start)   # seed 16 crashed here
            assert (X.std() < 1e-12).sum() == 1, (seed, start)                         # only the constant
            coef = np.linalg.lstsq(X.to_numpy(), I.to_numpy(), rcond=None)[0]
            r2 = 1 - ((I - X @ coef) ** 2).sum() / ((I - I.mean()) ** 2).sum()
            assert r2 < 0.5, (seed, start, r2)          # no regressor combination can stand in for I itself
    for seed in pc_seeds:
        r = positive_control(seed)
        assert r["same_I"] and r["n_held"] > 0 and r["shift"] > 0, r
        assert r["share"] >= 0.8 and r["t"] > 2, r


# =============================================================================
# 2. attention_signal
# =============================================================================
def attention_signal(env, overall, z_burn_in="nonzero", window=60, min_nonzero=48,
                     max_zero_frac=0.10, pct=0.80, min_prior_z=60):
    """
    Rule steps 1-3: ratio, zero handling, rolling z, past-only percentile threshold, extreme flag.

    Inputs
      env, overall : Series of EMVENRGYENVREG and EMVOVERALLEMV indexed by data month (month-end).
                     Months before the first observation do not exist: unavailable, not zero.
      z_burn_in    : "nonzero"  (default, your reading): z_t exists once the trailing 60 months hold
                                 >= 48 nonzero months, counting only months inside the sample.
                     "calendar" (my previous hard-coded reading): also require 60 calendar months since
                                 the first observation.
      Rules shared by both readings: a zero month has no z and cannot be extreme; the month is "off"
      (no z) when the trailing 60 calendar months contain more than 6 zeros (10% of 60) or fewer than
      48 nonzero months; the z window and the percentile history use only months with a z.
    Outputs
      DataFrame by data month t: s, zero, n_nonzero, n_zero, on, z, n_prior_z, thr, extreme.
      z_t = (log s_t - mean) / std(ddof=1) over nonzero months t-59..t.
      thr_t = 80th percentile (linear) of all z before t, needing >= 60 of them. extreme_t = on_t and z_t > thr_t.
    Timing
      Row t uses EMV data for months <= t only. Publication lag: row t is usable at the end of t+1.
      The lag is applied in the hold step (unchanged), which reads extreme[t-1] at the end of t.
    """
    if z_burn_in not in ("nonzero", "calendar"):
        raise ValueError("z_burn_in must be 'nonzero' or 'calendar'")
    env, overall = env.copy(), overall.copy()
    env.index = pd.DatetimeIndex(env.index) + pd.offsets.MonthEnd(0)
    overall.index = pd.DatetimeIndex(overall.index) + pd.offsets.MonthEnd(0)
    df = pd.concat({"env": env, "overall": overall}, axis=1).sort_index()
    idx = pd.date_range(df.dropna().index.min(), df.index.max(), freq=MONTH_END)
    df = df.reindex(idx)

    avail = df["env"].notna() & df["overall"].notna() & (df["overall"] > 0)
    zero = avail & (df["env"] == 0)
    good = avail & (df["env"] > 0)
    s = (df["env"] / df["overall"]).where(good)
    x = np.log(s)

    n_nonzero = good.astype(int).rolling(window, min_periods=1).sum()
    n_zero = zero.astype(int).rolling(window, min_periods=1).sum()
    max_zero = int(np.floor(max_zero_frac * window + 1e-9))
    on = good & (n_nonzero >= min_nonzero) & (n_zero <= max_zero)
    if z_burn_in == "calendar":
        on &= pd.Series(np.arange(len(idx)) >= window - 1, index=idx)

    roll = x.rolling(window, min_periods=min_nonzero)
    z = ((x - roll.mean()) / roll.std(ddof=1)).where(on)

    zv = z.to_numpy()
    thr = np.full(len(zv), np.nan)
    n_prior = np.zeros(len(zv), dtype=int)
    hist = []
    for i, zi in enumerate(zv):
        n_prior[i] = len(hist)
        if len(hist) >= min_prior_z:
            thr[i] = np.quantile(hist, pct)
        if np.isfinite(zi):
            hist.append(zi)
    thr = pd.Series(thr, index=idx)
    extreme = on & (z > thr)
    return pd.DataFrame({"s": s, "zero": zero, "n_nonzero": n_nonzero, "n_zero": n_zero, "on": on,
                         "z": z, "n_prior_z": n_prior, "thr": thr, "extreme": extreme})


def test_attention_signal():
    idx = pd.date_range("1985-01-31", "2005-12-31", freq=MONTH_END)
    rng = np.random.default_rng(1)
    overall = pd.Series(100 * np.exp(rng.normal(0, 0.2, len(idx))), idx)
    env = overall * np.exp(rng.normal(np.log(0.01), 0.3, len(idx)))
    a = attention_signal(env, overall, z_burn_in="nonzero")
    b = attention_signal(env, overall, z_burn_in="calendar")
    assert a["z"].first_valid_index() == idx[47]           # month 48
    assert b["z"].first_valid_index() == idx[59]           # month 60
    assert a["thr"].first_valid_index() == idx[107] and b["thr"].first_valid_index() == idx[119]
    assert np.allclose(a["z"].iloc[59:], b["z"].iloc[59:], rtol=0, atol=0)   # same z once the window is full
    t = 100                                                                    # hand-checked z
    x = np.log(env / overall).iloc[t - 59:t + 1]
    assert abs(a["z"].iloc[t] - (x.iloc[-1] - x.mean()) / x.std(ddof=1)) < 1e-10
    h = a["z"].iloc[:150].dropna().to_numpy()                                  # past-only threshold
    assert a["thr"].iloc[150] == np.quantile(h, 0.8) and a["n_prior_z"].iloc[150] == len(h)
    env4 = env.copy(); env4.iloc[150] *= 10
    a4 = attention_signal(env4, overall)
    pd.testing.assert_series_equal(a["thr"].iloc[:151], a4["thr"].iloc[:151])
    env2 = env.copy(); env2.iloc[[100, 101]] = 0.0                              # zeros dropped, not low
    c = attention_signal(env2, overall)
    assert c["z"].iloc[100:102].isna().all() and not c["extreme"].iloc[100:102].any()
    x2 = np.log((env2 / overall).where(env2 > 0)).iloc[102 - 59:103]
    assert abs(c["z"].iloc[102] - (x2.iloc[-1] - x2.mean()) / x2.std(ddof=1)) < 1e-10
    env3 = env.copy(); env3.iloc[150:157] = 0.0                                 # 7 zeros in 60 -> off
    d = attention_signal(env3, overall)
    assert not d["on"].iloc[150:210].any() and d["on"].iloc[210]
    assert (a["extreme"] == (a["z"] > a["thr"])).all()


# =============================================================================
# 3. leg_pipeline
# =============================================================================
def leg_pipeline(ff49, hedge_fac, brown=BROWN, beta_window=60, sd_window=36, vol_target=0.05):
    """
    Brown leg, rolling hedge, out-of-sample residual, sizing.

    Inputs
      ff49      : industry returns (decimal), month-end DatetimeIndex or a 'date' column.
      hedge_fac : the team's FF3 frame (Mkt-RF, SMB, HML, RF; decimal; month-end). It supplies R^B's RF,
                  the hedge regressors and the overlay factor returns. The attribution's `fac` (FF5+UMD)
                  never enters this function. An FF5 frame is refused.
      brown     : industries in the leg (pass four names for a leave-one-out rebuild).
    Outputs (DataFrame by month-end t)
      RB = mean(brown_t) - RF_t; f_<k> = hedge factor returns in t; a, b_<k> = OLS of RB on [1, f] over
      t-59..t; e = RB_t - a_{t-1} - b_{t-1}'f_t; sd = std(e_{t-35..t}, ddof=1);
      m = min(1, vol_target / (sqrt(12) sd_t)).
    Timing
      Row t uses data through t only; nothing is shifted forward. A position set at the end of t
      (size m_t, hedge b_t) earns in t+1: w_t RB_{t+1} - w_t b_t' f_{t+1}, with f_{t+1} from THIS frame.
    """
    if "date" in ff49.columns:
        ff49 = ff49.set_index("date")
    ff49 = ff49.copy()
    ff49.index = pd.DatetimeIndex(ff49.index)
    missing = [c for c in HEDGE_COLS + ["RF"] if c not in hedge_fac.columns]
    if missing:
        raise ValueError(f"hedge_fac lacks {missing}")
    if {"RMW", "CMA"} & set(hedge_fac.columns):
        raise ValueError("hedge_fac has RMW/CMA: that is the FF5 file; the team hedge uses its FF3 file")
    idx = ff49.index.intersection(hedge_fac.index).sort_values()
    per = idx.to_period("M").asi8
    if len(per) > 1 and (np.diff(per) != 1).any():
        raise ValueError("gaps in the monthly overlap of ff49 and hedge_fac")

    F = hedge_fac.loc[idx, HEDGE_COLS].astype(float)
    rb = ff49.loc[idx, list(brown)].mean(axis=1, skipna=False) - hedge_fac.loc[idx, "RF"]
    y = rb.to_numpy(float)
    X = np.column_stack([np.ones(len(idx)), F.to_numpy()])
    usable = np.isfinite(y) & np.isfinite(X).all(axis=1)
    coef = np.full((len(idx), 1 + len(HEDGE_COLS)), np.nan)
    for i in range(beta_window - 1, len(idx)):
        lo = i - beta_window + 1
        if usable[lo:i + 1].all():
            coef[i] = np.linalg.lstsq(X[lo:i + 1], y[lo:i + 1], rcond=None)[0]

    out = pd.DataFrame(index=idx)
    out["RB"] = rb
    for k in HEDGE_COLS:
        out[f"f_{k}"] = F[k]
    out["a"] = coef[:, 0]
    for j, k in enumerate(HEDGE_COLS):
        out[f"b_{k}"] = coef[:, j + 1]
    fit_prev = out["a"].shift(1) + sum(out[f"b_{k}"].shift(1) * out[f"f_{k}"] for k in HEDGE_COLS)
    out["e"] = out["RB"] - fit_prev
    out["sd"] = out["e"].rolling(sd_window, min_periods=sd_window).std(ddof=1)
    out["m"] = np.minimum(1.0, vol_target / (np.sqrt(12.0) * out["sd"]))
    return out


def test_leg_pipeline():
    raw = make_synthetic_inputs(seed=3)
    ff49, hf = raw["ff49"].set_index("date"), raw["hedge_fac"]
    leg = leg_pipeline(ff49, hf)
    T = pd.Timestamp("2000-06-30"); i = leg.index.get_loc(T); Tm1 = leg.index[i - 1]
    win = leg.index[i - 59:i + 1]
    Xw = np.column_stack([np.ones(60), hf.loc[win, HEDGE_COLS].to_numpy()])
    yw = (ff49.loc[win, BROWN].mean(axis=1) - hf.loc[win, "RF"]).to_numpy()
    bh = np.linalg.lstsq(Xw, yw, rcond=None)[0]
    assert np.allclose(leg.loc[T, ["a"] + [f"b_{k}" for k in HEDGE_COLS]].to_numpy(float), bh, rtol=0, atol=1e-12)
    e_hand = leg.at[T, "RB"] - leg.at[Tm1, "a"] - sum(leg.at[Tm1, f"b_{k}"] * hf.at[T, k] for k in HEDGE_COLS)
    assert abs(leg.at[T, "e"] - e_hand) < 1e-12                                 # residual uses t-1 coefficients
    ff49b, hfb = ff49.copy(), hf.copy()                                          # month T+1 cannot reach row T
    ff49b.loc[leg.index[i + 1], BROWN] += 0.2
    hfb.loc[leg.index[i + 1], HEDGE_COLS] += 0.2
    pd.testing.assert_frame_equal(leg.iloc[:i + 1], leg_pipeline(ff49b, hfb).iloc[:i + 1])
    hf2 = hf.copy()                                                              # hedge SMB moves b, not X
    hf2["SMB"] = hf2["SMB"] + 0.01 * np.random.default_rng(0).standard_normal(len(hf2))
    assert (leg_pipeline(ff49, hf2)["b_SMB"] - leg["b_SMB"]).abs().max() > 1e-3
    _, o1 = run_panel(raw, *TEST_WINDOW)
    _, o2 = run_panel(dict(raw, hedge_fac=hf2), *TEST_WINDOW)
    pd.testing.assert_frame_equal(o1["X"], o2["X"])
    try:
        leg_pipeline(ff49, raw["fac"])
    except ValueError:
        pass
    else:
        raise AssertionError("FF5 frame accepted as hedge_fac")


# =============================================================================
# 4. pass_fail_table
# =============================================================================
def _fin(v):
    try:
        return v is not None and bool(np.isfinite(float(v)))
    except (TypeError, ValueError):
        return False


def _fmt(v, spec=".6f"):
    return format(float(v), spec) if _fin(v) else "nan"


def pass_fail_table(alpha, t_alpha, p_alpha, shuffle_p, n_episodes, loo_alpha,
                    t_bar=2.0, p_bar=0.05, min_episodes=8, n_brown=5):
    """
    Inputs : alpha, t_alpha, p_alpha (NW(6) fit; p from t(n-k)), shuffle_p, n_episodes (merged holds),
             loo_alpha (Series of drop-one alphas indexed by the dropped industry).
    Output : DataFrame indexed (i), (ii), (iii), (iv), overall; columns name, statistic, bar, value,
             passed, result, detail.
    Timing : ex post summary of the single run.
    (iv)   : "keeps its sign" is read as "stays positive": the full alpha and all five drop-one alphas
             must be > 0. A negative alpha whose drop-one alphas are also negative fails. So does a
             missing or NaN drop-one alpha.
    """
    loo = pd.Series(loo_alpha, dtype=float)
    all_a = pd.concat([pd.Series({"full": alpha}, dtype=float), loo])
    ok1 = _fin(alpha) and _fin(t_alpha) and alpha > 0 and t_alpha >= t_bar
    ok2 = _fin(shuffle_p) and shuffle_p <= p_bar
    ok3 = _fin(n_episodes) and n_episodes >= min_episodes
    ok4 = len(loo) == n_brown and bool(all_a.notna().all()) and bool((all_a > 0).all())
    rows = [
        ("(i)", "timing alpha", "NW(6) t of alpha", f"alpha > 0 and t >= {t_bar:g}", t_alpha, ok1,
         f"alpha = {_fmt(alpha)}/month, p = {_fmt(p_alpha, '.4f')} (t, n-k df)"),
        ("(ii)", "calendar shuffle", "shuffle p", f"<= {p_bar:g}", shuffle_p, ok2, ""),
        ("(iii)", "independent episodes", "merged hold episodes", f">= {min_episodes}", n_episodes, ok3, ""),
        ("(iv)", "drop each Brown industry", "min(full, drop-one alphas)", f"> 0, all {n_brown + 1}",
         all_a.min() if len(all_a) else np.nan, ok4, ", ".join(f"{k}: {_fmt(v)}" for k, v in all_a.items())),
    ]
    n_pass = int(sum(r[5] for r in rows))
    rows.append(("overall", "verdict", "components passed", "4 of 4", n_pass, n_pass == 4, f"{n_pass} of 4"))
    t = pd.DataFrame(rows, columns=["component", "name", "statistic", "bar", "value", "passed", "detail"])
    t = t.set_index("component")
    t["result"] = np.where(t["passed"], "PASS", "FAIL")
    return t[["name", "statistic", "bar", "value", "passed", "result", "detail"]]


def test_pass_fail_table():
    pos, neg = pd.Series(0.002, index=BROWN), pd.Series(-0.001, index=BROWN)
    t = pass_fail_table(-0.002, -1.0, 0.3, 0.5, 5, neg)
    assert t.at["(iv)", "result"] == "FAIL" and t.at["overall", "result"] == "FAIL"   # the reported bug
    t = pass_fail_table(0.004, 2.0, 0.05, 0.05, 8, pos)                                # boundaries pass
    assert (t["result"] == "PASS").all()
    t = pass_fail_table(0.004, 2.5, 0.01, 0.01, 9, pd.Series([0.002] * 4 + [-1e-5], index=BROWN))
    assert t.at["(iv)", "result"] == "FAIL" and t.at["overall", "value"] == 3
    t = pass_fail_table(0.004, 2.5, 0.01, np.nan, 9, pos.iloc[:4])
    assert t.at["(ii)", "result"] == "FAIL" and t.at["(iv)", "result"] == "FAIL"
    t = pass_fail_table(0.004, 1.99, 0.05, 0.01, 7, pos)
    assert t.at["(i)", "result"] == "FAIL" and t.at["(iii)", "result"] == "FAIL"


# =============================================================================
# 5. check_no_lookahead
# =============================================================================
def check_no_lookahead(raw, run_fn, start, end, cut_to=None, delta=0.05, n_months=8, min_obs=48,
                       seed=0, tol=1e-10, state_cols=STATE_COLS, pos_of=POS_OF, brown=BROWN):
    """
    Look-ahead audit on both the state side and the return side.

    Inputs
      raw      : raw dict (read_raw / make_synthetic_inputs).
      run_fn   : raw -> panel by month-end t with w, w_ao (positions set at end of t), RT, RAO (net returns
                 earned in t), m, hold, b_<k> (end-of-t states).
      start,end: window audited.
      cut_to   : if given, every input is cut at this month BEFORE anything runs, so later data are never
                 loaded (seen mode: 2022-07-31).
    Tests (one report row per test and month T)
      post_window    inputs cut at `end`: all states and returns in the window identical.
      truncate       inputs cut at T: states through T-1 identical (the old check; T >= start + min_obs).
      rb_perturb     every Brown industry +delta in T: states and returns before T identical, and
                     RT_T moves by exactly w_{T-1}*delta (RAO_T by w_ao_{T-1}*delta).
      f_perturb      hedge-file Mkt-RF, SMB, HML +delta in T: same, with move -w_{T-1}*delta*sum_k b_{T-1,k}.
      signal_perturb EMVENRGYENVREG for T spiked: states through T and returns through T+1 identical,
                     because the value for T is usable only at the end of T+1.
      A month-t return hedged with b_t, or a position earning its own month, leaves every end-of-month
      state unchanged. Both do change how RT_T and RAO_T respond to month-T data, and the two perturb
      tests measure that response. D and alpha are not compared: pi is an ex post window constant by design.
    Months T : fixed-seed draw of up to n_months each from hold entries/exits, held months (w_{T-1} != 0)
               and flat months (w_{T-1} = w_T = 0) inside the window.
    Output   : (report DataFrame, ok). Stop if ok is False.
    """
    start, end = pd.Timestamp(start), pd.Timestamp(end)
    if cut_to is not None:
        raw = truncate_raw(raw, cut_to)
        end = min(end, pd.Timestamp(cut_to))
    rows = []

    def add(test, T, pre, err):
        rows.append(dict(test=test, month=T, pre_diff=pre, ret_err=err, passed=bool(pre <= tol and err <= tol)))

    ret_cols = list(pos_of)
    all_cols = state_cols + ret_cols
    base = run_fn(raw)
    if _last_date(raw) > end:
        add("post_window", end, _maxdiff(base.loc[start:end], run_fn(truncate_raw(raw, end)), all_cols), 0.0)

    w, wprev = base["w"], base["w"].shift(1)
    inwin = base.index.to_series().between(start, end)
    groups = {
        "change": inwin & (w != wprev) & ((w == 0) | (wprev == 0)) & wprev.notna(),
        "held": inwin & (wprev != 0) & wprev.notna(),
        "flat": inwin & (wprev == 0) & (w == 0),
    }
    rng = np.random.default_rng(seed)
    months = pd.DatetimeIndex([])
    for mask in groups.values():
        ix = base.index[mask.to_numpy()]
        if len(ix) > n_months:
            ix = ix[np.sort(rng.choice(len(ix), n_months, replace=False))]
        months = months.union(ix)
    if len(months) == 0:
        add("coverage", pd.NaT, np.inf, 0.0)

    for T in months:
        pos = base.index.get_loc(T)
        Tm1, Tp1 = base.index[pos - 1], base.index[min(pos + 1, len(base) - 1)]
        before = base.loc[start:Tm1]
        if T >= start + pd.DateOffset(months=min_obs):
            add("truncate", T, _maxdiff(before, run_fn(truncate_raw(raw, T)), state_cols), 0.0)

        p = run_fn(_bump(raw, "ff49", T, list(brown), lambda v: v + delta))
        exp = {r: base.at[Tm1, c] * delta for r, c in pos_of.items()}
        add("rb_perturb", T, _maxdiff(before, p, all_cols), _ret_err(base, p, T, exp))

        p = run_fn(_bump(raw, "hedge_fac", T, HEDGE_COLS, lambda v: v + delta))
        bsum = sum(base.at[Tm1, f"b_{k}"] for k in HEDGE_COLS)
        exp = {r: -base.at[Tm1, c] * delta * bsum for r, c in pos_of.items()}
        add("f_perturb", T, _maxdiff(before, p, all_cols), _ret_err(base, p, T, exp))

        p = run_fn(_bump(raw, "EMVENRGYENVREG", T, ["EMVENRGYENVREG"], lambda v: 5.0 * v + 0.5))
        pre = max(_maxdiff(base.loc[start:T], p, state_cols), _maxdiff(base.loc[start:Tp1], p, ret_cols))
        add("signal_perturb", T, pre, 0.0)

    rep = pd.DataFrame(rows)
    return rep, bool(len(rep)) and bool(rep["passed"].all())


def _fred_month_end(df, sid):
    s = df.set_index("observation_date")[sid].astype(float)
    s.index = pd.DatetimeIndex(s.index) + pd.offsets.MonthEnd(0)
    return s


def _toy_run(raw, bug=None):
    """Minimal pipeline used only to test the auditor. bug: None, 'hedge_bt' (month-t return hedged with b_t),
    'own_month' (position set at end of t earns month t), 'signal_lag' (extreme flag of t used at end of t)."""
    leg = leg_pipeline(raw["ff49"], raw["hedge_fac"])
    sig = attention_signal(_fred_month_end(raw["EMVENRGYENVREG"], "EMVENRGYENVREG"),
                           _fred_month_end(raw["EMVOVERALLEMV"], "EMVOVERALLEMV"))
    ext = sig["extreme"].reindex(leg.index, fill_value=False).astype(bool)
    cross = ext & ~ext.shift(1, fill_value=False)
    usable = cross if bug == "signal_lag" else cross.shift(1, fill_value=False)
    hold = usable.astype(float).rolling(6, min_periods=1).max().astype(bool)
    m = leg["m"].fillna(0.0)
    B = leg[[f"b_{k}" for k in HEDGE_COLS]].fillna(0.0)
    F = leg[[f"f_{k}" for k in HEDGE_COLS]].to_numpy()
    panel = pd.DataFrame({"m": leg["m"], "hold": hold}, index=leg.index)
    for k in HEDGE_COLS:
        panel[f"b_{k}"] = leg[f"b_{k}"]
    for wcol, rcol, wpos in (("w", "RT", -m.where(hold, 0.0)), ("w_ao", "RAO", -m)):
        panel[wcol] = wpos
        if bug == "own_month":
            p, bb = wpos, B.shift(1)
        elif bug == "hedge_bt":
            p, bb = wpos.shift(1), B
        else:
            p, bb = wpos.shift(1), B.shift(1)
        gross = p * (leg["RB"] - (bb.to_numpy() * F).sum(axis=1))
        ov = B.mul(wpos, axis=0)
        trade = (0.0010 * wpos.diff().abs() + 0.0005 * ov["b_Mkt-RF"].diff().abs()
                 + 0.0025 * (ov["b_SMB"].diff().abs() + ov["b_HML"].diff().abs()))
        panel[rcol] = gross - trade.shift(1)                  # trade at end of t pays in t+1
    return panel


def test_check_no_lookahead():
    raw = make_synthetic_inputs(seed=0)
    kw = dict(start="1995-01-31", end="2009-12-31", n_months=4)
    rep, ok = check_no_lookahead(raw, _toy_run, **kw)
    assert ok, rep.loc[~rep["passed"]]
    for bug, must_fail in (("hedge_bt", {"rb_perturb", "f_perturb"}),
                           ("own_month", {"rb_perturb", "f_perturb"}),
                           ("signal_lag", {"signal_perturb"})):
        rep, ok = check_no_lookahead(raw, lambda r, b=bug: _toy_run(r, bug=b), **kw)
        failed = set(rep.loc[~rep["passed"], "test"])
        assert not ok and must_fail <= failed, (bug, failed)
        assert rep.loc[rep["test"] == "truncate", "passed"].all(), bug   # the old state-only check misses all three


def test_seen_cut():
    raw = make_synthetic_inputs(seed=0)
    assert _last_date(raw) > pd.Timestamp("2022-12-31")
    assert _last_date(truncate_raw(raw, SEEN_WINDOW[1])) <= pd.Timestamp("2022-07-31")


def run_unit_tests():
    for f in (test_attention_signal, test_pass_fail_table, test_seen_cut, test_leg_pipeline,
              test_check_no_lookahead, test_make_synthetic_inputs):
        f()
        print(f"{f.__name__}: ok")


# =============================================================================
# 6. __main__
# =============================================================================
if __name__ == "__main__":
    MODE = "synthetic"                 # "synthetic" -> "seen" -> "test"
    Z_BURN_IN = "nonzero"              # your reading; "calendar" reproduces my previous version
    WINDOWS = {"synthetic": TEST_WINDOW, "seen": SEEN_WINDOW, "test": TEST_WINDOW}
    PATHS = {"ff49": "ff49_industry_monthly.csv", **{sid: f"{sid}.csv" for sid in FRED_IDS}}
    fac = None                         # <- your FF5+UMD frame, as before
    hedge_fac = None                   # <- NEW: the team's FF3 frame (Mkt-RF, SMB, HML, RF), decimals

    start, end = WINDOWS[MODE]
    if MODE == "synthetic":
        run_unit_tests()
        raw = make_synthetic_inputs(seed=0)
    else:
        if fac is None or hedge_fac is None:
            sys.exit("Set `fac` (FF5+UMD) and `hedge_fac` (team FF3) first.")
        raw = read_raw(PATHS, fac, hedge_fac)

    cut_to = end if MODE == "seen" else None
    if cut_to is not None:
        raw = truncate_raw(raw, cut_to)            # nothing after 2022-07 reaches the audit or the run

    rep, ok = check_no_lookahead(raw, lambda r: run_panel(r, start, end, z_burn_in=Z_BURN_IN)[0],
                                 start, end, cut_to=cut_to)
    print(rep.groupby("test")["passed"].agg(["size", "sum"]).to_string())
    if not ok:
        print(rep.loc[~rep["passed"]].to_string())
        sys.exit("STOP: look-ahead audit failed.")

    run_panel(raw, start, end, z_burn_in=Z_BURN_IN, n_shuffle=5000, loo=True, verbose=True)  # prints the five tables
```

**Checks to run first.** Each of these should match exactly:
- `attention_signal(env, overall, z_burn_in="calendar")` should give the same `z`, `thr` and `extreme` as the old function on the real EMV files. If it doesn't, the difference is one of the conventions in the docstring, and you should keep the old behaviour (your implementation matches it). Tell me which convention it is.
- `leg_pipeline(ff49, fac[["Mkt-RF","SMB","HML","RF"]])` should reproduce the old leg exactly.
- Run with `("nonzero", FF3)`, the full backtest should again match your implementation to 1e-13. (iv) used to pass when it shouldn't have, so compare each pass/fail row, not just the count of passes.

**Choices I made inside these functions**
1. **Zero rule:** I count zero months in the trailing 60 calendar months. Months before 1985-01 are not counted as zero. The month is off when there are more than 6 zeros. Both burn-in settings use this rule.
2. **Off months:** An off month gets no z and is left out of the percentile history, the same as a zero month.
3. **Risk-free rate:** RF for R^B comes from `hedge_fac`. In the French files it is the same series as `fac`'s RF, and the driver can assert that.
4. **Criterion (iv):** I read "keeps its sign" as "stays positive": the full alpha and all five drop-one alphas must be above zero.
5. **Positive control:** The "planted shift" is mean(D) with the plant minus mean(D) without it, using the same seed. The plant lowers the Brown leg by 3 percentage points in every held month. That is deliberately large, because the check tests the pipeline, not statistical power.
6. **Look-ahead audit:** It compares positions and returns, not D, because pi is a constant computed over the whole window after the fact, by design.

The verdict is FAIL in your run and in mine, under either burn-in reading, so the default (your reading) doesn't change any conclusion. It is set to reproduce your pre-registered implementation, not picked on results. The seen-window numbers should be reported with the 2022-07 end: alpha of -0.16%/yr, t -0.24.
\end{lstlisting}

\paragraph{Fact-check (verbatim excerpt)} Sections ``Summary'', ``6. Bug and issue list'', ``R2 Summary'', ``R2.9'' of \texttt{factcheck.md} (10,321 words in full).

\begin{lstlisting}[style=transcriptsmall]
## Summary

**Verdict: the code is correct and reaches the same verdict as the reference.** Its numbers differ from the reference only through two readings that the prompt left open. "Do not implement" stands under either reading.

1. **Same verdict as the reference.** I ran the code on the real data with glue only; no line of Claude's file was edited.
   - It returns FAIL, with 1 of 4 components met (episodes).
   - Timing alpha -0.19%/yr, NW(6) t -0.25, one-sided p 0.60. Shuffle p 0.536. 10 episodes. Drop-one: 4 of 5 alphas keep the sign.
   - The reference gives -0.18%/yr, t -0.27, one-sided p 0.605, shuffle p 0.478, 10 episodes, and 3 of 5.
2. **Two readings explain every numeric gap. Neither is a coding error.**
   - **(A) z burn-in.** Claude requires 60 full calendar months of EMV history before the first z (lines 274-275). Its first z is therefore 1989-12 and its first threshold 1994-12. The reference (C5) lets the 48-nonzero minimum act as the burn-in, which gives a first z in 1989-01 and a first threshold in 1994-01. Under A:
     - the window starts at 1995-02 instead of 1994-03 (179 months, not 190);
     - the threshold moves by up to 0.076;
     - one extra crossing (data month 2002-12) opens a new hold;
     - pi rises from 0.729 to 0.809.
   - **(B) Hedge factor file.** Claude hedges with the SMB and HML of the `fac` it was given (the Ken French FF5 file) and subtracts that file's RF (lines 377-378). The team rule uses the team FF3 file. The prompt described only `fac`, so this is a gap in the prompt. SMB differs by up to 3.5 pp a month between the two files. As a result, b_t moves by up to 0.09 and m_t by up to 0.012.
   - **With both switched to the reference reading (variant AB), Claude's code reproduces every reference number to 1e-13.** That covers:
     - z, the threshold and every crossing over 1985-2026;
     - the holds, b_t, m_t, and net R^T and R^AO;
     - pi and D;
     - all 15 coefficients with their NW(6) and NW(12) t-statistics and p-values;
     - the decomposition, the 10 episodes and the five drop-one alphas;
     - all 5,000 shuffle draws, draw by draw (largest gap 8e-15; same seed and the same placement algorithm).
   - On the seen 2010-01 to 2022-07 window, AB reproduces the reference dry run exactly.
3. **One genuine bug, in the test harness, not the backtest (BUG-1).**
   - The synthetic GS10 generator (line 881) is a clipped random walk. It sits on its 0.5% floor in 76% of 1995-2009 months for seed 0, and in all 192 months of 1994-2009 for seed 16.
   - BOND is then constant, so I x BOND acts as an I main effect.
   - As a result, Claude's own positive control fails its own stop rule: only 4% of the planted effect reaches alpha (t 0.98), and I x BOND has t = -10.8.
   - Seed 16 of its null-seed check crashes on a rank-deficient design.
   - One edit fixes it (v1fix). The positive control then passes: 105% of the planted effect reaches alpha, t 4.15, shuffle p 0.005.
4. **Claude's own battery would have caught (B).** Its check (b)2 (reproduce the team trade to 1e-8) stops on the as-delivered code at a gap of 2.3e-3 a month. With the team FF3 file, the gap is 6e-17 and the check passes.
5. **Limits in the test design (not bugs):**
   - The look-ahead test compares states only. It misses 2 of 6 injected timing bugs, both on the return side.
   - Component (iv) does not require a positive alpha.
   - The burn-in reading is not a parameter, although the prompt asked for every silent reading to be one.
   - The "seen" mode includes 5 holdout months, and they flip the sign of the seen-window alpha.

## 6. Bug and issue list (most severe first)

1. **BUG-1, genuine (test harness), line 881.**
   - Cause: the synthetic GS10 random walk sticks to its 0.5% clip, so BOND is constant.
   - Consequences: (a) I x BOND stands in for an I main effect and absorbs 96% of the planted effect, so Claude's positive control, check (b)10, would STOP the protocol before the real run; (b) seed 16 gives a constant BOND over the whole window, so the rank check (line 521) raises and the (b)9 null-seed loop aborts.
   - The docstring expectation at lines 948-949 (alpha of about delta(1 - pi)mean(I m)) holds only if no regressor spans I. That is true on real data (99%) but false for this generator.
   - Fixed in v1fix by one edit. There is no effect on real-data numbers.
2. **D1, spec reading, not a parameter (discrepancy A), lines 274-275.**
   - Claude requires a full 60-month calendar window, so the "at least 48 nonzero months" clause can never bind. Its assumption 11 says so explicitly.
   - The prompt asked that each silent reading be a parameter. The three points the prompt named are parameters (lines 66-72), but this fourth one is hard-coded.
   - The reference reading (C5) gives the clause a role as the burn-in.
   - Impact: window 1995-02 instead of 1994-03, pi 0.81 instead of 0.73, shuffle p 0.54 instead of 0.48. The verdict does not change.
3. **D2, input gap in the prompt (discrepancy B), lines 377-378.**
   - The hedge and the leg's RF use `fac`, the FF5 file, whose SMB is built differently from the FF3 SMB. The Mkt-RF vintages also differ by up to 8 bp; HML and RF are identical.
   - The prompt never said that the team hedge uses the team FF3 file, so this is not Claude's error.
   - Impact on alpha: 0.004%/yr. Claude's own check (b)2 stops on it.
4. **D3, test coverage, lines 895-941.** `check_no_lookahead` compares only states set at the end of T. It misses the injected same-month hedge (M2) and same-month earning (M6) bugs, and it is conditional on `cfg.pub_lag` itself being right. Claude's separate lag probe (b)3 covers that last point.
5. **D4, pass-bar logic, lines 730, 744, 754.**
   - (iv) passes when every drop-one alpha has the sign of the full alpha, even if that sign is negative. So a negative alpha with five negative drop-one alphas would print "(iv) PASS".
   - This is a literal reading of "keeps its sign". The reference (C26) also requires alpha > 0; the one-line fix is `bool(a > 0 and same == nl)`.
   - No effect here: (i) fails, and as delivered (iv) is 4/5.
6. **D5, "seen" mode, lines 979-980.**
   - The seen window ends at 2022-12, so it includes holdout months 2022-08 to 2022-12. They flip the seen alpha from -0.16%/yr (t -0.24, through 2022-07) to +0.26%/yr (t 0.37).
   - The look-ahead call in that mode sets no `cut_to`, so its cuts reach 2024-08 despite the "seen data only" label.
   - This is cosmetic for the frozen test.

## R2 Summary

**Verdict: the corrected code is correct.** With its defaults (`z_burn_in="nonzero"`, hedge on the team FF3 file), it reproduces the reference to rounding error on every component and reaches the same verdict: FAIL, 1 of 4 (episodes only). "Do not implement" stands.

1. **All six round-1 items are fixed: BUG-1 and D1 to D5.** Claude's six unit tests pass the first time, with no edit to Claude's code. So do its three "checks to run first" and its checks (b)2, (b)9 and (b)10.
2. **The test-window run equals the reference.**
   - Result: FAIL 1/4. Alpha -0.18%/yr (NW(6) t -0.27, NW(12) t -0.29, p 0.790, one-sided 0.605). Shuffle p 0.478. 10 episodes. (iv) FAIL.
   - Largest gaps over all 15 terms: coefficients 3.8e-14, t-statistics 3.3e-13, p-values 1.9e-13.
   - The shuffle p is exact, and the 5,000 null draws agree draw by draw to 8.3e-15.
   - Every pass/fail row matches: (i) FAIL, (ii) FAIL, (iii) PASS, (iv) FAIL, overall FAIL.
3. **The two readings are now switches, and each does what it says.**
   - `z_burn_in="calendar"` reproduces Claude's old signal exactly. The full run equals round-1 v1 + B to 1e-14.
   - The new `hedge_fac` input puts the team FF3 file into the leg, the hedge and the overlay. The attribution still uses FF5 + UMD.
   - The verdict is FAIL 1/4 under either burn-in reading.
4. **The seen window equals the reference dry run:** alpha -0.12%/yr, t -0.19, shuffle p 0.510, 6 episodes, 0 of 4. Claude's instruction to report -0.16%/yr and t -0.24 for this window is wrong for its own new code. Those are round-1 numbers with the FF5-file hedge (N2).
5. **The new look-ahead audit works.**
   - It passes all 361 rows on four runs (synthetic, test under both readings, seen). Its largest return-side error is 5e-17.
   - It flags 7 of 8 injected timing bugs, including round 1's two misses (M2, M6). The eighth (M5) makes the audited pipeline raise an error, which also stops the protocol.
6. **One new issue matters: an integration gap (N1).**
   - The new `attention_signal` changes its output layout, and the reply's driver instructions do not mention it.
   - A literal drop-in crashes (`KeyError 'off'`). The obvious one-line patch then runs silently on a window 14 months too long (1993-01, n = 204): alpha -0.13%/yr, t -0.21. The verdict is still FAIL 1/4.
   - I bridged the gap with a value-preserving adapter.
7. **Other items are minor:** three inaccurate claims (N2 to N4) and some dead or incompatible code (N5).
8. **Extra observation (not about Claude's code).** On 100 further null seeds, the calendar shuffle is correctly sized, but the pre-registered NW(6) t over-rejects (13 of 100 beyond the 5% two-sided critical value). This affects the reference's component (i) as much as Claude's. It cannot turn a FAIL into a PASS (R2.9).

## R2.9 Extra check: size of the two tests on the synthetic null (not about Claude's code)

R12 ran 100 further null seeds (21-120, 200 draws each, no planted effect; `out/r2_null_seeds_d200_all.csv`).

| Statistic | Nominal | Observed (100 seeds) | Reading |
|---|---|---|---|
| Shuffle p <= 0.05 | 5 | 6 (P(>= 6) = 0.38) | Correctly sized. Kolmogorov-Smirnov against uniform: p 0.23 |
| Shuffle p <= 0.10 | 10 | 11 | Correctly sized |
| \|NW(6) t\| above the two-sided 5% t(175) critical value | 5 | 13 (P(>= 13) = 0.0015) | Over-rejects: the sd of t is 1.24 over 120 seeds |
| NW(6) t >= 2 (the (i) bar) | 2.4 | 5 (P(>= 5) = 0.09) | Probably oversized |

This concerns the pre-registered inference (C21), which v2 and the reference share. It is not a coding error in either. The likely cause is that D_t switches regime in blocks of 6 to 63 months, and 6 Newey-West lags do not span that. The calendar shuffle is built from the same blocks and stays correctly sized. The verdict is unaffected: the real alpha is negative, and a PASS needs all four components, including (ii).
\end{lstlisting}

\subsection{Exchange 03: robustness design (the multiple-testing framework behind M7)}\label{app:ex03}

Files: \texttt{exchange/\allowbreak{}03\_\allowbreak{}robustness\_\allowbreak{}design/\allowbreak{}}. Answered by: Claude (see the provenance note). Date: 2026-09-26. Purpose: robustness design.

\begin{center}\fbox{\begin{minipage}{0.95\linewidth}\small
\begin{tabularx}{\linewidth}{@{}lX@{}}
\textbf{Asked.} & A robustness and multiple-testing framework that decides the final verdict: which test families to correct and how, a deflated Sharpe specification with every input, at most two candidates with exact pass bars, the one check most likely to overturn ``Do not implement'', a committee's remaining objections, and a checklist of at most ten items. \\
\textbf{Came back.} & A framework of 1,716 words with its two tables (1,481 without them; cap 1,500), with a hand-computed deflated Sharpe table, two candidates (the momentum optimizer book and the EPA-ranked spread) and a frozen pre-1970 run of the book as the one decisive check. \\
\textbf{Verified.} & 93 checkable claims scored against the module outputs (67 correct, 15 partly correct, 11 wrong); all ten cells of its deflated Sharpe table reproduced to within 0.005; its power, holdout and multiple-testing arithmetic rerun (checks fc03 and fc03b). \\
\textbf{Kept.} & Families defined by the decision a test could change, the deflated appraisal ratio, Romano--Wolf and superior predictive ability (SPA) tests over the searched strategies, a one-sided holdout reading against a worthwhile alpha, and the frozen 1931--1969 run of the book, all implemented in M7; not its 155-row family, its post hoc composite, its $\pm$2\% margin, its MCCC pass test, exchange-traded fund (ETF) proxies or its crash budget. \\
\end{tabularx}\end{minipage}}\end{center}

\paragraph{My evaluation}

\begin{quote}\small
\textbf{The reply's statistics hold up; its reading of our books does not.} The fact-check scored 93 claims: 67 correct, 15 partly correct, 11 wrong and none unverifiable. Every formula and citation checks out, all ten cells of its deflated Sharpe table reproduce to within 0.005, and its bottom line (no candidate reaches ``implement'') matches mine.

\textbf{The framework is what I kept.} Families are defined by the decision a test could change, so loadings, placebos and exploratory rows sit outside every family. The ratio to deflate is the appraisal ratio: the optimizer book earns 4.53\% over FF5 (t = 5.21) but 2.17\% once UMD enters, so its raw Sharpe flatters it. A best-of-many claim needs a max-t or SPA test over every series I could have presented. And the book's one untouched window, 1931--1969, deserved a frozen, hashed, single run.

\textbf{Eleven claims were wrong, and three trace back to my prompt.} I printed the equal-weight book's UMD beta (0.99) beside the optimizer's numbers, and two of its claims rest on it; the optimizer loads 0.27 (t = 15.2). I never named the book's carbon source, so it guessed ``2022 intensities''. The other eight are its own. Romano--Wolf is not stricter than BH: a lone BH rejection on the 8,091-cell grid needs p $\leq$ 6.2e-6, while max-t would flag the grid's best cell (p = 0.00018) at an effective count below about 285. Loadings are 32\% of the nominal hits, not ``most''. M2 and M5 did test 25 bp. ``The BOND control helps'' contradicts M3. Its window cannot start before 1931-07 (the optimizer needs 60 months), and the 5\% tracking target holds prewar residual volatility down (5.14\% against 5.44\%). ``Insufficient evidence, not evidence of no alpha'' is false for the attention thesis: each of the six rules rejects a +2\% holdout alpha at one-sided p $\leq$ 0.0016. It also broke the word limit (1,716 against 1,500).

\textbf{Several of its pass bars would have misread our data.} The deflated Sharpe ratio (DSR) clause in its pre-1970 bar is decided by post-1970 data and fails on its own arithmetic beyond 7 trials (beyond 3 on M7's cluster-based inputs), so its ``one check'' could never overturn anything. Its $\pm$2\% equivalence test (two one-sided tests) calls a significantly negative holdout ``insufficient evidence''. Its MCCC control raises the EPA spread's alpha (the shock loading is negative), and the pass is knife-edge (t = 2.13 or 1.85, by shock construction). ``No significant carbon cost'' rewards low power.

\textbf{It missed the conflict between its own yardsticks.} At 42 effective trials a 5\% max-t bar is 3.03 and a 0.95 DSR needs about 3.85; the book's t of 3.07 sits between them, and the reply never says which governs. It ignored decay (2.99\% in validation, $-$0.34\% in the holdout) and gave no reading of a fail at roughly 0.6 power.

\textbf{What I adopted} (adopted\_checks.md, whose text was hashed at 08:49 PDT, before any pre-1970 return of the book existed): all 23,923 tests sorted into seven families; one-sided Holm and BH on the 73 distinct primary alphas; BH and BY on the grids, by window; Romano--Wolf and SPA over two search families; the deflated appraisal ratio, required in both windows; a holdout margin of a quarter appraisal ratio; a 25 bp cost stress; the frozen run with three pre-set readings; and a verdict map. The rule it lacked: max-t decides whether an alpha survives the search, and only the DSR can move a verdict. I dropped the 155-row family (M1's Fisher test, p = 1e-32, would ``survive'' it), the post hoc composite, the $\pm$2\% margin, the MCCC pass test, ETF proxies and a crash budget set after 2009.

\textbf{The frozen run passed and the verdict held.} The book earned 1.95\% a year over FF3+UMD before 1970 (t = 2.45; halves 1.71\% and 0.90\%), but its deflated appraisal ratio is 0.559 and its holdout alpha $-$0.34\%, so it moves to paper trading at pilot size. Do not implement stands.

\end{quote}

{\raggedright\paragraph{Prompt (verbatim)} \texttt{prompt.md}, 997 words (994 excluding Markdown symbols).\par}

\begin{lstlisting}[style=transcript]
You are the skeptical member of an investment committee. Your job is to design the robustness and multiple-testing framework that decides my project's final verdict. Everything you need is below.

SETTING. This is my solo extension of our team's MFE 230GA (Active Asset Management, UC Berkeley) final project. The brief grades thesis 40%, execution 40% and originality 20%, and requires "Do not implement" when there is no alpha. Our team's thesis: short a Brown industry leg (Util, Ships, Aero, Steel, BldMt, from the Fama-French 49), hedged with rolling 60-month FF3 betas, for 3 or 6 months after "climate-transition attention" crosses its expanding 80th percentile. The six variants use a raw or macro-purified signal with either hold, or a continuous weight. Validation is 2010-01 to 2022-07, the holdout 2022-08 to 2026-07. The verdict is Do not implement. I ran eight modules; an independent verifier reproduced each. Each module names its primary tests in its code; only the 1994-2009 rule (item 7) has a hash and timestamps proving they came first.

EVIDENCE (annualized; Newey-West(6) t, p from t(n-k)).
1. Team record. Validation: the 6-month rules earn FF3 alphas of 2.1% to 2.5% (t 2.2 to 2.6). Holdout: all six lose 0.5% to 3.8% a year net (FF3 alpha t down to -2.19), and the IC flips from -0.09 to +0.12. Audit: the team's "2 of 32 cells survive Holm" were two 24-month COVID regressions with normal p-values (0 of 32 with t(n-4)); a missing October 2025 CPI print created the whole Pure-versus-Original holdout gap; the emissions file has no source, units or date.
2. The signal is FRED EMVENRGYENVREG, a volatility-news tracker for energy and environmental regulation. Its z-score correlates -0.0004 with the Ardia et al. Media Climate Change Concerns index (MCCC) and 0.12 with climate policy uncertainty (CPU), through overall volatility news (partial r 0.05). It is exactly zero in 31 of 48 holdout months.
3. MCCC and CPU in the identical machinery: 0 of 12 validation alphas reach t 1.96 (EMV: 4 of 6), all six MCCC holdout alphas are negative, and the smallest Holm p over 40 primary tests is 1.00. A signal-free always-short Brown position earns 71% of the COVID alpha.
4. Cristhian's tests (8-and-8 legs, FF5+UMD+commodity controls, HML source, 5/10/25 bp costs): 8,091 alpha tests, 0 positive survivors under Holm or BH; 0 of 900 attention-rule holdout alphas are positive.
5. Alpha versus beta. A traded 10-year Treasury factor (BOND) plus UMD makes every holdout alpha slightly more negative; rates moved holdout returns by 0.5 pp a year at most. The Lewellen-Nagel beta-timing term is negative (-0.6% to -1.1%), and the always-short benchmark carries it too. The COVID gain is market and size exposure the lagged hedge left open.
6. Emissions. On EPA supply-chain intensities Aero and Ships rank 31st and 20th of 49 (the team file: 3rd and 2nd of 41). The EPA-ranked Green-minus-Brown spread earns 6.4% post-2010 (t 1.72) and 9.6% in the holdout (t 1.61), but its pre-specified FF5+UMD alpha is 4.7% (t 1.52). Its full-sample alpha, 4.5% (t 2.58), comes from short CMA and UMD loadings and a 2022 sort applied back to 1970.
7. Frozen rule. A pre-registered EMV-share rule, scored once on the unseen 1994-03 to 2009-12 window, fails 3 of 4 pass-bar components: alpha -0.18% (t -0.27), calendar-shuffle p 0.48.
8. Shrinkage timing (Lehnherr-Mehta-Nagel style ridge). Cross-validation picks the maximum penalty in 100% of monthly refits, so every forecast is the historical mean; attention lowers out-of-sample R2 at every fixed penalty.
9. Carbon-aware industry momentum (8 long, 8 short, t-12 to t-2, 10 bp). Net Sharpe 0.47, but the pre-registered FF5+UMD alphas are 1.74% (t 1.09) and 1.86% post-2010 (t 0.71), with UMD beta 0.99. A Grinold-Kahn optimizer book (IC fixed at 0.05, 5% ex-ante tracking; its alpha was exploratory) has full-sample alpha 2.17% (t 3.07, BH p 0.027) and IR 0.555, post-2010 alpha 3.12% (t 2.18), but holdout alpha -0.34% (t -0.15). Rank IC: 0.054 full sample, 0.030 holdout. Carbon frontier: at bound b = -1 long-leg carbon intensity falls 60% for an IR change of -0.005 (95% CI -0.110 to +0.096); only b = -2 costs significantly (-0.284, p 0.038).

TEST COUNT. The eight ledgers hold 23,923 tests: 155 primary, 8,113 robustness, 15,485 exploratory, 170 placebo or reference. About 4,200 have nominal p < 0.05, many of them factor loadings. They overlap heavily: shared holdout months, near-identical variants, one grid of 8,091 alphas.

UNSEEN DATA. FF49 returns start in 1926-07 and FF5 in 1963-07, but no module computed a strategy return before 1970.

WHAT I NEED, under these headings.
1. Families. How to control the family-wise error rate or the false discovery rate across a project this size: which families to define, which get Holm or Romano-Wolf and which get BH or BY, whether loadings and descriptive tests belong in any family, and how to treat dependence between near-duplicate variants.
2. Deflated Sharpe. How to deflate the best Sharpe ratios for the number of trials: which ratios, what counts as a trial, how to estimate the effective number of independent trials, and every input the formula needs. I have the optimizer book (679 months, annualized Sharpe 0.555, skew -0.28, raw kurtosis 5.77; post-2010: 199 months, 0.70, 0.29, 3.66) and the EPA spread (post-2010: 0.52, 0.35, 4.77).
3. Candidates. At most two that deserve the strictest scrutiny, and the exact tests that would move each to "implement": sample, benchmark, statistic, pass bar.
4. The one robustness check most likely to overturn "Do not implement", and the result that would do it.
5. What an investment committee would still object to after all of this.
6. A prioritized checklist, at most ten items: test, pass bar, what a fail means.

RULES. Work from these numbers. Cite a paper only with authors, year and the specific claim. Label any number you have not computed as a guess. Do not propose new signals. "No candidate can reach implement" is an acceptable answer if you say why. Under 1,500 words.
\end{lstlisting}

{\raggedright\paragraph{Reply (verbatim; Claude)} \texttt{claude\_\allowbreak{}response.md}, 1,907 words (1,716 excluding Markdown symbols).\par}

\begin{lstlisting}[style=transcriptsmall]
**Bottom line:** no candidate reaches "implement" on the current evidence. The attention thesis is finished (items 1-3, 7 and 8). Two strategies from outside the thesis deserve strict scrutiny. Both fail today, and only one still has an unseen sample that could change that.

---

## 1. Families

Group tests by the decision they could change, not by module.

- **Confirmatory family, controlled for family-wise error.** This holds only the tests that could flip the verdict: each candidate's pre-specified alpha test. Make them one-sided, because only a positive alpha leads to implementation.
  - Use Romano-Wolf stepdown. Romano and Wolf (2005, *Econometrica*) show that a bootstrap stepdown controls the family-wise error rate while using the dependence between tests, so it has more power than Holm.
  - Bootstrap with a stationary block bootstrap that resamples the same months across every series. Report Holm alongside, since Holm is valid under any dependence.
- **The 155 "primary" labels.** Only item 7 can be shown to have been pre-registered. Treat all 155 as one project-level family, not eight module families. A committee cannot check labels a module gave itself.
- **Robustness grids, controlled for false discovery rate.** Cristhian's 8,091 alphas are a sensitivity analysis.
  - BH (Benjamini and Hochberg 1995) holds under positive dependence. Benjamini and Yekutieli (2001) prove this for PRDS test statistics, and near-duplicate one-sided alphas are plausibly PRDS.
  - BY, from the same 2001 paper, holds under any dependence but costs a factor of about ln m + 0.58. That is about 9.6 for m = 8,091 (my arithmetic, a guess).
  - BH already finds zero positive survivors, and BY and Romano-Wolf are stricter, so they will not rescue anything.
  - Report the grid as a distribution: the share of positive alphas and the median t.
- **Exploratory tests (15,485): no family.** They generate hypotheses. Nothing here can be confirmed on the data that produced it, and the optimizer book's full-sample alpha sits in this group.
- **Excluded: loadings, descriptive statistics and placebos.**
  - Loadings (UMD beta 0.99, the beta-timing term) are diagnostics of the mechanism. Report them with confidence intervals.
  - Placebos calibrate the null; they are not hypotheses.
  - Most of the roughly 4,200 nominal hits are loadings, so that count tells you nothing.
- **Near-duplicate variants**, in order of preference:
  1. Pre-specify one composite statistic per strategy family, for example the equal-weight average of the six variants' alphas. That is one test, with more power.
  2. Use a max-t bootstrap such as Romano-Wolf, which handles the correlation automatically.
  3. To count tests, use Nyholt (2004): M_eff = 1 + (M-1)(1 - Var(lambda)/M), where lambda are the eigenvalues of the variants' return-correlation matrix.
- **"Best of many" claims.** Use White's Reality Check (2000, *Econometrica*) or Hansen's SPA test (2005, *JBES*). Both test whether the best strategy beats the benchmark after allowing for the search. SPA is less distorted by poor, irrelevant alternatives.

## 2. Deflated Sharpe

**Formula.** Bailey and Lopez de Prado (2014, *JPM*) define the deflated Sharpe ratio as DSR = PSR(SR0), where:

PSR(SR0) = Phi[(SR^ - SR0)*sqrt(T-1) / sqrt(1 - gamma3*SR^ + ((gamma4-1)/4)*SR^^2)]

SR0 = sqrtV[SRn] * [(1-gamma)*Phi^-1(1 - 1/N) + gamma*Phi^-1(1 - 1/(N*e))]

**Inputs:**
- SR^ as a monthly figure, not annualized
- T, the number of months
- skewness gamma3
- raw kurtosis gamma4
- N, the effective number of independent trials
- V[SRn], the variance of Sharpe estimates across trials, in monthly units
- gamma ~ 0.5772 (the Euler-Mascheroni constant)
- pass bar: DSR >= 0.95

**Which ratios to deflate:**
1. The book's Sharpe ratio.
2. The book's appraisal ratio (alpha divided by residual volatility). The claim is alpha, and with a UMD beta of 0.99 the raw Sharpe is mostly momentum.
3. The EPA spread's post-2010 Sharpe ratio.

**What counts as a trial:** every return series you looked at that could have been presented as "the strategy." That includes the team variants, the MCCC and CPU versions, Cristhian's grid, momentum lookbacks, carbon bounds, the optimizer's IC and tracking-error settings, and the choice of sample window (full versus post-2010). Placebos and loadings are not trials.

**Effective number of trials:** cluster the trials' return series. Lopez de Prado and Lewis (2019, *Quantitative Finance*) use clustering of strategy returns to estimate how many trials are effectively independent. Alternatively, use the eigenvalue M_eff above.

**V[SRn]:** use the variance of Sharpe ratios across clusters. Setting it to the null value 1/(T-1) gives a lower bound on SR0.

**Autocorrelation:** PSR assumes independent returns. Your Newey-West t-statistics are below their iid equivalents (EPA: 1.72 against about 2.1). Correct with Lo (2002, *FAJ*), which gives the serial-correlation adjustment for annualized Sharpe ratios and their standard errors.

**Hand arithmetic with V = 1/(T-1).** These are my calculations from your summary statistics; treat every figure as a guess until you code it.

| Series | PSR(0) | DSR, N=10 | DSR, N=100 | DSR, N=1,000 |
|---|---|---|---|---|
| Book, full sample (0.555, T=679) | ~1.00 | ~0.99 | ~0.94 | ~0.81 |
| Book, post-2010 (0.70, T=199) | ~0.998 | ~0.90 | ~0.62 | ~0.34 |
| EPA, post-2010 (0.52, T=199) | ~0.98 | ~0.71 | - | - |

- **EPA spread:** DSR is only about 0.94 to 0.95 at N = 2 (the team ranking versus the EPA ranking), before the autocorrelation correction. It fails.
- **Book, appraisal ratio:** roughly t/sqrtyears ~ 0.41. It clears DSR >= 0.95 only if the effective number of trials is at most about 7 (guess). Your project almost certainly has more, so the full-sample alpha does not survive deflation.

## 3. Candidates

**A. Carbon-aware momentum optimizer book.** It has the strongest statistics, but its alpha was exploratory.
- **Holdout:** -0.34% with a standard error of about 2.3% (backed out from the t). Too short either to confirm the alpha or to kill it.
- **Holdout rank IC:** 0.030, about one standard error below 0.054. I guess that standard error at about 0.02 for 48 months x 49 industries. The signal has not clearly died.
- **Sample:** 1927-07 to 1969-12. FF49 starts in 1926-07 and the momentum signal needs 12 months of history. Nothing has touched this period.
- **Freeze before running:** the code, IC 0.05, 5% tracking error, b = -1, 10 bp costs, and the rule for industries with missing early data. Hash and timestamp it as in item 7, then run it once.
- **Benchmark:** FF3 + UMD, because FF5 starts in 1963-07. As a secondary test, FF5 + UMD on 1963-07 to 1969-12. The carbon ranks use 2022 intensities. That is acceptable here because carbon acts as a constraint, not as the source of return.
- **Statistics:** the one-sided NW(6) t on alpha, and the alpha of the book minus the unconstrained book.
- **Pass bar:**
  - alpha t >= 2.0;
  - alpha positive both before and after 1948;
  - no significant carbon cost;
  - project-wide DSR on the appraisal ratio >= 0.95.
- **Power (guess):** your t implies residual volatility of about 5% a year, so the standard error is about 0.8% a year over 43 years. A true 2% alpha would clear t = 2.0 about two-thirds of the time. Residual volatility before the war is probably higher, which lowers power.

**B. EPA Green-minus-Brown spread.** This one is weaker.
- Its pre-specified alpha has t 1.52.
- Its full-sample alpha comes from CMA and UMD loadings and a 2022 sort applied back to 1970.
- Pastor, Stambaugh and Taylor (2022, *JFE*) attribute green stocks' outperformance in 2012-2020 to unexpected rises in climate concern, not to higher expected returns.
- **Test:** post-2010 FF5 + UMD alpha with MCCC innovations added as a control. MCCC is an existing series, not a new signal.
- **Pass bar:** one-sided alpha t >= 2.0 and DSR >= 0.95.
- It has no clean unseen sample. Emissions intensities as they stood at the time do not exist before 2010, and a pre-1970 test says nothing about climate. It cannot reach "implement" within this project; its only route is live paper-trading.

## 4. The one check

**The frozen pre-1970 run of the book.** The result that would overturn "Do not implement":
- FF3 + UMD alpha of roughly 1.5% a year or more (guess) with one-sided t >= 2.0;
- the same sign in both halves;
- DSR >= 0.95 at a documented effective number of trials.

Even that would justify only "implement at pilot size with paper-trading," because the post-1970 evidence is exploratory and the holdout is flat. It would also change the verdict for a different strategy, not the attention thesis. EMV data start in 1985, and the frozen 1994-2009 test has already failed.

## 5. What the committee will still object to

- **It is just momentum.** The book is industry momentum with a carbon screen (UMD beta 0.99). Why not buy a cheaper momentum factor? The incremental alpha is small and was found by searching.
- **Crash risk.** Daniel and Moskowitz (2016, *JFE*) show momentum suffers rare, persistent crashes in panic states after market declines, when the market rebounds. The full-sample skew is -0.28.
- **Low power.** Forty-eight holdout months cannot tell 0% from a 2-4% alpha. "Do not implement" rests on insufficient evidence, not on evidence of no alpha.
- **Unverified pre-registration.** Only item 7 is provable. With 23,923 tests, every other "primary" label has to be taken on trust.
- **Data provenance.** One missing CPI print created an entire result. The emissions file has no source. The 2022 carbon ranks are applied back to 1970.
- **Implementability.** FF49 portfolios cannot be traded directly. ETF proxies, borrow costs for shorting Aero and Ships, and costs above 10 bp are untested.
- **Regime.** The IC flip and the across-the-board holdout losses coincide with the 2022 rate and energy shock. The BOND control helps, but 48 months cover a single regime.

## 6. Checklist

| # | Test | Pass bar | A fail means |
|---|---|---|---|
| 1 | Frozen pre-1970 run of the book | alpha t >= 2.0; positive in both halves | The book is a post-1970 artifact; close it |
| 2 | DSR on the book's appraisal ratio, effective trials from clustering | >= 0.95 | The full-sample alpha is a product of selection |
| 3 | Project-wide one-sided Romano-Wolf on the 155 primary tests | any positive survivor at 5% | Confirms the verdict |
| 4 | Hansen SPA over all strategy return series against zero alpha | p < 0.05 | The best result from the search is noise |
| 5 | Holdout power and equivalence (TOST, +/-2% a year margin, a guess) | 90% interval inside the margin | The holdout cannot settle it; say "insufficient evidence" |
| 6 | EPA spread with MCCC-innovation control | alpha t >= 2.0 | Its returns came from a shock in climate concern |
| 7 | Data audit: ALFRED vintages for CPI and EMV, emissions source | headline figures move < 0.1 pp | Rerun everything downstream |
| 8 | Pre-registered crash stress on the book (1932, 2009) | drawdown within the stated budget | Needs smaller size or an overlay |
| 9 | ETF proxies with 25 bp costs | net alpha > 0 | The alpha exists only on paper |
| 10 | Carbon frontier re-run on pre-1970 data | IR change at b = -1 not significant | The carbon constraint is not free |
\end{lstlisting}

This exchange had one round: no follow-up prompt was sent.

\paragraph{Fact-check (verbatim excerpt)} Section ``Summary'' of \texttt{factcheck.md} (6,343 words in full).

\begin{lstlisting}[style=transcriptsmall]
## Summary

I extracted 72 checkable claims and recommendations from the reply: 47 correct, 17 partly correct, 8 wrong, 0 unverifiable. The statistics are sound and the arithmetic is right. The main errors are about our own strategy, and they matter because they sit in the reply's lead candidate.

1. **The formulas and hand arithmetic check out.**
   - The PSR and DSR formulas and their inputs match Bailey and Lopez de Prado (2014), including monthly SR, raw kurtosis and the expected-maximum term.
   - The BY constant is right (9.576 at m = 8,091), and so is the Nyholt formula.
   - All ten cells of its DSR table reproduce to within 0.005.
   - Its appraisal ratio (0.41) and its break-even count of about 7 trials both reproduce: 0.400 to 0.408, and N at most 6 to 7.
2. **Two facts about the lead candidate are wrong.**
   - The UMD beta of 0.99 belongs to the EW book (recomputed: 0.987). The optimizer book's UMD beta is 0.27 (t 15.2, R2 0.39).
   - The book's carbon constraint does not use "2022 intensities". It uses the team's undated `emissions_ff_industry.csv` (M5 section 3; `m5lib.load_inputs`), the file M4 found has no source and ranks Aero and Ships implausibly.
   - The 2022 EPA factors are used only by the EPA spread.
3. **One statistical claim is wrong.** Romano-Wolf is not stricter than BH.
   - BH, like Holm, needs p at most q/m (6e-6 here) before it rejects anything. A max-t stepdown uses the dependence between tests and can reject when BH rejects nothing.
   - The grid's best positive cell has one-sided p 0.00019. A max-t test would flag it if the grid's effective number of tests were below about 270.
   - That cell is a 24-month COVID alpha, the known small-sample artefact, so the conclusion "nothing implementable is rescued" probably survives. The stated reason does not.
4. **The frozen pre-1970 test is feasible, but it is mis-specified and low-powered.**
   - The optimizer needs 60 months of returns, so the window is 1931-07 to 1969-12 (462 months), not 1927-07.
   - Only 37 of its 41 covered industries exist before 1963.
   - With a true 2% alpha, the chance of reaching t 2.0 is 0.63 at post-1970 residual volatility, 0.33 at 1.5 times that level and 0.20 at twice it. So a fail is the likely outcome even if the alpha is real.
   - The pass bar contradicts itself. Its "project-wide DSR on the appraisal ratio" clause either cannot be met (Claude's own arithmetic fails it for more than 7 trials) or should use N = 1 for a single frozen trial.
5. **Two checklist items can give opposite answers.** DSR at least 0.95 requires t at least E[max_N] + 1.645 (3.10 at N = 8). A 5% max-t or SPA-type test requires only the 95th percentile of the maximum (2.49 at N = 8).
   - The book's alpha has t 3.07 and one-sided p 0.0011.
   - It passes a max-t test for up to about 46 effective series, but it fails DSR from N = 7.
   - The reply does not say which yardstick governs.
6. **The "low power" objection overreaches for the attention thesis.** Every one of the four hold variants rejects a +2% holdout alpha at p at most 0.001, and their 90% intervals lie wholly below zero. The low-power point is right only for the book.
7. **Its EPA test and its stated reading of a fail are backwards on our data.**
   - I ran it as an exploratory check. With a same-month real-time MCCC shock added, the post-2010 FF5+UMD alpha rises from 4.9% (t 1.52) to 6.7% (t 2.13) over the same 186 months (MCCC ends 2025-06).
   - The shock loading is negative: -0.52% a month per 1-sd shock, t -2.21. So the spread lost when concern rose, the opposite of the Pastor, Stambaugh and Taylor (2022) channel.
   - The pass is knife-edge. Lagged shock: t 1.06. MCCC-transition shock: t 1.95. CPU shock: t 1.97. FF3 base: t 1.76.
8. **Checklist status.**
   - Three items have outcomes that the existing numbers already settle:
     - 3, Romano-Wolf on the primaries, cannot pass: the largest positive primary alpha t is 1.89.
     - 5, the TOST, I computed here.
     - 7, the EMV vintages, pass on ALFRED.
   - Two items are partly done already (8 and 9).
   - One cannot be run on our data (ETF proxies).
   - The reply is also over its word limit: 1,716 words, or 1,469 without the two tables and formula lines, against "under 1,500".
\end{lstlisting}

\subsection{Exchange 04: red team of the draft report (two rounds)}\label{app:ex04}

Files: \texttt{exchange/\allowbreak{}04\_\allowbreak{}red\_\allowbreak{}team/\allowbreak{}}. Answered by: Claude (see the provenance note). Date: 2026-09-26. Purpose: red-team review of the conclusion.

\begin{center}\fbox{\begin{minipage}{0.95\linewidth}\small
\begin{tabularx}{\linewidth}{@{}lX@{}}
\textbf{Asked.} & Round 1: a red team of the draft summary and its claims (five weakest claims, at most four alternatives with separating predictions, the case against ``Do not implement'' and its reversing result, three edits, two least certain assumptions). Round 2: the same on the revised summary, now with the best evidence against the verdict, at most three alternatives, and edits no longer than the text they replace. \\
\textbf{Came back.} & Round 1: 1,340 words (cap 1,200) with a reversal test on the untimed EPA-ranked short. Round 2: 1,091 words (cap 1,000) arguing ``not proven'' rather than ``disproven'', with three alternatives (wrong legs, direction-free regulation news, a crisis hedge) and three replacement texts. \\
\textbf{Verified.} & Every criticism checked against the module outputs, and the missing tests run (checks fc04 and fc04b): holdout paths, intervals and timing, frozen-test power, a COVID bootstrap, the rules rebuilt without each Brown industry, MCCC and CPU on EPA-ranked legs, a direction split of the holdout entries, a crisis split of the timing gain, and our team's signal on both windows it was not designed on. \\
\textbf{Kept.} & Thirteen summary rewordings, the holdout interval and timing loss, and three alternatives tested and retired; not its edits as written, its cost comparison (the timing alpha is already net), its one-industry reading of the holdout, or its case for capital in the momentum book. \\
\end{tabularx}\end{minipage}}\end{center}

\paragraph{Round 1, on the first draft} Kept verbatim in \texttt{report/\allowbreak{}transcripts\_\allowbreak{}src/\allowbreak{}}.

\paragraph{My evaluation of round 1}

\begin{quote}\small
The reply refereed my wording well and my evidence badly. Every criticism aimed at a sentence landed, but almost every number it needed from our files was guessed, and four of the six guesses I could check were wrong. I checked each point against the modules and ran the missing tests (\texttt{checks/fc04\_checks.py}, exploratory).

\textbf{All five criticisms of my wording were right, and the executive summary now follows them.} ``900 of 900'' is one result: corrected Original and Pure coincide in all 48 holdout months, so the six rules have four return paths, and the first principal component carries 90\% of their variance. The summary and Section 2.3 now lead with Pure 6m (FF3 alpha $-$1.97\% a year, t = $-$1.30, six entries; 95\% interval $-$5.0\% to +1.1\%). Its power arithmetic for the frozen test is right (standard error 0.68\%, so only an edge of about 1.9\% a year passes with 80\% probability) and exposed an error of mine: ``+0.9\%'' was the shuffle's bar, while the t bar needed +1.4\%. The COVID edge's interval runs from $-$2.6\% to +6.2\%, and a block bootstrap puts the 71\% between $-$36\% and 233\% (90\%), so the ratio left the summary. The omitted CPU correlation and the ``real-time inputs'' overclaim are fixed in its words.

\textbf{Three corrections were wrong.} Its first fix, lead with the frozen rule's holdout alpha, is impossible: the frozen rule's zero clause switches it off from 2022-04 onward, so it never trades there. Its guess that the holdout |t| is below 1 fails for every rule (t = $-$1.30 to $-$1.90 corrected, $-$1.66 to $-$2.19 as our team ran them); only the residual term qualifies (t = $-$0.86). And its second edit would have put a false claim in the summary: ``2.4 standard errors below its validation alpha'' compares a timing alpha with a mean strategy alpha (1.45\%); the frozen rule's own seen-window timing alpha was $-$0.12\% (t = $-$0.19), so its proposed lesson, that the 1994--2009 test is informative, has no footing either. It flagged this as its least certain assumption. A true 1\% edge fails 70\% of the time, not half as it guessed.

\textbf{Three of its four alternatives were already answered by modules my prompt left out.} Same-month slopes on MCCC and CPU shocks have the wrong sign (t = $-$0.49 and $-$1.23). M2's commodity hedge left every holdout alpha negative; oil in the evaluation trims Pure 6m's to $-$1.50\% (t = $-$0.88), and breakevens add nothing once oil is in. The split is not what it guessed: Util, which it blamed, contributed a gain (+1.07 points), and Steel's $-$1.37 points match the $-$1.38-point gross loss only because that gain offsets Ships, Aero and BldMt ($-$1.08). On the EPA leg the frozen rule earns $-$0.01\% over 1994--2009.

\textbf{Its case against the verdict was aimed at the wrong result.} Its reversal test, a positive holdout alpha for the untimed EPA-ranked short, is met (+1.07\%, t = 0.79) by a position whose 1994--2009 alpha is $-$1.57\% (t = $-$0.74): a sign on 48 months, the counting its own first point warned against. The strongest case against sits in my own results: our team's signal, lagged, has a timing alpha of +0.69\% over 1994--2009 (t = 1.00) and +0.36\% (t = 0.70) pooled over both windows it was not designed on. My prompt left it out.

Claude is a good referee of prose and arithmetic and a weak one of data it cannot see. Its derived numbers were right to rounding and its citations accurate; it ran 12\% over its word cap (1,340 of 1,200). A red team is only as hard as its evidence: the next prompt sends the best result against my verdict, not only the results for it.

\end{quote}

{\raggedright\paragraph{Round 1 prompt (verbatim)} \texttt{04\_\allowbreak{}prompt.txt}, 795 words (782 excluding Markdown symbols).\par}

\begin{lstlisting}[style=transcript]
Read my draft twice: first as a skeptical investment-committee member who must sign off on "Do not implement", then as a finance-journal referee. Break the conclusion before a grader does. No summary, no praise.

SETTING. My UC Berkeley MFE 230GA final project (thesis 40%, execution 40%, originality 20%). I extended our team's halfway project alone. "What did you learn" is at its page limit, so an edit there must replace text.

THE EXECUTIVE SUMMARY, AS DRAFTED.
Our team's halfway project proposed a state-dependent climate trade: short the factor-hedged Brown industry leg (Util, Ships, Aero, Steel, BldMt) for three or six months after "climate-transition attention" crosses its past-only 80th percentile, sized to 5% residual volatility. Its best-specified form, which I froze before testing it, triggers on the energy-and-environmental-regulation share of volatility news, lagged one month. I asked three questions of the strategy: what the signal measures, what the legs are, and whether any version survives a window it was not designed on. The study uses Fama-French industry and factor returns, FRED news and macro series, two climate-concern indices (MCCC and CPU) and EPA supply-chain emission factors; every inference uses real-time inputs and small-sample p-values, each module fixes a primary test, and an independent verifier rebuilt each module's headline numbers. The recommendation is Do not implement, for four reasons. First, the "attention" series is FRED's equity-volatility news tracker for energy and environmental regulation (identical in 500 of 500 months), its z-score is uncorrelated with media climate concern (r = -0.0004), and the same machinery fed genuine climate measures earns 0 of 12 validation alphas with t >= 1.96. Second, the rule frozen before its test window earns -0.18% a year over 1994-2009 (t = -0.27, shuffle p = 0.48). Third, the COVID gain belongs to the position: an untimed short of the same leg earns 71% of it, and the rule's edge over that short is +1.84% a year (t = 0.87). Fourth, every timed rule has a negative 2022-2026 holdout alpha in all 900 net-of-cost configurations, all 28 zero-cost runs are negative too, and rates do not explain the loss. Carbon-aware industry momentum and shrinkage-based timing on the same data fare no better. In short: read the label, benchmark timing against the always-on position, and spend the last clean window on one frozen rule.

THE CLAIMS BEHIND IT (Newey-West t, 6 lags; p from t(n-k)).
1. Signal. On VIX and overall-EMV z-scores, the team z-score has R^2 = 0.169, all through overall EMV; its CPU correlation is 0.123 (Holm p = 0.074). Its zeros (39 of 59 months since 2021-10) changed no holdout entry.
2. Placebo. Overall EMV, a pure volatility measure, earns a 1.07% mean validation alpha, against the team signal's 1.45%.
3. Exposures. Green-minus-Brown is short HML (-0.233, t = -4.63), mostly through Brown; the untimed short shares the rules' beta-timing drag.
4. Holdout. Pure 6m nets -1.84% = hedge leakage -0.24% + residual -1.14% - cost 0.46%; Treasury and UMD factors make every holdout alpha more negative.
5. Frozen test. Its pre-registration was hashed before its code existed, but the rule was chosen after 2010-2026 was seen; it meets 1 of 4 pass-bar components, and passing needed about +0.9% a year.
6. Legs. The team emissions file is undocumented; Ships and Aero rank 2nd and 3rd of 41 there, 20th and 31st of 49 on EPA factors, and with Util carry 86% of the COVID gross P&L. An EPA-ranked spread earns a post-2010 FF5+UMD alpha of 4.7% (t = 1.52).
7. Testing. No positive alpha survives its pre-specified family among 23,923 logged tests; only the frozen test's pre-specification is timestamped, and NW(6) over-rejects on synthetic nulls (13 of 100 at 5%).
Already disclosed: static emissions snapshots, final vintages, an assumed EMV lag, a 48-month holdout.

WHAT I NEED, under these headings.
1. Weakest claims. Exactly five, ranked by effect on the verdict or the grade: the sentence, why the evidence is weaker than the wording, and the test or rewording that fixes it. Overreach toward "Do not implement" counts (t = 0.87 does not prove timing adds nothing).
2. Alternative explanations I have not ruled out. At most four, each with a prediction that separates it from mine on this data.
3. The case against "Do not implement": the strongest argument that the verdict is wrong or too broad, and the one result that would reverse it; if none survives these numbers, say so.
4. The three edits that would most improve the grade, each with its section and replacement text.
5. The two assumptions in your answer you are least sure of.

RULES. Work from these numbers; label any number you have not computed as a guess. Cite a paper only with authors, year and the specific claim. Stay under 1,200 words.
\end{lstlisting}

{\raggedright\paragraph{Round 1 reply (verbatim; Claude)} \texttt{04\_\allowbreak{}claude\_\allowbreak{}response.txt}, 1,385 words (1,340 excluding Markdown symbols).\par}

\begin{lstlisting}[style=transcriptsmall]
## 1. Weakest claims (ranked)

**1. "every timed rule has a negative 2022-2026 holdout alpha in all 900 net-of-cost configurations, all 28 zero-cost runs are negative too"**
*Why weak:* All 900 configurations run on the same 48-month path of the same Brown leg. That makes "900 of 900" one observation reported 900 times. You give no t-statistic, no interval and no count of entries. Since 39 of the last 59 months are zeros, the rule may have traded only a few episodes. Guess: at 5% residual volatility over 48 months, -1.84% a year has |t| < 1.
*Fix:* Lead with the frozen specification's holdout alpha, its NW t, a 95% interval and the number of entry months. Report how much of the variance in the 900 return series the first principal component explains. Call the rest "robustness of sign."

**2. "the rule frozen before its test window earns -0.18% a year over 1994-2009 (t = -0.27, shuffle p = 0.48)"**
*Why weak:* Failing to reject zero is not evidence of zero. The implied standard error is about 0.67% a year (0.18/0.27, derived), so the +0.9% pass bar needed t ~ 1.35. Guess: a true 1%-a-year edge would fail about half the time. The window also comes before transition risk was plausibly priced, and the legs use today's emissions snapshot.
*Fix:* Test against the effect you are rejecting. If the frozen rule's own validation alpha is near the 1.45% mean (guess), -0.18% is about 2.4 standard errors below it. Report that, plus the minimum detectable effect at 80% power: about 2.8 x 0.67 ~ 1.9% a year (derived).

**3. "the COVID gain belongs to the position ... +1.84% a year (t = 0.87)"**
*Why weak:* The implied standard error is about 2.1%, so the 95% interval runs from roughly -2.3% to +6.0% a year. The point estimate favours timing. "71%" is a ratio of two noisy P&Ls with no interval.
*Fix:* "Most of the COVID gain came from the position (71%). The timing edge of +1.84% a year cannot be told apart from zero or from 5% a year." Add a bootstrap interval on the 71%.

**4. "its z-score is uncorrelated with media climate concern (r = -0.0004)"**
*Why weak:* You quote MCCC and leave out CPU (r = 0.123, Holm p = 0.074). CPU is the index closest in concept to "environmental regulation." A Pearson r on a zero-inflated series is fragile. Also, R^2 = 0.169 means 83% of the signal's variance is not volatility. Counting newspaper articles that tie market volatility to environmental regulation is plausibly a transition-policy measure; the label alone does not rule that out.
*Fix:* Report both correlations, as Spearman and on innovations. Add the overlap between trigger months and top-quintile CPU and MCCC months. Reword to "weakly related to climate-policy uncertainty, unrelated to media climate concern."

**5. "every inference uses real-time inputs ... each module fixes a primary test"**
*Why weak:* Your own disclosures contradict both halves. Final vintages and static emissions snapshots are not real-time inputs, and only the frozen test's pre-specification is timestamped. A referee who catches this will discount everything after it.
*Fix:* "Every threshold, percentile and hedge uses only past data; macro series are final vintages and emissions a static snapshot. Each module names a primary test; only the frozen test's was timestamped before its data were seen."

## 2. Alternative explanations not ruled out

1. **Wrong horizon, not wrong signal.**
   - Pastor, Stambaugh and Taylor (2022) find green stocks outperformed in 2012-2020 mainly because climate concern rose more than expected. Ardia, Bluteau, Boudt and Inghelbrecht (2023) find brown stocks fall on days when climate concern rises unexpectedly.
   - If prices adjust in the same month, a short entered one month later and held 3-6 months arrives after the move.
   - *Prediction:* Green-minus-Brown loads positively on lag-0 CPU and MCCC innovations, and at zero or below at lags 1-6. Your explanation predicts no loading at any lag.

2. **Industry-specific holdout shocks.**
   - Guess: utilities repricing on data-centre power demand in 2024-25, and Ships and Steel in 2021-22, drive the -1.14% residual.
   - *Prediction:* splitting the residual by industry and year puts most of it in one or two legs and periods. Your explanation predicts losses spread across legs and years.

3. **Unhedged inflation or commodity exposure.**
   - You tested rates, not inflation. Steel, BldMt and Ships are commodity-sensitive.
   - *Prediction:* adding FRED WTI returns or 10-year breakeven changes to the hedge pulls the holdout residual toward zero. Your explanation predicts no change.

4. **Wrong leg.**
   - Ships and Aero rank 20th and 31st of 49 on EPA factors.
   - *Prediction:* the frozen rule applied to the EPA-ranked leg shows a validation timing edge over the untimed EPA short. Your explanation predicts about zero edge there as well.

## 3. The case against "Do not implement"

**The strongest argument:**
- The evidence rejects one timing rule on one poorly classified leg, but the verdict reads as a rejection of the climate trade.
- The untimed short captured 71% of the COVID gain, and the timing edge's point estimate is positive.
- The EPA-ranked spread earns 4.7% a year after 2010. Its implied standard error is about 3.1%, so the interval runs from roughly -1.4% to +10.8%.
- On that evidence the trade is unproven, not refuted.

**Why it mostly fails:**
- A t of 1.52 does not survive a family of 23,923 tests.
- The EPA spread was probably specified after the post-2010 results were known (my assumption).
- Bolton and Kacperczyk (2021) find that higher-emission stocks earn higher returns, so an always-on Brown short expects to pay a carbon premium.
- Pastor, Stambaugh and Taylor (2022) argue that green stocks' realized outperformance does not mean higher expected returns.

**Verdict:**
- "Do not implement the timed rule" survives. Nothing in the current data can reverse it, because both clean windows are spent.
- The broad version should be narrowed, not defended.
- The one result that would reverse the broad version is a positive 2022-2026 alpha for the untimed EPA-ranked short, which you do not report. If it is negative, say so and the broad verdict stands.

## 4. Three edits for the grade

**1. Executive summary, verdict sentence.** Replace "The recommendation is Do not implement, for four reasons." with:
> "The recommendation is Do not implement the timed rule. The untimed Brown short is unproven rather than refuted; an EPA-ranked Brown leg (post-2010 alpha 4.7% a year, t = 1.52) is the only open lead and does not survive the test family. Four reasons:"

**2. Executive summary, reasons two to four.** Replace with:
> "Second, the frozen rule earns -0.18% a year over 1994-2009 (t = -0.27), about [2.4] standard errors below its validation alpha, though the window cannot exclude an edge near 1.9% a year. Third, the untimed short earns 71% of the COVID gain; the timing edge, +1.84% a year (95% interval about -2.3% to +6.0%), is not distinguishable from zero. Fourth, the frozen specification's 2022-2026 alpha is [x]% (t = [y], [n] entries); all 900 configurations share that 48-month path and its sign, and rates do not explain it."

Drop "Carbon-aware industry momentum and shrinkage-based timing ... fare no better", or give each one's best alpha and adjusted p.

**3. "What did you learn".** Replace whichever lesson restates a count of tests or configurations:
> "Test the claim, not zero. Counting negative configurations felt decisive, but 900 configurations share one 48-month path. What separates results is how far the out-of-sample alpha sits from the in-sample one in standard errors, and the smallest effect the window could detect. By that standard the 1994-2009 test is informative and the holdout only weakly so."

## 5. The two assumptions I am least sure of

1. **The 2.4-standard-error figure.** It assumes the frozen rule's own validation alpha is near the 1.45% mean and that alpha divided by t gives a usable standard error. Your finding that NW(6) rejects 13 of 100 synthetic nulls at 5% suggests the standard errors are understated. If so, every interval above widens and the 2.4 shrinks.
2. **What drove the holdout.** I assumed there are few holdout entries, and that Util, Ships and Steel drove the loss. I inferred both from the zero count and from market history, not from your files.
\end{lstlisting}

\paragraph{Round 2, on the revised draft} The prompt sent the revised summary and the best evidence against it.

\paragraph{My evaluation of round 2}

\begin{quote}\small
This round did what round 1's evaluation asked for: the prompt sent the best evidence against my verdict as well as the evidence for it. The reply used that evidence as an advocate would, selectively, and that is where it went wrong. I checked every point against the modules and ran the tests it proposed (\texttt{checks/fc04b\_checks.py}, exploratory).

\textbf{All five weak claims found real gaps in the wording, and the summary now closes them.} ``Five high-emission industries'' hid that EPA factors rank Ships and Aero 20th and 31st of 49. ``Most of the COVID gain came from the position'' rested on a point estimate: the rule's edge is +1.84\% over the untimed short (t = 0.87) and 3.0\% over an exposure-matched one (t = 1.67), so the window cannot split the gain. ``p = 0.48'' needed its null, whose median is $-$0.22\%. ``Value-tilted'' described the raw leg rather than the hedged rule, and a wrong label explains why the thesis is untested, not why the returns fail, so the verdict sentence now says both. ``The last unused window'' is the live record.

\textbf{Four of its supporting claims were wrong.} The holdout does not indict ``the position, not the timing'': against an exposure-matched short the timing lost a further 1.0\% to 1.7\% a year (t = $-$0.99 to $-$1.72). The loss does not come from one industry either: rebuilt without any one industry, the rule still loses in 10 of 10 cases (FF3 alpha $-$1.26\% to $-$2.61\%). Both results now stand in the summary. The figures it called unreconciled add up in my own claims list. And the residual test it asked for is the frozen rule, which trades the signal with overall volatility news divided out, the same rule it had just dismissed as one I am not proposing.

\textbf{Each alternative failed its own prediction.} On EPA-ranked legs the same-month slopes on MCCC and CPU shocks became more wrong-signed (t = $-$2.95 and $-$2.17), and CPU timing added no detectable edge over the untimed EPA short (holdout +0.05\% to +1.22\%, largest t = 1.43; negative in 5 of 6 rules over 1994--2009); Section 2.1 and Appendix C now say so. With the administration as a crude proxy for policy direction, the short lost more after entries under Biden ($-$4.5 gross points) than under Trump ($-$1.0). And 80\% of the 6-month rule's 1994--2026 timing gain came in months with VIX at or below 30, 95\% outside NBER recessions.

\textbf{Its case against the verdict misread what I sent.} ``Positive in both unseen windows'' reads a pooled +0.36\% as two positives; the windows give +0.69\% (t = 1.00) and $-$1.11\% (t = $-$1.13). Its committee argument set that timing alpha against the 0.46\% holdout cost, but the alpha is computed on net returns, so the cost is already inside it. The deflated appraisal ratio is applied to the windows the momentum book was chosen on, and 2010--2026 is one of them, so it is no second out-of-sample pass. None of its three edits went in as written: two carry these errors, and the third runs 18 words against a 17-word original while claiming 17.

\textbf{Claude referees sentences well and reads numbers carelessly.} Sending the counter-evidence bought sharper, testable alternatives, but an advocate quotes the favorable half: 5.4\% without the 3.0\%, and the 1994--2009 drop-one result without the holdout. As in round 1, the error that reached its replacement text sat in the assumption it flagged as least certain, and it again ran over its cap (1,091 words of 1,000). Read its flags, test its predictions, and never paste its edits.

\end{quote}

{\raggedright\paragraph{Round 2 prompt (verbatim)} \texttt{round2/\allowbreak{}prompt.md}, 798 words.\par}

\begin{lstlisting}[style=transcript]
Read my draft conclusion as the skeptical investment-committee member who must sign "Do not implement", then as a finance-journal referee. No summary, no praise.

SETTING. My solo extension of our team's MFE 230GA final project (40% thesis, 40% execution, 20% originality).

EXECUTIVE SUMMARY (DRAFT).
Our team's halfway project proposed a state-dependent climate trade, and this report is my solo extension of it. The strategy shorts a factor-hedged Brown leg of five high-emission industries (utilities, shipbuilding and railroad equipment, aircraft, steel, construction materials) for three or six months after "climate-transition attention" crosses its past-only 80th percentile, sized to 5% residual volatility. I tested it on Fama-French industry and factor returns, FRED news and macro series, two climate-concern indices, media climate change concern (MCCC) and climate policy uncertainty (CPU), and EPA supply-chain emission factors. Every threshold, percentile and hedge uses only past data, p-values come from small-sample t distributions, and independent code rebuilt each module's headline numbers. The recommendation is Do not implement, for four reasons. First, the "attention" series is FRED's equity-volatility news tracker for energy and environmental regulation (identical in 500 of 500 months). Its z-score is uncorrelated with MCCC (r=-0.0004), its weak link to CPU (r=0.123, Holm p=0.074) runs through overall volatility news, and the same rules fed either climate index earn 0 of 12 validation alphas with t>=1.96. In plain terms, the rule shorts value-tilted utility and industrial stocks after stock-market volatility makes the news. Second, a cleaner version frozen before its test (the regulation share of volatility news, lagged one month) has a timing alpha of -0.18% a year over 1994-2009 (t=-0.27; p=0.48 against random re-datings of its holding periods); that window detects an edge of about 1.9% a year with 80% power, so the fail excludes a large edge, not a small one. Third, most of the COVID gain came from the position: the rule's alpha over an untimed short of the same leg is +1.84% a year (t=0.87; 95% interval -2.6% to +6.2%). Fourth, every one of our team's timed rules lost over the 2022-2026 holdout, before costs as well as after (6-month rule: FF3 alpha -1.97% a year, t=-1.30), the untimed short lost too (-1.72%, t=-1.09), and rates do not explain the loss; the 900 net-of-cost variants share that one 48-month window, so their common sign is one result. Of the alternatives, shrinkage timing collapses to the historical mean, and an optimized industry-momentum book that passed a pre-registered 1931-1969 test (alpha 1.95% a year, t=2.45) fails the deflated appraisal ratio and lost in the holdout (-0.34%), which earns it a paper-trading pilot, not capital. In short: read the label, benchmark timing against the always-short position, and spend the last unused window on one frozen rule.

CLAIMS BEHIND IT (annualized; Newey-West t, 6 lags).
1. Signal: R^2=0.169 on VIX and overall-EMV z-scores, all through overall EMV; its non-volatility topic share correlates 0.03 with CPU; same-month Green-minus-Brown slopes on MCCC and CPU shocks are wrong-signed (t=-0.49, -1.23).
2. Exposures: Green-minus-Brown loads -0.233 on HML (t=-4.63), mostly through Brown; the hedged rules keep -0.015 to -0.053 since 2010.
3. Holdout: the 6-month rule's -1.84% net is leakage -0.24%, residual -1.14% and cost 0.46%, from six entries; Steel carries -1.37 of -1.38 gross points. BOND, UMD or a commodity factor leave every alpha negative.
4. Frozen rule: chosen after 2010-2026 was seen; in position 142 of 190 months; off from 2022-04, so flat through the holdout.
5. Legs: Ships and Aero rank 2nd and 3rd of 41 in the undocumented team file, 20th and 31st of 49 on EPA factors. The EPA-ranked spread's post-2010 FF5+UMD alpha is 4.7% (t=1.52).
6. Testing: 73 primary alphas, best one-sided p=0.030, smallest Holm p=1.00; that t rejects 13 of 100 synthetic nulls at 5%.

BEST EVIDENCE AGAINST ME.
a. Our team's signal, lagged: timing alpha +0.69% over 1994-2009 (t=1.00; all five drop-one-industry alphas positive), +0.36% (t=0.70) pooled over both unseen windows.
b. COVID, against an exposure-matched benchmark: 5.4% (t=2.35), 3.0% (t=1.67) with in-window factors.
c. Momentum book, 2010-2026: alpha 3.12% (t=2.18).

WHAT I NEED, under these headings.
1. Weakest claims: five summary sentences, ranked by effect on the verdict (overreach toward "Do not implement" counts), each with the gap between evidence and wording, and the test or rewording that closes it.
2. Alternatives I have not ruled out: at most three, each with a separating prediction on this data.
3. The case against "Do not implement": the strongest argument that it is wrong or too broad, and the one result that would reverse it (or "none").
4. Three edits that would most improve the report, with replacement text no longer than the original.
5. Your two least certain assumptions.

RULES. Label any number you have not computed as a guess. Cite papers by authors, year and claim. Under 1,000 words.
\end{lstlisting}

{\raggedright\paragraph{Round 2 reply (verbatim; Claude)} \texttt{round2/\allowbreak{}claude\_\allowbreak{}response.md}, 1,108 words (1,091 excluding Markdown symbols).\par}

\begin{lstlisting}[style=transcriptsmall]
**1. Weakest claims (ranked by effect on the verdict)**

1. **Reason 2, the frozen rule.** This rejects a rule you are not proposing. It was chosen after 2010-2026 was seen and was in position 142 of 190 months, so it is mostly the untimed short. Your team's lagged rule earns +0.69% in the same window (t=1.00, all drop-one alphas positive), and the 1.9% power caveat applies to that result too. A p of 0.48 for a negative alpha means either the re-dating null is centred below zero or the test is two-sided. Say which. Fix: Edit 1, and name which window is "the last unused" one.

2. **Reason 4, the holdout.** "Every one of our team's timed rules lost" suggests breadth. The evidence is one industry (Steel, -1.37 of -1.38 gross points), six entries and t=-1.30. The timed and untimed shorts lost similar amounts, so the holdout indicts the position, not the timing the thesis sells. The -1.97% FF3, -1.84% net and -1.38% gross figures are never reconciled. Fix: report the holdout timing alpha (timed regressed on untimed) and the ex-Steel alpha, and name the model behind each number (Edit 2).

3. **Reason 1, "In plain terms..."**
   - R^2=0.169 leaves 83% of the signal unexplained by VIX and overall EMV.
   - The hedged rules keep at most -0.053 HML.
   - So both halves of the sentence overstate.
   - A mislabelled signal can still pay. This reason supports "don't call it climate", not "don't implement".
   - Fix: rerun the rule on the signal's residual after VIX and EMV. If alpha is unchanged, move the label point out of the verdict (Edit 3).

4. **Reason 3, COVID.** "Most of the gain came from the position" rests on a point estimate. The interval (-2.6% to +6.2%) allows a timing share of anything from none to all, and the exposure-matched benchmark gives 5.4% (t=2.35). Fix: write "the data cannot split the COVID gain between position and timing", report both benchmarks and say why you prefer one.

5. **Opening, "five high-emission industries."** Your EPA ranks put Ships and Aero 20th and 31st of 49. The summary also leaves out the EPA-ranked spread's 4.7% (t=1.52). The verdict therefore covers this construction, not the thesis. Fix: write "five industries our team labelled high-emission (two rank mid-table on EPA factors)", and add the EPA result to the evidence against you.

**2. Alternatives not ruled out**

1. **Wrong legs.** Your wrong-signed MCCC and CPU slopes contradict two papers:
   - Ardia, Bluteau, Boudt & Inghelbrecht (2023): green beats brown when climate concern rises unexpectedly.
   - Pastor, Stambaugh & Taylor (2022): green outperformance in 2012-2020 came from concern shocks.

   Prediction: on EPA-ranked legs the slopes turn positive and the MCCC- and CPU-timed validation alphas improve. If they do not, the thesis fails at industry level whatever the signal.

2. **Regulation news without a direction.** EMV (Baker, Bloom, Davis & Kost 2019) counts articles, so it carries no policy direction. Guess: in 2022-2025 the coverage was mostly IRA subsidies, tariffs and rollbacks that favoured Steel. Prediction: split entries by direction, using the administration in power as a crude proxy or hand-coding the articles. The short should earn after tightening news and lose after loosening news, and the Steel loss should cluster in tariff and IRA months.

3. **Crisis hedge.** Brown cyclicals fall in crises, and the signal loads on overall volatility news. The payoff may be crisis convexity (2008, 2020) rather than alpha. Prediction: timing alpha concentrates in months with VIX above 30 or in NBER recessions, and is near zero elsewhere. If so, price it as a hedge.

**3. The case against "Do not implement"**

The verdict is right for this rule, but it is written as if the climate thesis was tested. It was not: the signal is volatility news, two Brown legs are mid-emission, and the legs' climate slopes are wrong-signed.

For the rule itself, the evidence reads "not proven", not "disproven":
- positive in both unseen windows;
- robust to dropping any one industry;
- 5.4% (t=2.35) against an exposure-matched COVID benchmark;
- a holdout loss that comes from one industry.

On momentum: the deflated Sharpe ratio (Bailey & Lopez de Prado 2014) is built for in-sample selection. Applying it to a pre-registered pass double-counts trials. Suppose 2010-2026 was not used for optimisation. Then the book has two out-of-sample passes (t=2.45 and 2.18), and with the industry-momentum prior of Moskowitz & Grinblatt (1999) that argues for small capital, not only paper trading.

The committee's case for signing is stronger than the one you wrote. Your best unseen estimate (+0.36%) is below the holdout cost (0.46%), so the rule loses money even if that estimate is true. Bolton & Kacperczyk (2021) also find a carbon premium, so the always-short leg the timing sits on has a negative expected return.

Reversing result: none on existing data, because any rule built now has already seen both unseen windows. Only a frozen rule's forward record can reverse the verdict.

**4. Three edits**

*Edit 1, replacing Reason 2 (58 words; the original has 60):*
"Second, on data unseen when rules were set, our team's rule lagged one month earns a timing alpha of +0.36% a year (t=0.70), below its 0.46% annual holdout cost; a cleaner frozen variant earns -0.18% over 1994-2009 (t=-0.27). These tests detect only edges above about 1.4-1.9% a year, so they exclude a large edge, not a small one."

The 1.4% is my arithmetic from your t=0.70, not a computed result.

*Edit 2, replacing Reason 4 (52 words; the original has 57):*
"Fourth, over the 2022-2026 holdout the timed rules lost (6-month rule: FF3 alpha -1.97% a year, t=-1.30) about as much as the untimed short (-1.72%, t=-1.09), so timing added little; Steel accounts for nearly all of the gross loss, from six entries, and the 900 variants' common sign is one 48-month result."

*Edit 3, replacing "In plain terms..." (17 words; the original has 17):*
"In plain terms, the signal counts regulation mentions in volatility news; no test links it to climate concern."

**5. My two least certain assumptions**

1. I read "13 of 100 synthetic nulls" as meaning your nominal 5% test rejects about 13% of the time. If so, every t in the report is overstated, including 2.35, 2.45 and 2.18, and the summary's sentence on p-values needs a caveat. If it means a family-wise max-t p of about 0.13, ignore this point.
2. I assumed "timing alpha" is measured the same way for the team rule and the frozen rule. I also assumed the 0.46% holdout cost reflects the lagged rule's turnover in other windows. The cost comparison in Edit 1 depends on both.
\end{lstlisting}

Each round was one prompt and one reply; no follow-up was sent.

{\raggedright\paragraph{Fact-check} The check scripts are in \texttt{exchange/\allowbreak{}04\_\allowbreak{}red\_\allowbreak{}team/\allowbreak{}checks/\allowbreak{}} (\texttt{fc04\_\allowbreak{}apply\_\allowbreak{}revisions.py}, \texttt{fc04\_\allowbreak{}checks.py}, \texttt{fc04b\_\allowbreak{}apply\_\allowbreak{}revisions.py}, \texttt{fc04b\_\allowbreak{}checks.py}, \texttt{verify\_\allowbreak{}fc04b\_\allowbreak{}independent.py}), with their outputs in \texttt{checks/out/}.\par}
